% Les lentilles minces — Physique 1ère E — Cours signé OnBuch+
% Programme officiel MINESEC. Compiler avec : tectonic N07-lentilles-minces.tex
\newcommand{\DOCMATIERE}{Physique}
\newcommand{\DOCNIVEAU}{1ère E}
\newcommand{\DOCMODULE}{Module 3 : Optique}
\newcommand{\DOCLECON}{N07}
\newcommand{\DOCTITRE}{Les lentilles minces}
% OnBuch+ — préambule « cours signés » ENRICHI (compilé avec tectonic / XeLaTeX)
% Système visuel « Pop » d'OnBuch+ : fond crème (jamais blanc), bordures encre
% épaisses, ombres « dures » décalées, titres Archivo Black en capitales.
% Version MATHÉMATIQUES : graphes (pgfplots/tikz), boîtes pédagogiques
% (définition, propriété, méthode, attention, le savais-tu, exercice résolu,
% exercice, TP, bilan), page de garde et signature OnBuch+ ancrée en bas.
%
% Macros attendues dans main.tex AVANT \input{preamble} :
%   \DOCMATIERE  \DOCNIVEAU  \DOCMODULE  \DOCLECON (numéro)  \DOCTITRE
\documentclass[11pt]{article}

\usepackage{fontspec}
\usepackage[french]{babel}
\usepackage[a4paper,margin=2cm,top=3.4cm,bottom=2.8cm,headheight=1.4cm,footskip=1.3cm]{geometry}
\usepackage{amsmath,amssymb,amsfonts}
\usepackage{mathtools} % \xmapsto, \coloneqq... souvent utilisés par les modèles
\usepackage{cancel} % simplifications barrées (bilans de forces)
% \abs{z} (module, valeur absolue) et \norm{u} (norme), utilisés par les modèles
\providecommand{\abs}{}\let\abs\relax\DeclarePairedDelimiter{\abs}{\lvert}{\rvert}
\providecommand{\norm}{}\let\norm\relax\DeclarePairedDelimiter{\norm}{\lVert}{\rVert}
\usepackage[table]{xcolor} % [table] : fournit aussi \rowcolors (alternance de lignes)
\usepackage{pagecolor}
\usepackage{tikz}
\usetikzlibrary{arrows.meta,positioning,calc,shapes.geometric,shapes.misc,shapes.arrows,decorations.pathmorphing,decorations.pathreplacing,decorations.markings,patterns,shadows,mindmap,trees,fit,backgrounds,matrix,intersections,angles,quotes}
\usepackage{pgfplots}
\usepackage[version=4]{mhchem} % équations nucléaires \ce{^{235}_{92}U}
\usepackage[european,siunitx]{circuitikz} % schémas électriques
\pgfplotsset{compat=1.18}
\usepgfplotslibrary{fillbetween,groupplots} % aires entre courbes (calcul intégral)
\usepackage{enumitem}
\usepackage{fancyhdr}
% [most] charge toutes les bibliothèques tcolorbox (listings, minted, raster,
% documentation...) et gonfle inutilement la mémoire TeX sur un document long
% (~300 boîtes) : on ne charge que ce qui est réellement utilisé.
\usepackage[skins,breakable]{tcolorbox}
\usepackage{graphicx}
\usepackage{colortbl}
\usepackage{booktabs}
\usepackage{tabularx}
\usepackage{diagbox} % case d'en-tête diagonale des tableaux à double entrée
% Colonnes extensibles centrée / gauche / droite (C, L, R), souvent utilisées par les modèles
\newcolumntype{C}{>{\centering\arraybackslash}X}
\newcolumntype{L}{>{\raggedright\arraybackslash}X}
\newcolumntype{R}{>{\raggedleft\arraybackslash}X}
\usepackage{array}
\usepackage{multicol}
\usepackage{siunitx}
% parse-numbers=false : accepte aussi \SI{2,5 \times 10^{-3}}{...}
\sisetup{locale=FR,output-decimal-marker={,},group-minimum-digits=4,parse-numbers=false}
\DeclareSIUnit\ohm{\ensuremath{\symup{\Omega}}} % U+2126 absent de la police
\DeclareSIUnit\division{div} % réglages d'oscilloscope (ms/div, V/div)
\DeclareSIUnit\tr{tr} % tours (vitesse de rotation, ex. \SI{1500}{\tr\per\minute})
\DeclareSIUnit\ch{ch} % cheval-vapeur (puissance moteur, unité usuelle hors SI)
\DeclareSIUnit\franccfa{FCFA} % franc CFA (coûts, contexte camerounais)
\DeclareSIUnit\dioptre{\ensuremath{\delta}} % vergence d'une lentille (dioptrie)
\usepackage{qrcode}
% \degree : commande fréquemment utilisée par les modèles pour un angle ou une
% température en dehors de \SI{}{} (non fournie par les paquets chargés).
\providecommand{\degree}{^{\circ}}
% \qed (fin de démonstration), utilisé par les modèles sans amsthm
\providecommand{\qed}{\hfill$\square$}
% \barre : double barre des valeurs interdites dans un tableau de variations (tkz-tab)
\providecommand{\barre}{\ensuremath{\|}}
% environnement proof (amsthm non chargé) utilisé par les modèles
\ifdefined\proof\else\newenvironment{proof}[1][Démonstration]{\par\noindent\textit{#1.}\ }{\qed\par}\fi
\usepackage{lastpage}
\usepackage{hyperref}
\hypersetup{hidelinks}

% ============================================================
% Palette « Pop » OnBuch+ — mêmes hex que la classe P (plus_ui.dart)
% ============================================================
\definecolor{poporange}{HTML}{F59321}
\definecolor{popdark}{HTML}{E07A0C}
\definecolor{popcream}{HTML}{FFF6E8}
\definecolor{popink}{HTML}{1C1714}
\definecolor{poppurple}{HTML}{7A5AE0}
\definecolor{popgreen}{HTML}{2E9E6A}
\definecolor{poppink}{HTML}{E0457B}
\definecolor{popblue}{HTML}{2F7DE1}
\definecolor{popgold}{HTML}{C98A12}
\definecolor{popmuted}{HTML}{8A7F76}
\definecolor{popline}{HTML}{EADFD2}
\definecolor{popgray}{HTML}{8A7F76}
\colorlet{popgrey}{popgray}
% Alias tolérants (le modèle nomme parfois les couleurs différemment)
\colorlet{popwhite}{white}
\colorlet{popblack}{popink}
\colorlet{popred}{poppink}
\colorlet{popyellow}{popgold}
% Teintes claires pour fonds de tableaux / zones
\colorlet{popblueL}{popblue!10!white}
\colorlet{poporangeL}{poporange!14!white}
\colorlet{popgreenL}{popgreen!12!white}
\colorlet{poppurpleL}{poppurple!12!white}
\colorlet{poppinkL}{poppink!12!white}
\colorlet{popgoldL}{popgold!14!white}

\pagecolor{popcream}

% ============================================================
% Typographie
% ============================================================
\defaultfontfeatures{Ligatures=TeX}
\setmainfont{PlusJakartaSans-Medium.ttf}[Path=../../../fonts/,BoldFont=PlusJakartaSans-Bold.ttf]
% Maths en sans-serif (Fira Math) avec chiffres et lettres droites de Plus
% Jakarta, pour que les formules se fondent dans le texte.
\usepackage[math-style=ISO,bold-style=ISO]{unicode-math}
\setmathfont{FiraMath-Regular.otf}[Path=../../../fonts/]
\setmathfont{PlusJakartaSans-Medium.ttf}[Path=../../../fonts/,range={up/num,up/{Latin,latin},\mathup}]
\usepackage{icomma} % virgule décimale sans espace en mode math
% Caractères mathématiques Unicode absents des polices de texte (Plus Jakarta,
% Archivo Black) : rendus par leur équivalent mathématique, en texte comme en
% formule — sinon ils s'affichent en carré vide (ex. « Arithmétique dans ℤ »).
\usepackage{newunicodechar}
\newunicodechar{ℝ}{\ensuremath{\mathbb{R}}}
\newunicodechar{ℤ}{\ensuremath{\mathbb{Z}}}
\newunicodechar{ℂ}{\ensuremath{\mathbb{C}}}
\newunicodechar{ℕ}{\ensuremath{\mathbb{N}}}
\newunicodechar{ℚ}{\ensuremath{\mathbb{Q}}}
\newunicodechar{𝔻}{\ensuremath{\mathbb{D}}}
\newunicodechar{α}{\ensuremath{\alpha}}
\newunicodechar{β}{\ensuremath{\beta}}
\newunicodechar{γ}{\ensuremath{\gamma}}
\newunicodechar{δ}{\ensuremath{\delta}}
\newunicodechar{ε}{\ensuremath{\varepsilon}}
\newunicodechar{θ}{\ensuremath{\theta}}
\newunicodechar{λ}{\ensuremath{\lambda}}
\newunicodechar{μ}{\ensuremath{\mu}}
\newunicodechar{σ}{\ensuremath{\sigma}}
\newunicodechar{τ}{\ensuremath{\tau}}
\newunicodechar{φ}{\ensuremath{\varphi}}
\newunicodechar{ψ}{\ensuremath{\psi}}
\newunicodechar{ω}{\ensuremath{\omega}}
\newunicodechar{Σ}{\ensuremath{\Sigma}}
% majuscules produites par \MakeUppercase dans les titres (α-aminés -> Α)
\newunicodechar{Α}{\ensuremath{\alpha}}
\newunicodechar{Β}{\ensuremath{\beta}}
\newunicodechar{Γ}{\ensuremath{\Gamma}}
\newunicodechar{Δ}{\ensuremath{\Delta}}
\newunicodechar{Ε}{\ensuremath{\varepsilon}}
\newunicodechar{Θ}{\ensuremath{\Theta}}
\newunicodechar{Λ}{\ensuremath{\Lambda}}
\newunicodechar{Μ}{\ensuremath{\mu}}
\newunicodechar{Τ}{\ensuremath{\tau}}
\newunicodechar{Φ}{\ensuremath{\Phi}}
\newunicodechar{Ψ}{\ensuremath{\Psi}}
\newunicodechar{Ω}{\ensuremath{\Omega}}
\newunicodechar{↦}{\ensuremath{\mapsto}}
\newunicodechar{∈}{\ensuremath{\in}}
\newunicodechar{∉}{\ensuremath{\notin}}
\newunicodechar{∧}{\ensuremath{\wedge}}
\newunicodechar{⇔}{\ensuremath{\Leftrightarrow}}
\newunicodechar{⟺}{\ensuremath{\Longleftrightarrow}}
\newunicodechar{⇒}{\ensuremath{\Rightarrow}}
\newunicodechar{⟹}{\ensuremath{\Longrightarrow}}
\newunicodechar{∘}{\ensuremath{\circ}}
\newunicodechar{∪}{\ensuremath{\cup}}
\newunicodechar{∩}{\ensuremath{\cap}}
\newunicodechar{≡}{\ensuremath{\equiv}}
\newunicodechar{⊂}{\ensuremath{\subset}}
\newunicodechar{⊥}{\ensuremath{\perp}}
\newunicodechar{∃}{\ensuremath{\exists}}
\newunicodechar{∀}{\ensuremath{\forall}}
\newunicodechar{∛}{\ensuremath{\sqrt[3]{\,}}}
\newunicodechar{∜}{\ensuremath{\sqrt[4]{\,}}}
\newunicodechar{⃗}{\ensuremath{{}^{\to}}}
\newunicodechar{ⁿ}{\ifmmode^{n}\else\textsuperscript{\ensuremath{n}}\fi}
\newunicodechar{ˣ}{\ifmmode^{x}\else\textsuperscript{\ensuremath{x}}\fi}
\newunicodechar{ᵇ}{\ifmmode^{b}\else\textsuperscript{\ensuremath{b}}\fi}
\newunicodechar{ʳ}{\ifmmode^{r}\else\textsuperscript{\ensuremath{r}}\fi}
\newunicodechar{ᵘ}{\ifmmode^{u}\else\textsuperscript{\ensuremath{u}}\fi}
\newunicodechar{ʸ}{\ifmmode^{y}\else\textsuperscript{\ensuremath{y}}\fi}
\newunicodechar{ᵃ}{\ifmmode^{a}\else\textsuperscript{\ensuremath{a}}\fi}
\newunicodechar{ᵉ}{\ifmmode^{e}\else\textsuperscript{\ensuremath{e}}\fi}
\newunicodechar{ᵏ}{\ifmmode^{k}\else\textsuperscript{\ensuremath{k}}\fi}
\newunicodechar{ᵖ}{\ifmmode^{p}\else\textsuperscript{\ensuremath{p}}\fi}
\newunicodechar{ᴺ}{\ifmmode^{N}\else\textsuperscript{\ensuremath{N}}\fi}
\newunicodechar{ᶜ}{\ifmmode^{c}\else\textsuperscript{\ensuremath{c}}\fi}
\newunicodechar{ʰ}{\ifmmode^{h}\else\textsuperscript{\ensuremath{h}}\fi}
\newunicodechar{ᵠ}{\ifmmode^{q}\else\textsuperscript{\ensuremath{q}}\fi}
\newunicodechar{ₙ}{\ifmmode_{n}\else\textsubscript{\ensuremath{n}}\fi}
\newunicodechar{ₐ}{\ifmmode_{a}\else\textsubscript{\ensuremath{a}}\fi}
\newunicodechar{ᵢ}{\ifmmode_{i}\else\textsubscript{\ensuremath{i}}\fi}
\newunicodechar{ᵣ}{\ifmmode_{r}\else\textsubscript{\ensuremath{r}}\fi}
\newunicodechar{ₖ}{\ifmmode_{k}\else\textsubscript{\ensuremath{k}}\fi}
\newunicodechar{ᵦ}{\ifmmode_{\beta}\else\textsubscript{\ensuremath{\beta}}\fi}
\newunicodechar{ₓ}{\ifmmode_{x}\else\textsubscript{\ensuremath{x}}\fi}
\newfontfamily\popdisplay{ArchivoBlack-Regular.ttf}[Path=../../../fonts/]
\newfontfamily\poplogo{SpaceGrotesk-Bold.ttf}[Path=../../../fonts/]
\newfontfamily\popsemi{PlusJakartaSans-SemiBold.ttf}[Path=../../../fonts/]
\newfontfamily\popxbold{PlusJakartaSans-ExtraBold.ttf}[Path=../../../fonts/]
\newfontfamily\popxboldls{PlusJakartaSans-ExtraBold.ttf}[Path=../../../fonts/,LetterSpace=3.0]

\setlist[enumerate]{leftmargin=1.7em,itemsep=5pt,topsep=4pt}
\setlist[itemize]{leftmargin=1.7em,itemsep=4pt,topsep=4pt,label={\color{poporange}\textbullet}}
\setlength{\parindent}{0pt}
\setlength{\parskip}{5pt}
\renewcommand{\arraystretch}{1.25}

% pgfplots : style Pop par défaut
\pgfplotsset{
  every axis/.append style={axis line style={popink,line width=1pt},tick style={popink},
    grid style={popline,line width=0.5pt},label style={font=\small},tick label style={font=\footnotesize}},
  popcurve/.style={poporange,line width=1.8pt,no markers,smooth},
}
% Même style disponible dans un \draw TikZ simple (hors pgfplots)
\tikzset{popcurve/.style={poporange,line width=1.8pt}}
\tikzset{popgrid/.style={popline,very thin}}
\colorlet{popcurve}{poporange} % le nom du style est parfois utilisé comme couleur

% ============================================================
% Pill / squiggle
% ============================================================
\newcommand{\poppill}[3]{%
  \tikz[baseline=-0.55ex]\node[draw=popink,line width=1.2pt,rounded corners=2.6pt,%
    fill=#2,inner xsep=6.5pt,inner ysep=3pt]{%
    \popxboldls\fontsize{7}{7}\selectfont\color{#3}\MakeUppercase{#1}};%
}
\newcommand{\popsquiggle}[1]{%
  \tikz[baseline=-0.4ex]{
    \draw[#1,line width=2.6pt,line cap=round]
      plot [smooth] coordinates {(0,0.09) (0.34,0.01) (0.68,0.21) (1.02,0.01) (1.36,0.10)};
  }%
}
% Mot-clé surligné (marqueur orange sous le texte)
\newcommand{\cle}[1]{{\popxbold\color{popdark}#1}}

% ============================================================
% En-tête / pied
% ============================================================
\pagestyle{fancy}
\fancyhf{}
\renewcommand{\headrulewidth}{0pt}
\renewcommand{\footrulewidth}{0pt}
\lhead{\raisebox{-0.15ex}{\poplogo\fontsize{13}{13}\selectfont\color{popink}OnBuch\color{poporange}+}}
\rhead{\poppill{\DOCMATIERE\ \textbullet\ \DOCNIVEAU\ \textbullet\ Leçon \DOCLECON}{white}{popink}}
\lfoot{\popxbold\fontsize{7.6}{7.6}\selectfont\color{popmuted}\MakeUppercase{Cours signé OnBuch+}\ \textbullet\ {\popsemi\DOCTITRE}}
\rfoot{\tikz[baseline=-0.6ex]\node[rounded corners=8.5pt,fill=popink,minimum height=17pt,minimum width=17pt,inner xsep=5pt,inner sep=0]{\popxbold\fontsize{8}{8}\selectfont\color{popcream}\thepage\,/\,\pageref*{LastPage}};}

% ============================================================
% Titres
% ============================================================
\newcommand{\chaptitre}[2]{%
  \begin{tcolorbox}[enhanced,colback=poporange,colframe=popink,boxrule=2.6pt,arc=11pt,
      drop shadow=popink,left=15pt,right=15pt,top=12pt,bottom=11pt,before skip=3pt,after skip=18pt]
    \poppill{#2}{popink}{popcream}\par\vspace{9pt}
    {\raggedright\hyphenpenalty=10000\exhyphenpenalty=10000\popdisplay\fontsize{22}{24}\selectfont\color{popcream}\MakeUppercase{#1}\par}
  \end{tcolorbox}
}
% Section numérotée : pastille orange avec le numéro + titre Archivo
\newcounter{popsec}
\newcommand{\coursec}[1]{%
  \stepcounter{popsec}\setcounter{popsub}{0}%
  \par\vspace{16pt}\noindent
  \begin{minipage}{\linewidth}
  \tikz[baseline=-0.7ex]\node[draw=popink,line width=1.6pt,fill=poporange,rounded corners=4pt,minimum size=22pt,inner sep=0]{\popdisplay\fontsize{12}{12}\selectfont\color{popcream}\thepopsec};%
  \hspace{8pt}{\popdisplay\fontsize{15}{17}\selectfont\color{popink}\MakeUppercase{#1}}\par
  \vspace{4pt}\noindent{\color{poporange}\rule{36pt}{3.6pt}}
  \end{minipage}\par\nopagebreak\vspace{8pt}
}
\newcounter{popsub}
\newcommand{\courssub}[1]{%
  \stepcounter{popsub}%
  \par\vspace{9pt}\noindent{\popxbold\fontsize{11.5}{13}\selectfont\color{popink}{\color{poporange}\thepopsec.\thepopsub}\hspace{6pt}#1}\par\nopagebreak\vspace{4pt}
}

% ============================================================
% Boîtes pédagogiques — même recette : fond blanc, bordure épaisse colorée,
% ombre dure encre, badge plein. Titre optionnel [#1].
% ============================================================
\tcbset{popbase/.style={breakable,enhanced,colback=white,boxrule=2.3pt,arc=9pt,
  shadow={3pt}{-3pt}{0pt}{popink},left=14pt,right=14pt,top=11pt,bottom=11pt,before skip=11pt,after skip=15pt,
  fontupper=\popsemi\fontsize{10.6}{14.5}\selectfont\color{popink},
  fontlower=\popsemi\fontsize{10.6}{14.5}\selectfont\color{popink}}}
% Coche dessinée (le glyphe ✓ n'existe pas dans les polices du document)
\newcommand{\popcheck}[1]{\tikz[baseline=-0.5ex]\draw[#1,line width=1.6pt,line cap=round,line join=round](-0.13,0.0)--(-0.04,-0.09)--(0.13,0.1);}
\providecommand{\coche}{\popcheck{popgreen}}
\providecommand{\resultat}[1]{\par\smallskip\noindent\fcolorbox{popgreen}{popgreenL}{\parbox{\dimexpr\linewidth-2\fboxsep-2\fboxrule\relax}{\textbf{Résultat.} #1}}\par\smallskip}
\newcommand{\popboxhead}[3]{\poppill{#1}{#2}{white}\ifx\relax#3\relax\else\hspace{6pt}{\popxbold\fontsize{10.6}{12}\selectfont\color{popink}#3}\fi\par\vspace{7pt}}

\newtcolorbox{aretenir}[1][]{popbase,colframe=popgreen,before upper={\popboxhead{À retenir}{popgreen}{#1}}}
\newtcolorbox{exemplebox}[1][]{popbase,colframe=poporange,before upper={\popboxhead{Exemple}{poporange}{#1}}}
\newtcolorbox{definition}[1][]{popbase,colframe=popblue,before upper={\popboxhead{Définition}{popblue}{#1}}}
\newtcolorbox{propriete}[1][]{popbase,colframe=poppurple,before upper={\popboxhead{Propriété}{poppurple}{#1}}}
\newtcolorbox{methode}[1][]{popbase,colframe=popgold,before upper={\popboxhead{Méthode}{popgold}{#1}}}
\newtcolorbox{attention}[1][]{popbase,colframe=poppink,before upper={\popboxhead{Attention}{poppink}{#1}}}
\newtcolorbox{experience}[1][]{popbase,colframe=popblue,colback=popblueL,before upper={\popboxhead{Expérience}{popblue}{#1}}}
% « Le savais-tu ? » : carte violette pleine (contexte, culture, Cameroun)
\newtcolorbox{savaistu}[1][]{popbase,colback=poppurple,colframe=popink,
  fontupper=\popsemi\fontsize{10.4}{14.2}\selectfont\color{white},
  before upper={\poppill{Le savais-tu\,?}{white}{poppurple}\ifx\relax#1\relax\else\hspace{6pt}{\popxbold\color{white}#1}\fi\par\vspace{7pt}}}
% Exercice résolu : énoncé en haut, \tcblower puis la solution sur fond crème
\newcounter{popexo}\newcounter{popexr}
\newtcolorbox{exoresolu}[1][]{popbase,colframe=poporange,colbacklower=popcream,
  segmentation style={popink,line width=1.2pt,dashed},
  before upper={\stepcounter{popexr}\popboxhead{Exercice résolu \thepopexr}{poporange}{#1}},
  before lower={\poppill{Solution}{popink}{popcream}\par\vspace{6pt}}}
% Exercice (non résolu en place) — la correction est dans la partie Corrigés
\newtcolorbox{exercice}[1][]{popbase,colframe=popink,
  before upper={\stepcounter{popexo}\popboxhead{Exercice \thepopexo}{popink}{#1}}}
% Corrigé
\newtcolorbox{corrige}[1][]{popbase,colframe=popgreen,colback=popgreenL,
  before upper={\popboxhead{Corrigé}{popgreen}{#1}}}
% Objectifs / prérequis / bilan
\newtcolorbox{objectifs}{popbase,colframe=popink,colback=poporangeL,
  before upper={\popboxhead{Objectifs de la leçon}{popink}{}}}
\newtcolorbox{prerequis}{popbase,colframe=popink,colback=popblueL,
  before upper={\popboxhead{Prérequis}{popblue}{}}}
\newtcolorbox{bilan}{popbase,colframe=popink,colback=white,boxrule=2.8pt,
  before upper={\popboxhead{Fiche bilan}{poporange}{L'essentiel à connaître pour l'examen}}}
% Figure encadrée avec légende
\newtcolorbox{popfigure}[1][]{enhanced,breakable=false,colback=white,colframe=popline,boxrule=1.4pt,arc=8pt,
  left=10pt,right=10pt,top=10pt,bottom=8pt,before skip=10pt,after skip=12pt,halign=center,
  lower separated=false,halign lower=center,
  fontlower=\popsemi\fontsize{9}{11.5}\selectfont\color{popmuted},
  #1}
\newcommand{\legende}[1]{\par\vspace{6pt}{\popsemi\fontsize{9}{11.5}\selectfont\color{popmuted}#1}}

% ============================================================
% Signature OnBuch+
% ============================================================
\newcommand{\poplogoinline}[1]{{\poplogo\fontsize{#1}{#1}\selectfont\color{popink}OnBuch\color{poporange}+}}
\newcommand{\signatureonbuch}{%
  \begin{tcolorbox}[enhanced,colback=popink,colframe=popink,boxrule=2.2pt,arc=10pt,
      drop shadow=poporange,left=14pt,right=14pt,top=10pt,bottom=10pt,before skip=0pt,after skip=0pt]
    \tikz[baseline=-0.7ex]\node[circle,fill=popgreen,draw=popcream,line width=1.3pt,minimum size=22pt,inner sep=0]{\popcheck{white}};%
    \hspace{9pt}\begin{minipage}[c]{0.62\linewidth}
      {\popxbold\fontsize{12}{13}\selectfont\color{popcream}Signé\ \textbullet\ {\poplogo OnBuch\color{poporange}+}}\par\vspace{2pt}
      {\raggedright\popsemi\fontsize{8}{10.5}\selectfont\color{popline}\MakeUppercase{Équipe pédagogique OnBuch+}\\\MakeUppercase{Conforme au programme officiel MINESEC}\par}
    \end{minipage}\hfill
    {\poplogo\fontsize{22}{22}\selectfont\color{popcream}OnBuch\color{poporange}+}
  \end{tcolorbox}%
}

% ============================================================
% Page de garde
% #1 titre · #2 matière · #3 niveau · #4 module · #5 n° leçon · #6 durée ·
% #7 description / objectifs (texte ou itemize)
% ============================================================
\newcommand{\pagegarde}[7]{%
  \thispagestyle{empty}
  \newgeometry{a4paper,margin=1.7cm,top=1.6cm,bottom=1.4cm}%
  \begin{tikzpicture}[remember picture,overlay]
    \fill[poporange!18] ([xshift=-3.2cm,yshift=-3.6cm]current page.north east) circle (4.6cm);
    \fill[poppurple!14] ([xshift=2.4cm,yshift=7.6cm]current page.south west) circle (3.8cm);
  \end{tikzpicture}%
  {\poplogo\fontsize{30}{30}\selectfont\color{popink}OnBuch\color{poporange}+}\hfill
  \raisebox{0.6ex}{\poppill{Cours signé \textbullet\ Premium}{popink}{popcream}}\par
  \vspace{1pt}\popsquiggle{poppurple}\par
  \vspace{8pt}
  \begin{tcolorbox}[enhanced,colback=poporange,colframe=popink,boxrule=2.8pt,arc=13pt,
      drop shadow=popink,left=18pt,right=18pt,top=12pt,bottom=13pt,after skip=10pt]
    \poppill{#2 \textbullet\ #3}{popink}{popcream}\hspace{5pt}\poppill{#4}{white}{popink}\par\vspace{8pt}
    {\popdisplay\fontsize{14}{14}\selectfont\color{popink}LEÇON #5}\par\vspace{4pt}
    {\raggedright\hyphenpenalty=10000\exhyphenpenalty=10000\popdisplay\fontsize{25}{27}\selectfont\color{popcream}\MakeUppercase{#1}\par}
    \vspace{8pt}
    {\popxbold\fontsize{9.5}{9.5}\selectfont\color{popink}\MakeUppercase{Durée conseillée : #6}}
  \end{tcolorbox}
  \begin{tcolorbox}[enhanced,colback=white,colframe=popink,boxrule=2.2pt,arc=10pt,
      drop shadow=popink,left=16pt,right=16pt,top=10pt,bottom=10pt,after skip=8pt]
    \poppill{Dans cette leçon}{poppurple}{white}\par\vspace{5pt}
    \popsemi\fontsize{9.3}{12.6}\selectfont\color{popink}#7
  \end{tcolorbox}
  \noindent\begin{minipage}[t]{0.487\linewidth}
    \begin{tcolorbox}[enhanced,colback=popgreen,colframe=popink,boxrule=2.2pt,arc=10pt,
        drop shadow=popink,left=11pt,right=11pt,top=10pt,bottom=10pt,height=3.1cm,valign=center]
      \tcbox[colback=white,colframe=popink,boxrule=1.3pt,arc=5pt,boxsep=0pt,nobeforeafter,
        left=3pt,right=3pt,top=3pt,bottom=3pt,valign=center]{\qrcode[height=1.9cm]{https://play.google.com/store/apps/details?id=com.luvvix.onbuch&hl=fr}}%
      \hspace{8pt}\begin{minipage}[c]{0.5\linewidth}
        {\popdisplay\fontsize{11}{12.5}\selectfont\color{white}\MakeUppercase{Télécharge OnBuch}}\par\vspace{4pt}
        {\raggedright\popsemi\fontsize{8.4}{11}\selectfont\color{white}Gratuit \textbullet\ Google Play\\Tuteur IA, annales, concours, résultats\par}
      \end{minipage}
    \end{tcolorbox}
  \end{minipage}\hfill
  \begin{minipage}[t]{0.487\linewidth}
    \begin{tcolorbox}[enhanced,colback=poppurple,colframe=popink,boxrule=2.2pt,arc=10pt,
        drop shadow=popink,left=11pt,right=11pt,top=10pt,bottom=10pt,height=3.1cm,valign=center]
      \tikz[baseline=-0.7ex]\node[circle,fill=white,draw=popink,line width=1.3pt,minimum size=1.6cm,inner sep=0]{\popdisplay\fontsize{24}{24}\selectfont\color{poppurple}+};%
      \hspace{8pt}\begin{minipage}[c]{0.6\linewidth}
        {\popdisplay\fontsize{11}{12.5}\selectfont\color{white}\MakeUppercase{Passe à OnBuch+}}\par\vspace{4pt}
        {\raggedright\popsemi\fontsize{8.4}{11}\selectfont\color{white}Tous les cours signés, Tuteur IA illimité, fiches PDF — onglet OnBuch+ de l'app\par}
      \end{minipage}
    \end{tcolorbox}
  \end{minipage}
  \vspace{8pt}
  \popcontenu
  \vfill
  \signatureonbuch
  \restoregeometry
  \newpage
}

% Rangée « Ce cours contient » (page de garde)
\newcommand{\poptile}[3]{% #1 couleur #2 gros texte #3 légende
  \begin{tcolorbox}[enhanced,colback=#1,colframe=popink,boxrule=2pt,arc=9pt,shadow={2.5pt}{-2.5pt}{0pt}{popink},
      left=6pt,right=6pt,top=9pt,bottom=9pt,halign=center,valign=center,height=1.75cm,width=\linewidth,nobeforeafter]
    {\popdisplay\fontsize{12.5}{13}\selectfont\color{white}\MakeUppercase{#2}}\par\vspace{3pt}
    {\centering\popsemi\fontsize{7.8}{9.5}\selectfont\color{white}#3\par}
  \end{tcolorbox}}
\newcommand{\popcontenu}{%
  \par\noindent{\popxboldls\fontsize{7.5}{7.5}\selectfont\color{popmuted}CE COURS SIGNÉ CONTIENT}\par\vspace{7pt}
  \noindent\begin{minipage}[t]{0.235\linewidth}\poptile{popblue}{Cours}{complet, illustré, pas à pas}\end{minipage}\hfill
  \begin{minipage}[t]{0.235\linewidth}\poptile{poporange}{Méthodes}{et exercices résolus}\end{minipage}\hfill
  \begin{minipage}[t]{0.235\linewidth}\poptile{poppink}{Exercices}{type examen + corrigés}\end{minipage}\hfill
  \begin{minipage}[t]{0.235\linewidth}\poptile{popgreen}{Fiche bilan}{l'essentiel à retenir}\end{minipage}\par}

% Sommaire
\newtcolorbox{sommaire}{enhanced,colback=white,colframe=popink,boxrule=2.2pt,arc=10pt,
  drop shadow=popink,left=18pt,right=18pt,top=13pt,bottom=13pt,before skip=6pt,after skip=20pt,
  before upper={\poppill{Sommaire}{poppurple}{white}\par\vspace{9pt}\popsemi\fontsize{10.6}{14.5}\selectfont\color{popink}}}

% Signature de fin de cours — poussée tout en bas de la dernière page
\newcommand{\findecours}{%
  \par\vfill\nopagebreak
  \begin{center}\popsquiggle{poporange}\end{center}\vspace{4pt}
  \signatureonbuch
}
\AtBeginDocument{\def\llbracket{\mathopen{[\mkern-3.5mu[}}\def\rrbracket{\mathclose{]\mkern-3.5mu]}}} % intervalles d'entiers (glyphe ⟦ absent de la police)
% Glyphes absents de Fira Math : redéfinis à partir de glyphes disponibles
\AtBeginDocument{%
  \def\setminus{\mathbin{\backslash}}\let\smallsetminus\setminus
  \def\vdots{\mathord{\vbox{\baselineskip3.2pt\lineskiplimit0pt\kern4pt\hbox{.}\hbox{.}\hbox{.}}}}%
  \def\ddots{\mathinner{\mkern1mu\raise6.5pt\hbox{.}\mkern2mu\raise3.6pt\hbox{.}\mkern2mu\raise0.7pt\hbox{.}\mkern1mu}}%
  \def\triangle{\mathord{\scalebox{0.9}{$\symup{\Delta}$}}}%
  \def\nabla{\mathord{\raisebox{\depth}{\scalebox{1}[-1]{$\symup{\Delta}$}}}}%
  \def\checkmark{\popcheck{popdark}}%
  \def\star{\mathbin{\ast}}\def\bigstar{\mathord{\ast}}%
  \def\diamond{\mathbin{\scalebox{0.6}{$\Diamond$}}}%
  \def\bigcup{\mathop{\mathchoice{\raisebox{-0.3em}{\scalebox{1.7}{$\cup$}}}{\scalebox{1.25}{$\cup$}}{\cup}{\cup}}\displaylimits}%
  \def\bigcap{\mathop{\mathchoice{\raisebox{-0.3em}{\scalebox{1.7}{$\cap$}}}{\scalebox{1.25}{$\cap$}}{\cap}{\cap}}\displaylimits}%
  \def\lesseqqgtr{\mathrel{\lessgtr}}\def\bigoplus{\mathop{\oplus}\displaylimits}%
}
\providecommand{\ssi}{si et seulement si} % abréviation fréquente
\AtBeginDocument{\providecommand{\wideparen}{\overparen}} % arc de cercle

%% FIGFIX-BEGIN v3 — correctifs globaux des figures (docs/onbuch-generation/tools/figfix)
\usepackage{adjustbox}\usepackage{etoolbox}
% toute figure TikZ/pgfplots plus large que la ligne est réduite à la largeur disponible
\BeforeBeginEnvironment{tikzpicture}{\begin{adjustbox}{max width=\linewidth}}
\AfterEndEnvironment{tikzpicture}{\end{adjustbox}}
% légendes pgfplots lisibles par-dessus les courbes
\pgfplotsset{every axis/.append style={legend style={fill=white,fill opacity=0.92,text opacity=1,draw=black!15,inner sep=3pt}}}
% étiquettes d'arêtes (syntaxe edge["..."]) avec fond blanc
\tikzset{every edge quotes/.append style={fill=white,inner sep=1.2pt}}
% bulles de cartes mentales : texte à la ligne dans la bulle, bulle un peu plus grande
\tikzset{every concept/.append style={text width=2.0cm,align=center,minimum size=2.3cm,inner sep=2pt}}
%% FIGFIX-END
\begin{document}

\pagegarde{Les lentilles minces}{Physique}{1ère E}{Module 3 : Optique}{N07}{8 h}{Cette leçon étudie les lentilles minces : définition, éléments caractéristiques (foyers, vergence, distance focale), conditions de Gauss et constructions géométriques des images. Elle établit et exploite les formules de position et de grandissement, le théorème des vergences, et présente des applications concrètes (loupe, œil, appareil photographique).
\vspace{4pt}
\begin{multicols}{2}\small
\begin{itemize}
\item À la fin de cette leçon, je sais définir une lentille mince et la reconnaître par son aspect (convergente ou divergente)
\item Je sais schématiser une lentille et situer son centre optique, ses foyers et ses plans focaux
\item Je sais énoncer et réaliser les conditions de Gauss
\item Je sais construire géométriquement l'image d'un objet donnée par une lentille mince
\item Je sais utiliser les formules de position et de grandissement pour résoudre des problèmes
\item Je sais appliquer le théorème des vergences à un système de deux lentilles accolées
\end{itemize}\end{multicols}}

\begin{sommaire}
\begin{multicols}{2}
\begin{enumerate}
\item Définition, classification et éléments caractéristiques
\item Conditions de Gauss et propriétés optiques
\item Construction géométrique des images
\item Formules de position et de grandissement
\item Association de lentilles : théorème des vergences
\item Applications des lentilles minces
\item Activité d'intégration
\item Méthodes et exercices résolus
\item Exercices
\item Corrigés des exercices
\item Fiche bilan
\end{enumerate}
\end{multicols}
\end{sommaire}

\begin{objectifs}
\begin{itemize}
\item À la fin de cette leçon, je sais définir une lentille mince et la reconnaître par son aspect (convergente ou divergente)
\item Je sais schématiser une lentille et situer son centre optique, ses foyers et ses plans focaux
\item Je sais énoncer et réaliser les conditions de Gauss
\item Je sais construire géométriquement l'image d'un objet donnée par une lentille mince
\item Je sais utiliser les formules de position et de grandissement pour résoudre des problèmes
\item Je sais appliquer le théorème des vergences à un système de deux lentilles accolées
\item Je sais construire l'image d'un objet à travers un système de deux lentilles
\item Je sais déterminer expérimentalement la vergence d'une lentille et vérifier la formule de position
\item Je sais citer et décrire quelques applications des lentilles (loupe, œil, appareil photo)
\end{itemize}
\end{objectifs}

\begin{prerequis}
\begin{itemize}
\item Propagation rectiligne de la lumière et modèle du rayon lumineux (Seconde)
\item Lois de la réfraction de Snell-Descartes (leçon sur la réfraction)
\item Notion d'image réelle et d'image virtuelle
\item Calcul algébrique sur les grandeurs algébriques et les inverses (1/x)
\item Lecture et exploitation de schémas géométriques à l'échelle
\end{itemize}
\end{prerequis}

\newpage

\coursec{Définition, classification et éléments caractéristiques}

Moussa, élève en Première C à Yaoundé, remarque que sa grand-mère plisse les yeux pour lire le journal, mais qu'avec les lunettes prêtées par un voisin, les lettres deviennent nettes. Intrigué, il rapproche une loupe d'un mur et obtient, à une certaine distance, l'image renversée et nette de la fenêtre. Comment un simple morceau de verre bombé peut-il redresser la vue ou projeter une image ? Où faut-il placer exactement la lentille pour obtenir une image nette ?

Pour répondre à ces questions, nous allons étudier dans cette leçon les \cle{lentilles minces} : leur définition, leurs éléments caractéristiques, leurs propriétés optiques, les formules qui permettent de prévoir la position et la taille des images, puis leurs applications, de la loupe à l'appareil photo. Commençons par le commencement : qu'est-ce qu'une lentille, et comment la reconnaître ?

\courssub{Définition d'une lentille mince}

\textbf{Observation.} Une loupe, un verre de lunette, l'objectif d'un appareil photo ou d'un projecteur utilisé au cinéma Canal Olympia de Yaoundé : tous ces objets ont un point commun. Ce sont des morceaux de verre (ou de plastique) transparents, dont au moins une face n'est pas plane mais bombée, comme une calotte sphérique. La lumière qui les traverse est \emph{déviée} par réfraction à chaque face : c'est cette double réfraction qui permet de diriger les rayons lumineux là où on le souhaite.

\begin{definition}[Lentille]
Une \cle{lentille} est un milieu transparent (verre, plastique), homogène et isotrope, limité par deux \cle{dioptres} dont l'un au moins est \emph{sphérique} (l'autre pouvant être plan ou sphérique).
\end{definition}

Rappel : un \cle{dioptre} est une surface de séparation entre deux milieux transparents d'indices différents. Ici, la lumière traverse successivement le dioptre d'entrée (air $\to$ verre) puis le dioptre de sortie (verre $\to$ air). À chaque traversée, la loi de la réfraction de Snell-Descartes s'applique : le rayon est dévié vers la normale en entrant dans le verre (milieu plus réfringent), puis s'en écarte en ressortant.

\begin{definition}[Lentille mince]
Une lentille est dite \cle{mince} lorsque son épaisseur $e$ au centre est très petite devant les rayons de courbure $R_1$ et $R_2$ de ses deux faces :
\[ e \ll R_1 \quad \text{et} \quad e \ll R_2. \]
Dans ce cas, on assimile la lentille à un plan : tout se passe comme si les deux faces se touchaient en un point $O$ appelé \cle{centre optique}. Les sommets $S_1$ et $S_2$ des deux dioptres sont confondus avec $O$.
\end{definition}

\begin{attention}[Pourquoi l'hypothèse « mince » est-elle essentielle ?]
Toutes les formules de cette leçon (formule de position, grandissement, vergences) ne sont valables que pour des lentilles \emph{minces} utilisées dans les \emph{conditions de Gauss} (que nous verrons à la section suivante). Une lentille épaisse, comme la lentille d'un phare de voiture, n'obéit pas à ces formules simples. Au Probatoire, sauf mention contraire, toute lentille étudiée est supposée mince.
\end{attention}

\courssub{Classification géométrique des lentilles}

\textbf{Observation.} Passe le doigt sur le bord d'une loupe : le bord est \emph{plus mince} que le centre. Fais de même avec le verre correcteur d'un myope : le bord est \emph{plus épais} que le centre. Ce simple test au toucher suffit à classer toutes les lentilles en deux grandes familles.

\begin{definition}[Les deux familles de lentilles]
\begin{itemize}
\item Les lentilles à \cle{bords minces} : le centre est plus épais que les bords. Ce sont les lentilles \textbf{biconvexe}, \textbf{plan-convexe} et \textbf{ménisque convergent}. Elles sont \emph{convergentes}.
\item Les lentilles à \cle{bords épais} : le centre est plus mince que les bords. Ce sont les lentilles \textbf{biconcave}, \textbf{plan-concave} et \textbf{ménisque divergent}. Elles sont \emph{divergentes}.
\end{itemize}
\end{definition}

\begin{popfigure}
\begin{center}
\begin{tikzpicture}[scale=0.9]
% --- Bords minces ---
% Biconvexe
\begin{scope}[shift={(0,0)}]
\fill[popblueL] (0,0) circle[x radius=0.55, y radius=1.1];
\draw[popblue, thick] (0,0) circle[x radius=0.55, y radius=1.1];
\node[below, popdark] at (0,-1.25) {\small biconvexe};
\end{scope}
% Plan-convexe
\begin{scope}[shift={(2.2,0)}]
\fill[popblueL] (0,-1.1) -- (0,1.1) arc[start angle=90, end angle=-90, x radius=0.55, y radius=1.1] -- cycle;
\draw[popblue, thick] (0,-1.1) -- (0,1.1) arc[start angle=90, end angle=-90, x radius=0.55, y radius=1.1] -- cycle;
\node[below, popdark] at (0.25,-1.25) {\small plan-convexe};
\end{scope}
% Ménisque convergent
\begin{scope}[shift={(4.4,0)}]
\fill[popblueL] (0,-1.1) arc[start angle=-90, end angle=90, x radius=0.5, y radius=1.1]
arc[start angle=90, end angle=-90, x radius=0.18, y radius=1.1] -- cycle;
\draw[popblue, thick] (0,-1.1) arc[start angle=-90, end angle=90, x radius=0.5, y radius=1.1]
arc[start angle=90, end angle=-90, x radius=0.18, y radius=1.1] -- cycle;
\node[below, popdark] at (0.35,-1.25) {\small ménisque convergent};
\end{scope}
\node[popblue, font=\bfseries] at (2.2,1.7) {\small Bords minces (convergentes)};
% --- Bords épais ---
% Biconcave
\begin{scope}[shift={(7.2,0)}]
\fill[poporangeL] (-0.5,-1.1) rectangle (0.5,1.1);
\fill[white] (-0.5,0) circle[x radius=0.35, y radius=1.1];
\fill[white] (0.5,0) circle[x radius=0.35, y radius=1.1];
\draw[poporange, thick] (-0.5,-1.1) arc[start angle=-90, end angle=90, x radius=0.35, y radius=1.1];
\draw[poporange, thick] (0.5,-1.1) arc[start angle=270, end angle=90, x radius=0.35, y radius=1.1];
\node[below, popdark] at (0,-1.25) {\small biconcave};
\end{scope}
% Plan-concave
\begin{scope}[shift={(9.4,0)}]
\fill[poporangeL] (0,-1.1) rectangle (0.5,1.1);
\fill[white] (0,0) circle[x radius=0.35, y radius=1.1];
\draw[poporange, thick] (0,-1.1) arc[start angle=-90, end angle=90, x radius=0.35, y radius=1.1];
\draw[poporange, thick] (0.5,-1.1) -- (0.5,1.1);
\node[below, popdark] at (0.25,-1.25) {\small plan-concave};
\end{scope}
% Ménisque divergent
\begin{scope}[shift={(11.6,0)}]
\fill[poporangeL] (0,-1.1) arc[start angle=-90, end angle=90, x radius=0.5, y radius=1.1]
-- (0.32,1.1) arc[start angle=90, end angle=-90, x radius=0.85, y radius=1.1] -- cycle;
\draw[poporange, thick] (0,-1.1) arc[start angle=-90, end angle=90, x radius=0.5, y radius=1.1];
\draw[poporange, thick] (0.32,1.1) arc[start angle=90, end angle=-90, x radius=0.85, y radius=1.1];
\node[below, popdark] at (0.15,-1.25) {\small ménisque divergent};
\end{scope}
\node[poporange, font=\bfseries] at (9.4,1.7) {\small Bords épais (divergentes)};
\end{tikzpicture}
\end{center}
\legende{Les six profils de lentilles minces : à gauche, les trois formes à bords minces (convergentes) ; à droite, les trois formes à bords épais (divergentes).}
\end{popfigure}

\begin{methode}[Reconnaître une lentille convergente ou divergente par son aspect]
\begin{enumerate}
\item \textbf{Test au toucher} : passe un doigt du centre vers le bord de la lentille.
\begin{itemize}
\item Bords plus \emph{minces} que le centre $\Rightarrow$ lentille \textbf{convergente}.
\item Bords plus \emph{épais} que le centre $\Rightarrow$ lentille \textbf{divergente}.
\end{itemize}
\item \textbf{Test sur un texte} : pose la lentille sur une ligne de texte (journal, manuel) et soulève-la légèrement.
\begin{itemize}
\item Les lettres paraissent \emph{grossies} $\Rightarrow$ lentille \textbf{convergente}.
\item Les lettres paraissent \emph{rétrécies} $\Rightarrow$ lentille \textbf{divergente}.
\end{itemize}
\item \textbf{Test au soleil} (avec précaution, sans regarder le Soleil) : une lentille convergente concentre la lumière du Soleil en un point lumineux très chaud ; une divergente ne le fait jamais.
\end{enumerate}
\end{methode}

\begin{savaistu}[Le briquet de verre]
En pleine saison sèche, une bouteille en verre abandonnée dans une savane peut jouer le rôle de lentille convergente et concentrer les rayons du Soleil sur l'herbe sèche : c'est une cause connue de départs de feux de brousse. La physique explique parfois les conseils de prudence des campagnes de sensibilisation !
\end{savaistu}

\courssub{Schématisation normalisée et convention de signe}

Pour les schémas et les calculs, on remplace la lentille réelle par un symbole simple.

\begin{definition}[Schématisation d'une lentille mince]
Une lentille mince est représentée par un \emph{segment vertical} traversé en son milieu par l'\cle{axe optique principal} (droite passant par les centres de courbure des faces et par le centre optique) :
\begin{itemize}
\item lentille \textbf{convergente} : segment terminé par des flèches vers l'\emph{extérieur} ;
\item lentille \textbf{divergente} : segment terminé par des flèches vers l'\emph{intérieur}.
\end{itemize}
Le point $O$, intersection de la lentille et de l'axe, est le \cle{centre optique}.
\end{definition}

\begin{aretenir}[Convention de signe]
On choisit comme \textbf{sens positif le sens de propagation de la lumière} (en général de la gauche vers la droite). Toutes les distances sont alors \emph{algébriques}, mesurées à partir du centre optique $O$ : on note $\overline{OA}$, $\overline{OF'}$, etc. Une grandeur mesurée dans le sens de la lumière est positive ; dans le sens contraire, elle est négative. De même, les distances \emph{verticales} (tailles d'objets et d'images) sont comptées positivement vers le haut : $\overline{AB} > 0$ si $B$ est au-dessus de l'axe.
\end{aretenir}

\courssub{Centre optique, foyers et distance focale}

\textbf{Le centre optique.} Comme la lentille est mince, un rayon lumineux qui arrive sur $O$ traverse une zone où les deux faces sont pratiquement parallèles : il se comporte comme à travers une lame à faces parallèles. Or une lame mince décale très légèrement le rayon mais ne le \emph{dévie} pas. On admet donc la propriété fondamentale suivante.

\begin{propriete}[Propriété du centre optique]
Tout rayon lumineux passant par le \cle{centre optique} $O$ d'une lentille mince traverse la lentille \textbf{sans être dévié}. (Propriété admise, justifiée par la faible épaisseur de la lentille : les faces sont localement parallèles au voisinage de $O$.)
\end{propriete}

\textbf{Les foyers.} Envoyons sur une lentille convergente un faisceau de rayons \emph{parallèles à l'axe optique}. À la sortie, les rayons se rencontrent réellement en un point de l'axe : c'est le foyer image. Avec une lentille divergente, au contraire, les rayons s'écartent, mais leurs \emph{prolongements} se coupent en un point de l'axe situé avant la lentille.

\begin{definition}[Foyers principaux]
\begin{itemize}
\item Le \cle{foyer image} $F'$ est le point de l'axe optique où convergent (réellement, ou par leurs prolongements) les rayons incidents \emph{parallèles à l'axe}.
\item Le \cle{foyer objet} $F$ est le point de l'axe optique tel que tout rayon passant par $F$ (ou dirigé vers $F$) émerge de la lentille \emph{parallèle à l'axe}.
\end{itemize}
Pour une lentille mince plongée dans l'air (même milieu de part et d'autre), $F$ et $F'$ sont \textbf{symétriques par rapport à $O$} : $\overline{OF} = -\overline{OF'}$.
\end{definition}

\begin{definition}[Distance focale]
La \cle{distance focale} (ou distance focale image) d'une lentille mince est la grandeur algébrique
\[ f' = \overline{OF'}. \]
Elle s'exprime en mètres (m). Pour une lentille convergente, $F'$ est après la lentille donc $f' > 0$ ; pour une lentille divergente, $F'$ est avant la lentille donc $f' < 0$. On définit de même la distance focale objet $f = \overline{OF}$, avec $f = -f'$ dans l'air.
\end{definition}

\begin{attention}[La distance focale est algébrique !]
C'est le piège numéro 1 du chapitre. $f' = \overline{OF'}$ porte un \textbf{signe} : $f' > 0$ pour une convergente, $f' < 0$ pour une divergente. Si l'énoncé dit « lentille divergente de distance focale 25 cm », tu dois écrire $f' = -0{,}25$ m dans tes calculs, et donc $C = \dfrac{1}{f'} = -4\ \delta$. Oublier le signe fausse tous les résultats.
\end{attention}

\begin{popfigure}
\begin{center}
\begin{tikzpicture}[scale=1.0, >=Stealth]
% Axe optique
\draw[->, popdark] (-5.2,0) -- (5.6,0) node[below left] {\small axe optique};
% Lentille convergente
\draw[popblue, very thick, <->] (0,-2.1) -- (0,2.1);
\node[popblue, above] at (0,2.15) {\small convergente};
% O, F, F'
\fill[popdark] (0,0) circle(1.5pt) node[below left=1pt] {$O$};
\fill[popink] (-2.5,0) circle(1.5pt) node[below] {$F$};
\fill[popink] (2.5,0) circle(1.5pt) node[below] {$F'$};
\draw[<->, poppurple, thick] (0,-0.35) -- (2.5,-0.35) node[midway, below] {$f'=\overline{OF'}$};
% Faisceau parallèle
\foreach \y in {0.7, 1.4}{
  \draw[->, poporange, thick] (-4.5,\y) -- (0,\y);
  \draw[->, poporange, thick] (0,\y) -- (2.5,0);
  \draw[poporange, thick] (2.5,0) -- (3.6,-0.28*\y);
}
\foreach \y in {-0.7, -1.4}{
  \draw[->, poporange, thick] (-4.5,\y) -- (0,\y);
  \draw[->, poporange, thick] (0,\y) -- (2.5,0);
  \draw[poporange, thick] (2.5,0) -- (3.6,-0.28*\y);
}
\node[poporange, above] at (-3.2,1.45) {\small faisceau parallèle à l'axe};
% Plans focaux
\draw[dashed, popmuted] (-2.5,-2.1) -- (-2.5,2.1) node[above, popdark] {\small plan focal objet};
\draw[dashed, popmuted] (2.5,-2.1) -- (2.5,2.1) node[above, popdark] {\small plan focal image};
\end{tikzpicture}
\end{center}
\legende{Lentille convergente : un faisceau parallèle à l'axe optique converge réellement au foyer image $F'$. La distance focale $f' = \overline{OF'}$ est positive.}
\end{popfigure}

\begin{popfigure}
\begin{center}
\begin{tikzpicture}[scale=1.0, >=Stealth]
\draw[->, popdark] (-5.2,0) -- (5.6,0) node[below left] {\small axe optique};
% Lentille divergente (flèches vers l'intérieur)
\draw[poporange, very thick] (0,-2.1) -- (0,2.1);
\draw[poporange, very thick, ->] (0,-2.1) -- (0,-1.6);
\draw[poporange, very thick, ->] (0,2.1) -- (0,1.6);
\node[poporange, above] at (0,2.15) {\small divergente};
\fill[popdark] (0,0) circle(1.5pt) node[below left=1pt] {$O$};
\fill[popink] (-2.5,0) circle(1.5pt) node[below] {$F'$};
\fill[popink] (2.5,0) circle(1.5pt) node[below] {$F$};
\draw[<->, poppurple, thick] (-2.5,-0.35) -- (0,-0.35) node[midway, below] {$f'=\overline{OF'}<0$};
% Rayons
\foreach \y in {0.7, 1.4}{
  \draw[->, popblue, thick] (-4.5,\y) -- (0,\y);
  \draw[->, popblue, thick] (0,\y) -- (3.8,2.2*\y);
  \draw[dashed, popblue] (0,\y) -- (-2.5,0);
}
\foreach \y in {-0.7, -1.4}{
  \draw[->, popblue, thick] (-4.5,\y) -- (0,\y);
  \draw[->, popblue, thick] (0,\y) -- (3.8,2.2*\y);
  \draw[dashed, popblue] (0,\y) -- (-2.5,0);
}
\node[popblue, above] at (-3.2,1.45) {\small faisceau parallèle à l'axe};
\node[popblue, right] at (3.8,1.6) {\small rayons émergents divergents};
\end{tikzpicture}
\end{center}
\legende{Lentille divergente : les rayons parallèles à l'axe émergent en s'écartant ; ce sont leurs \emph{prolongements} (pointillés) qui se coupent au foyer image $F'$, \emph{virtuel}, situé avant la lentille. Ici $f' < 0$.}
\end{popfigure}

\textbf{Plans focaux et foyers secondaires.} Le \cle{plan focal image} est le plan perpendiculaire à l'axe optique passant par $F'$ ; le \cle{plan focal objet} est celui passant par $F$. Un faisceau de rayons parallèles entre eux mais \emph{inclinés} par rapport à l'axe converge, après traversée d'une lentille convergente, en un point du plan focal image appelé \cle{foyer secondaire image} $\varphi'$. Symétriquement, tout rayon passant par un \cle{foyer secondaire objet} $\varphi$ (point du plan focal objet) émerge parallèle à la direction $O\varphi$. Ces propriétés seront précieuses pour construire les images d'objets situés à l'infini (section 3).

\courssub{Vergence d'une lentille mince}

\textbf{Intuition.} Deux loupes ne « concentrent » pas la lumière avec la même force : plus une lentille est bombée, plus elle dévie fortement les rayons et plus son foyer $F'$ est proche de $O$. La grandeur qui mesure cette « capacité à faire converger (ou diverger) la lumière » n'est donc pas $f'$ elle-même, mais son inverse.

\begin{definition}[Vergence]
La \cle{vergence} $C$ d'une lentille mince est l'inverse de sa distance focale image exprimée en mètres :
\[ C = \frac{1}{f'} \qquad \text{avec } f' \text{ en mètres (m)}. \]
L'unité de vergence est la \cle{dioptrie}, de symbole $\delta$ : $1\ \delta = 1\ \text{m}^{-1}$.
\begin{itemize}
\item Lentille convergente : $f' > 0$ donc $C > 0$.
\item Lentille divergente : $f' < 0$ donc $C < 0$.
\end{itemize}
Plus $|C|$ est grande, plus la lentille est « puissante » (elle dévie fortement les rayons).
\end{definition}

\textbf{Analyse dimensionnelle.} $C = \dfrac{1}{f'}$ est l'inverse d'une longueur : $[C] = \text{m}^{-1}$, ce qui est cohérent avec la définition de la dioptrie. Pense à toujours convertir $f'$ en \emph{mètres} avant de calculer $C$ !

\begin{exemplebox}[Calcul d'une distance focale à partir de la vergence]
Une lentille porte l'indication $C = +5\ \delta$. Sa distance focale est :
\[ f' = \frac{1}{C} = \frac{1}{5}\ \text{m} = 0{,}20\ \text{m} = 20\ \text{cm}. \]
Le signe $+$ indique une lentille \emph{convergente} : le foyer image $F'$ est situé 20 cm après la lentille. De même, une lentille de vergence $C = -2\ \delta$ est divergente, de distance focale $f' = \dfrac{1}{-2} = -0{,}50\ \text{m} = -50\ \text{cm}$ : son foyer image est virtuel, situé 50 cm avant la lentille.
\end{exemplebox}

\begin{exoresolu}[Reconnaissance et caractéristiques d'une lentille]
\textbf{Énoncé.} Moussa dispose de deux lentilles $L_1$ et $L_2$. Posée sur le journal, $L_1$ grossit les lettres ; $L_2$ les rétrécit. Le fabricant indique pour $L_1$ une distance focale de valeur absolue $10$ cm et pour $L_2$ une vergence de valeur absolue $4\ \delta$.
\begin{enumerate}
\item Donner la nature (convergente ou divergente) de chaque lentille.
\item Déterminer la distance focale algébrique et la vergence de chaque lentille.
\end{enumerate}
\tcblower
\textbf{Solution détaillée.}
\begin{enumerate}
\item $L_1$ grossit le texte : c'est une lentille \textbf{convergente} (bords minces). $L_2$ rétrécit le texte : c'est une lentille \textbf{divergente} (bords épais).
\item Pour $L_1$, convergente : $f'_1 = +10\ \text{cm} = +0{,}10\ \text{m}$, d'où
\[ C_1 = \frac{1}{f'_1} = \frac{1}{0{,}10} = +10\ \delta. \]
Pour $L_2$, divergente : sa vergence est négative, donc $C_2 = -4\ \delta$, d'où
\[ f'_2 = \frac{1}{C_2} = \frac{1}{-4} = -0{,}25\ \text{m} = -25\ \text{cm}. \]
\end{enumerate}
Vérification : les signes de $f'$ et de $C$ sont cohérents avec la nature de chaque lentille, et les unités (m pour $f'$, $\delta$ pour $C$) sont correctes.
\end{exoresolu}

\begin{attention}[Erreurs fréquentes à éviter]
\begin{itemize}
\item Calculer $C = \dfrac{1}{f'}$ avec $f'$ en \emph{centimètres} : on obtient un résultat 100 fois trop grand. Convertis toujours en mètres.
\item Confondre $F$ et $F'$ : le foyer \emph{image} $F'$ est du côté où émerge la lumière pour une convergente.
\item Croire qu'une lentille divergente « n'a pas de foyer » : elle en a un, mais il est \emph{virtuel} (intersection des prolongements des rayons).
\item Oublier que la dioptrie s'écrit avec le symbole $\delta$ (et non « d » ou « D » en physique, même si les opticiens écrivent souvent « D »).
\end{itemize}
\end{attention}

\begin{savaistu}[La dioptrie chez l'opticien]
L'unité « dioptrie » est utilisée quotidiennement par les opticiens de Yaoundé et de Douala pour prescrire les verres correcteurs. Une ordonnance mentionnant « $-2{,}5\ \delta$ » corrige une \emph{myopie} avec une lentille divergente ; « $+1{,}5\ \delta$ » corrige une \emph{hypermétropie} ou la presbytie (celle de la grand-mère de Moussa !) avec une lentille convergente. La physique de cette leçon est donc littéralement sous tes yeux chaque fois que quelqu'un porte des lunettes.
\end{savaistu}

\begin{aretenir}[L'essentiel de la section]
\begin{itemize}
\item Une lentille est un milieu transparent limité par deux dioptres dont l'un au moins est sphérique ; elle est \emph{mince} si $e \ll R_1, R_2$.
\item Bords minces $\Rightarrow$ convergente ; bords épais $\Rightarrow$ divergente. Une convergente grossit un texte proche, une divergente le rétrécit.
\item Tout rayon passant par le centre optique $O$ n'est pas dévié.
\item Foyer image $F'$ : point de convergence des rayons parallèles à l'axe ; distance focale $f' = \overline{OF'}$ (algébrique), avec $\overline{OF} = -\overline{OF'}$ dans l'air.
\item Vergence : $C = \dfrac{1}{f'}$ ($f'$ en mètres), en dioptries ($\delta$) ; $C > 0$ pour une convergente, $C < 0$ pour une divergente.
\item Sens positif = sens de propagation de la lumière ; toutes les distances sont algébriques.
\end{itemize}
\end{aretenir}

Maintenant que nous savons définir, reconnaître et caractériser une lentille mince, il reste une question : dans quelles conditions ces belles propriétés donnent-elles des images \emph{nettes} ? C'est l'objet de la section suivante, consacrée aux conditions de Gauss et aux propriétés optiques des points remarquables.

\coursec{Conditions de Gauss et propriétés optiques}

Dans la section précédente, nous avons défini les lentilles minces, classé les lentilles convergentes et divergentes, et introduit leurs éléments caractéristiques : centre optique $O$, foyers $F$ et $F'$, plans focaux, distance focale $f'$ et vergence $C$. Mais une question fondamentale reste en suspens : une lentille réelle donne-t-elle toujours une image nette ? L'expérience montre que non : si l'on éclaire largement une lentille avec des rayons très inclinés, l'image d'un point n'est plus un point mais une tache floue. Le physicien allemand Carl Friedrich Gauss a identifié les conditions dans lesquelles une lentille se comporte de manière « idéale ». C'est l'objet de cette section, qui te fournira ensuite les trois rayons particuliers, outils indispensables de toutes les constructions d'images.

\courssub{Les conditions de Gauss}

\textbf{Observation préalable : pourquoi l'image est-elle parfois floue ?}\par

\begin{experience}[Image nette ou image floue ?]
\textbf{Matériel :} banc d'optique, lanterne avec objet (lettre F découpée), lentille convergente de vergence $C = 5\ \delta$, écran, diaphragme à iris.

\textbf{Protocole :}
\begin{enumerate}
\item Former l'image de l'objet sur l'écran avec la lentille largement éclairée (sans diaphragme) et l'objet placé loin de l'axe optique.
\item Observer la netteté des bords de l'image.
\item Diaphragmer la lentille (réduire l'ouverture à quelques millimètres) et recentrer l'objet près de l'axe. Observer à nouveau.
\end{enumerate}

\textbf{Observations :} dans le premier cas, les contours de l'image sont flous, irisés, déformés ; dans le second cas, l'image est beaucoup plus nette, bien que moins lumineuse.

\textbf{Interprétation :} les rayons très inclinés sur l'axe ou passant loin du centre de la lentille ne convergent pas exactement au même point que les rayons centraux : ce sont les \cle{aberrations géométriques}. En limitant le faisceau aux rayons proches de l'axe, on supprime presque totalement ces défauts.
\end{experience}

\textbf{Énoncé des conditions de Gauss}\par

\begin{definition}[Conditions de Gauss — rayons paraxiaux]
Une lentille mince est utilisée dans les \cle{conditions de Gauss} (ou approximation \cle{paraxiale}) lorsque :
\begin{itemize}
\item les rayons lumineux sont \textbf{peu inclinés} sur l'axe optique (angles $i$ de quelques degrés au plus, de sorte que $\sin i \approx \tan i \approx i$ en radians) ;
\item les rayons passent \textbf{près du centre optique} de la lentille (faible distance à l'axe) ;
\item l'objet est \textbf{petit} et placé \textbf{au voisinage de l'axe optique}.
\end{itemize}
Les rayons vérifiant ces conditions sont appelés \cle{rayons paraxiaux}.
\end{definition}

\begin{methode}[Réalisation pratique des conditions de Gauss au banc d'optique]
\begin{enumerate}
\item \textbf{Diaphragmer} la lentille : placer contre elle un diaphragme à iris ou un cache percé d'un trou de quelques millimètres, centré sur le centre optique $O$.
\item \textbf{Centrer} l'objet, la lentille et l'écran : leurs centres doivent être à la même hauteur, sur l'axe optique du banc.
\item Utiliser un \textbf{objet de petite dimension} (quelques centimètres) placé près de l'axe.
\item Employer un \textbf{éclairage modéré} : un faisceau trop large et trop intense favorise les rayons marginaux.
\end{enumerate}
\end{methode}

\begin{exemplebox}[En TP de banc d'optique]
Au lycée, lors de la détermination de la vergence d'une lentille, le professeur exige que tu diaphragmes la lentille et que tu alignes soigneusement l'objet lumineux (lettre F), le centre de la lentille et le centre de l'écran sur l'axe du banc gradué. Ce n'est pas un caprice : sans ces précautions, la position de l'image nette est mal définie (la netteté s'étale sur plusieurs centimètres) et la mesure de $\overline{OA'}$ devient imprécise, ce qui fausse la vérification de la formule de conjugaison.
\end{exemplebox}

\textbf{Conséquences : stigmatisme et aplanétisme approchés}\par

\begin{definition}[Stigmatisme et aplanétisme]
Un système optique est dit \cle{stigmatique} pour un couple de points $(A, A')$ si tout rayon issu de $A$ traversant le système passe par $A'$ après traversée : à un \textbf{point objet} correspond un \textbf{point image}. Il est dit \cle{aplanétique} si l'image d'un objet plan perpendiculaire à l'axe optique est elle-même plane et perpendiculaire à l'axe.
\end{definition}

\begin{propriete}[Comportement de la lentille dans les conditions de Gauss — admise]
Dans les conditions de Gauss, une lentille mince est \textbf{approximativement stigmatique et aplanétique} :
\begin{itemize}
\item tout point objet $A$ donne un point image $A'$ (image nette) ;
\item l'image d'un objet $AB$ perpendiculaire à l'axe optique est un segment $A'B'$ perpendiculaire à l'axe optique.
\end{itemize}
Cette propriété, admise en Première, justifie toutes les constructions géométriques et les formules de conjugaison qui suivront.
\end{propriete}

\begin{attention}[Hors conditions de Gauss : image floue]
Si les conditions de Gauss ne sont pas respectées, la lentille n'est plus stigmatique : l'image d'un point est une tache (aberrations sphériques, coma, astigmatisme) et l'image d'un objet plan peut être déformée. \textbf{Diaphragmer n'est pas facultatif en TP} : c'est la condition de validité de toutes les mesures. De même, les formules des lentilles (position, grandissement) ne s'appliquent que dans l'approximation de Gauss.
\end{attention}

\courssub{Propriétés optiques des points remarquables}

Dans les conditions de Gauss, trois rayons particuliers ont des comportements simples et prévisibles. Ce sont eux qui rendent les constructions d'images possibles.

\textbf{Propriété du centre optique O}\par

\begin{propriete}[Rayon passant par le centre optique]
Tout rayon lumineux passant par le \cle{centre optique} $O$ d'une lentille mince traverse celle-ci \textbf{sans être dévié}.
\end{propriete}

\textbf{Justification intuitive :} au voisinage de $O$, les deux faces d'une lentille mince sont quasiment parallèles ; la lentille s'y comporte comme une lame à faces parallèles d'épaisseur négligeable. Or une lame à faces parallèles déplace latéralement le rayon sans le dévier, et comme la lentille est \emph{mince}, ce déplacement est négligeable : le rayon ressort dans le prolongement du rayon incident.

\textbf{Propriété du foyer image F'}\par

\begin{propriete}[Rayon parallèle à l'axe optique]
Tout rayon incident \textbf{parallèle à l'axe optique} émerge de la lentille en passant par le \cle{foyer image} $F'$ (pour une lentille convergente ; pour une divergente, son prolongement passe par $F'$).
\end{propriete}

C'est d'ailleurs la définition expérimentale de $F'$ vue à la section précédente : le foyer image est le point de convergence d'un faisceau parallèle à l'axe.

\textbf{Propriété du foyer objet F}\par

\begin{propriete}[Rayon passant par le foyer objet]
Tout rayon incident passant par le \cle{foyer objet} $F$ émerge de la lentille \textbf{parallèlement à l'axe optique} (pour une lentille divergente, c'est le rayon dirigé vers $F$ qui émerge parallèle à l'axe).
\end{propriete}

\textbf{Justification :} c'est une conséquence directe du \emph{principe de retour inverse de la lumière} : en renversant le sens de propagation, un rayon parallèle à l'axe qui convergeait en $F'$ devient un rayon issu de $F$ qui émerge parallèle à l'axe (par symétrie, $\overline{OF} = -\overline{OF'}$).

\textbf{Propriété des plans focaux : foyers secondaires}\par

\begin{definition}[Foyer secondaire]
Un \cle{foyer secondaire image} $F'_1$ est le point d'intersection du \cle{plan focal image} (plan perpendiculaire à l'axe passant par $F'$) avec la direction d'un faisceau parallèle incliné. De même, un foyer secondaire objet $F_1$ appartient au plan focal objet.
\end{definition}

\begin{propriete}[Faisceau parallèle incliné]
Un faisceau de rayons parallèles entre eux, incliné d'un angle $\theta$ sur l'axe optique, converge après traversée d'une lentille convergente en un foyer secondaire image $F'_1$ situé dans le plan focal image. Pour le construire, on utilise le rayon du faisceau passant par $O$, qui n'est pas dévié : $F'_1$ est l'intersection de ce rayon avec le plan focal image.
\end{propriete}

\begin{popfigure}
\begin{center}
\begin{tikzpicture}[scale=0.9,>=Stealth]
% Axe optique
\draw[->,thick] (-7,0) -- (7,0) node[below left]{$\Delta$};
% Lentille convergente
\draw[popblue,very thick] (0,-2.8) -- (0,2.8);
\draw[popblue,->,very thick] (0,2.8) -- (-0.28,2.52);
\draw[popblue,->,very thick] (0,-2.8) -- (0.28,-2.52);
% Points caractéristiques
\fill (0,0) circle (1.5pt) node[below right]{$O$};
\fill (-3,0) circle (1.5pt) node[below]{$F$};
\fill (3,0) circle (1.5pt) node[below]{$F'$};
\draw[dashed,popmuted] (-3,-0.3) -- (-3,0.3);
\draw[dashed,popmuted] (3,-0.3) -- (3,0.3);
% Objet AB (u = -6 cm, f' = 3 cm => v = 6 cm, gamma = -1)
\draw[very thick,poporange,->] (-6,0) -- (-6,1.6) node[above]{$B$};
\fill (-6,0) circle (1.5pt) node[below]{$A$};
% Rayon 1 : parallèle à l'axe -> F'
\draw[popink,thick,->] (-6,1.6) -- (0,1.6);
\draw[popink,thick] (0,1.6) -- (3,0) -- (6,-1.6);
% Rayon 2 : centre optique O (non dévié)
\draw[popgreen,thick,->] (-6,1.6) -- (0,0);
\draw[popgreen,thick] (0,0) -- (6,-1.6);
% Rayon 3 : foyer objet F -> parallèle à l'axe
\draw[poppurple,thick,->] (-6,1.6) -- (0,-1.6);
\draw[poppurple,thick] (0,-1.6) -- (6,-1.6);
% Image A'B'
\fill (6,0) circle (1.5pt) node[below]{$A'$};
\draw[very thick,poporange,->] (6,0) -- (6,-1.6) node[below]{$B'$};
\draw[dashed,popmuted] (6,-1.9) -- (6,1.9);
% Légendes des rayons
\node[popink] at (-3,1.95){\small 1};
\node[popgreen] at (-3,0.95){\small 2};
\node[poppurple] at (-3,-1.95){\small 3};
\end{tikzpicture}
\end{center}
\legende{Les trois rayons particuliers pour une lentille convergente ($f'=3$ cm, objet en $u=-6$ cm) : (1) rayon issu de $B$ parallèle à l'axe, émergeant par $F'$ ; (2) rayon passant par $O$, non dévié ; (3) rayon passant par $F$, émergeant parallèle à l'axe. Leur intersection donne $B'$, image de $B$ (réelle, inversée, même taille).}
\end{popfigure}

\begin{methode}[Construction d'une image : deux rayons suffisent, le troisième vérifie]
Pour construire l'image $A'B'$ d'un objet $AB$ perpendiculaire à l'axe ($A$ sur l'axe) :
\begin{enumerate}
\item Placer l'objet $AB$, la lentille, $O$, $F$ et $F'$ sur un schéma à l'échelle.
\item Tracer le rayon issu de $B$ \textbf{parallèle à l'axe} : il émerge en passant par $F'$.
\item Tracer le rayon issu de $B$ \textbf{passant par $O$} : il n'est pas dévié.
\item L'intersection de ces deux rayons émergents est $B'$ ; abaisser la perpendiculaire à l'axe pour obtenir $A'$.
\item \textbf{Vérification :} le rayon issu de $B$ passant par $F$ doit émerger parallèle à l'axe et passer exactement par $B'$. Si ce n'est pas le cas, la construction contient une erreur.
\end{enumerate}
Deux rayons suffisent donc ; le troisième sert de contrôle. Pour une lentille divergente, ou si l'objet est entre $F$ et la lentille (image virtuelle), ce sont les \emph{prolongements} des rayons émergents (tracés en pointillés) qui se coupent.
\end{methode}

\begin{popfigure}
\begin{center}
\begin{tikzpicture}[scale=0.85,>=Stealth]
% Axe optique
\draw[->,thick] (-6,0) -- (5,0) node[below left]{$\Delta$};
% Lentille
\draw[popblue,very thick] (0,-3) -- (0,3);
\draw[popblue,->,very thick] (0,3) -- (-0.28,2.72);
\draw[popblue,->,very thick] (0,-3) -- (0.28,-2.72);
% Plan focal image
\draw[dashed,popgold,thick] (3,-3) -- (3,3) node[above]{Plan focal image};
% Points
\fill (0,0) circle (1.5pt) node[below right]{$O$};
\fill (3,0) circle (1.5pt) node[below]{$F'$};
% Faisceau parallèle incliné (angle theta ~ -20 deg, tan = -0.364)
% Rayon 1 (rose) : entre à y=2.6 en x=-5
\draw[popink,thick,->] (-5,2.6) -- (0,0.78);
\draw[popink,thick] (0,0.78) -- (3,-0.31);
% Rayon 2 (vert) : entre à y=1.6 en x=-5
\draw[popgreen,thick,->] (-5,1.6) -- (0,-0.22);
\draw[popgreen,thick] (0,-0.22) -- (3,-1.31);
% Rayon 3 (violet) : passe par O (non dévié), même pente
\draw[poppurple,thick,->] (-5,1.82) -- (0,0);
\draw[poppurple,thick] (0,0) -- (3,-1.09);
% Foyer secondaire F'1 (intersection rayon violet / plan focal)
\fill[poporange] (3,-1.09) circle (2pt) node[right,popdark]{$F'_1$};
% Angle theta
\draw[popmuted] (-3.2,0) arc (180:200:1.8);
\node[popmuted] at (-3.6,-0.45){$\theta$};
% Cotes
\draw[|<->|,popmuted] (0,-2.5) -- (3,-2.5) node[midway,below]{$f'$};
\end{tikzpicture}
\end{center}
\legende{Faisceau parallèle incliné d'un angle $\theta$ sur l'axe : tous les rayons émergents convergent au foyer secondaire $F'_1$ du plan focal image. Le rayon passant par $O$ (violet), non dévié, permet de localiser $F'_1$ : $F'F'_1 = f'\tan\theta$.}
\end{popfigure}

\begin{exoresolu}[Position d'un foyer secondaire]
Une lentille convergente de distance focale $f' = \overline{OF'} = \SI{10}{cm}$ reçoit un faisceau de rayons parallèles incliné d'un angle $\theta = 5^{\circ}$ sur l'axe optique, dans les conditions de Gauss. Déterminer la position du point de convergence $F'_1$ du faisceau.
\tcblower
\textbf{Solution.} D'après la propriété des plans focaux, le faisceau converge au foyer secondaire $F'_1$, intersection du plan focal image avec le rayon passant par $O$ (non dévié). Dans le triangle rectangle formé par $O$, $F'$ et $F'_1$ :
\[ \tan\theta = \dfrac{F'F'_1}{OF'} = \dfrac{F'F'_1}{f'} \quad\Longrightarrow\quad F'F'_1 = f'\tan\theta. \]
Application numérique : $F'F'_1 = 10 \times \tan 5^{\circ} = 10 \times 0,0875 \approx \SI{0,87}{cm}$.

Le faisceau converge donc en $F'_1$, situé dans le plan focal image, à $\SI{10}{cm}$ de $O$ et à environ $\SI{0,9}{cm}$ de $F'$ (du côté opposé à l'inclinaison du faisceau). Vérification de cohérence : $\theta$ étant petit, $\tan\theta \approx \theta \approx 0,087$ rad, ce qui confirme que l'on reste dans l'approximation paraxiale.
\end{exoresolu}

\begin{attention}[Pièges fréquents sur les rayons particuliers]
\begin{itemize}
\item Ne pas confondre $F$ et $F'$ : le foyer \textbf{objet} $F$ est du côté de la lumière incidente, le foyer \textbf{image} $F'$ de l'autre côté. Un rayon parallèle à l'axe émerge par $F'$, jamais par $F$.
\item Le rayon passant par $O$ n'est pas dévié, mais il n'est pas « réfracté deux fois sans effet » par magie : c'est l'épaisseur nulle de la lentille mince qui justifie l'absence de décalage.
\item Pour une lentille \textbf{divergente}, les rayons émergents divergent : ce sont leurs \textbf{prolongements} (en pointillés) qui passent par $F'$ ou se coupent pour former l'image virtuelle.
\item Les constructions ne sont valables que dans les conditions de Gauss : garde les angles petits sur tes schémas conceptuels, même si l'échelle du dessin exagère les inclinaisons pour la lisibilité.
\end{itemize}
\end{attention}

\begin{savaistu}[Gauss, les lunettes et les télescopes]
Carl Friedrich Gauss (1777--1855) formula l'approximation paraxiale en étudiant les instruments d'optique de son époque. C'est grâce à elle que l'on calcule aujourd'hui les lunettes correctrices : un œil myope est corrigé par une lentille divergente dont l'opticien ajuste la vergence au quart de dioptrie près. Les grands télescopes, eux, diaphragment naturellement la lumière des étoiles : les rayons arrivant quasi parallèles et proches de l'axe, ils travaillent presque parfaitement dans les conditions de Gauss, ce qui explique la netteté remarquable des images stellaires.
\end{savaistu}

\begin{aretenir}[L'essentiel de la section]
\begin{itemize}
\item \textbf{Conditions de Gauss :} rayons paraxiaux (peu inclinés, proches de l'axe, passant près de $O$), objet petit et centré ; en pratique, on diaphragme la lentille.
\item \textbf{Conséquence :} la lentille est approximativement \textbf{stigmatique} (un point objet donne un point image) et \textbf{aplanétique} (un objet perpendiculaire à l'axe donne une image perpendiculaire à l'axe).
\item \textbf{Trois rayons particuliers :} rayon par $O$ non dévié ; rayon parallèle à l'axe émergeant par $F'$ ; rayon passant par $F$ émergeant parallèle à l'axe. Deux suffisent pour construire l'image, le troisième sert de vérification.
\item \textbf{Plans focaux :} un faisceau parallèle incliné converge en un foyer secondaire $F'_1$ du plan focal image, localisé grâce au rayon non dévié passant par $O$ ($F'F'_1 = f'\tan\theta$).
\end{itemize}
\end{aretenir}

\coursec{Construction géométrique des images}

Dans la section précédente, nous avons établi les \cle{conditions de Gauss} et découvert les propriétés remarquables du centre optique $O$ et des foyers $F$ et $F'$ : un rayon passant par $O$ n'est pas dévié, un rayon parallèle à l'axe émerge en passant par $F'$, un rayon passant par $F$ émerge parallèle à l'axe. Ces trois résultats, simples en apparence, sont en réalité des outils extraordinairement puissants : ils vont nous permettre de \emph{prévoir}, par une simple construction à la règle, la position, la nature, le sens et la taille de l'image qu'une lentille donne d'un objet. C'est exactement ce que fait l'opticien qui règle un projecteur, ou l'ingénieur qui conçoit l'objectif d'un appareil photo. Apprenons ce geste, étape par étape.

\courssub{Principe général de la construction}

\textbf{Objet et image : les conventions}\par

Dans toute cette leçon, l'objet est un petit segment $\overrightarrow{AB}$, \cle{perpendiculaire à l'axe optique}, le point $A$ étant \emph{sur} l'axe. Cette situation est en réalité très générale : un objet étendu (une flamme de bougie, un visage) peut toujours être décomposé en points, et le point $B$ le plus éloigné de l'axe est celui qui fixe la taille de l'image.

L'idée directrice est la suivante : tous les rayons issus de $B$ qui traversent la lentille dans les conditions de Gauss convergent (réellement ou par leurs prolongements) en un \emph{unique} point $B'$, image de $B$. Pour trouver $B'$, il suffit donc de tracer \textbf{deux} rayons particuliers dont on connaît le comportement ; leur intersection (ou celle de leurs prolongements) donne $B'$. Comme l'objet est perpendiculaire à l'axe, l'image l'est aussi : le point $A'$, image de $A$, est simplement le \cle{pied de la perpendiculaire} abaissée de $B'$ sur l'axe.

\begin{methode}[Construire une image en quatre étapes]
\begin{enumerate}
\item \textbf{Placer le squelette} : l'axe optique, la lentille (schématisée par un segment à flèches), son centre optique $O$, et les foyers $F$ et $F'$ symétriques par rapport à $O$, avec $\overline{OF'} = f'$.
\item \textbf{Premier rayon} : depuis $B$, tracer le rayon \textbf{parallèle à l'axe} ; il émerge de la lentille en passant par le foyer image $F'$.
\item \textbf{Deuxième rayon} : depuis $B$, tracer le rayon passant par le \textbf{centre optique $O$} ; il traverse la lentille \emph{sans déviation}.
\item \textbf{Conclure} : l'intersection des deux rayons émergents donne $B'$. Si les rayons se coupent réellement, l'image est \textbf{réelle} ; s'ils divergent, prolonger en pointillés du côté objet : l'image est \textbf{virtuelle}. Abaisser la perpendiculaire de $B'$ sur l'axe pour obtenir $A'$.
\end{enumerate}
On peut vérifier la construction avec un troisième rayon : celui qui passe par le foyer objet $F$ émerge parallèle à l'axe et doit passer par $B'$.
\end{methode}

\begin{propriete}[Image réelle et image virtuelle]
Une image est dite \cle{réelle} si les rayons lumineux émergents \emph{convergent réellement} en ses points : elle peut alors être recueillie sur un écran (mur, papier, capteur). Elle est dite \cle{virtuelle} si ce sont les \emph{prolongements} des rayons émergents qui se coupent : aucun écran ne peut la recevoir, mais on la \emph{voit} en regardant à travers la lentille, car l'œil prolonge mentalement les rayons qu'il reçoit.
\end{propriete}

\begin{attention}[Réelle ou virtuelle : ne pas se tromper de critère]
La question à te poser est toujours la même : \emph{la lumière passe-t-elle réellement par le point image ?} Une image réelle est toujours située \emph{après} la lentille (du côté où sort la lumière) pour un objet réel ; une image virtuelle est située \emph{avant} la lentille (du côté de l'objet). Sur tes schémas, les prolongements virtuels se tracent \textbf{en pointillés} et ne doivent jamais être confondus avec de vrais rayons.
\end{attention}

\courssub{Les différents cas pour une lentille convergente}

\textbf{Cas 1 : objet réel placé avant le foyer ($\overline{OA} < -f'$)}\par

C'est le cas le plus fréquent : photographier un ami, projeter un diaporama. L'objet est placé devant la lentille, au-delà du foyer objet. La construction montre que les rayons émergents \emph{convergent réellement} de l'autre côté de la lentille.

\begin{propriete}[Objet réel avant F, lentille convergente]
Si l'objet réel est placé avant le foyer objet ($\overline{OA} < -f'$), l'image est \textbf{réelle}, \textbf{renversée} et située après la lentille. Sa taille dépend de la position de l'objet :
\begin{itemize}
\item si l'objet est \emph{au-delà de $2f'$} ($\overline{OA} < -2f'$) : image \textbf{réduite}, entre $F'$ et $2f'$ ;
\item si l'objet est \emph{entre $F$ et $2f'$} : image \textbf{agrandie}, au-delà de $2f'$.
\end{itemize}
\end{propriete}

\begin{popfigure}
\begin{center}
\begin{tikzpicture}[scale=0.85,>=Stealth]
\draw[popmuted,thick,->] (-7,0) -- (7.5,0) node[below left] {axe optique};
\draw[popblue,very thick] (0,-2.6) -- (0,2.6);
\draw[popblue,->] (0,2.6) -- (0,2.2);
\draw[popblue,->] (0,-2.6) -- (0,-2.2);
\fill (0,0) circle (1.5pt) node[below right] {$O$};
\fill (-2,0) circle (1.5pt) node[below] {$F$};
\fill (2,0) circle (1.5pt) node[below] {$F'$};
\fill (-4,0) circle (1pt) node[below] {$2F$};
\fill (4,0) circle (1pt) node[below] {$2F'$};
\draw[poporange,very thick,->] (-5,0) node[below] {$A$} -- (-5,1.5) node[above] {$B$};
\draw[popgreen,thick] (-5,1.5) -- (0,1.5);
\draw[popgreen,thick,->] (0,1.5) -- (2,0);
\draw[popgreen,thick] (2,0) -- (3.34,-1);
\draw[popgold,thick] (-5,1.5) -- (0,0);
\draw[popgold,thick,->] (0,0) -- (3.34,-1);
\draw[poppink,thick,dashed] (-5,1.5) -- (-2,0) -- (0,-0.7);
\draw[poppink,thick,dashed] (0,-0.7) -- (3.34,-1);
\draw[popink,very thick,->] (3.34,0) node[above] {$A'$} -- (3.34,-1) node[below] {$B'$};
\draw[popmuted,dashed] (3.34,-1) -- (3.34,0);
\end{tikzpicture}
\end{center}
\legende{Objet réel au-delà de $2f'$ : l'image $A'B'$ est réelle, renversée et réduite, située entre $F'$ et $2F'$.}
\end{popfigure}

\textbf{Position particulière : l'objet à $2f'$}\par

Il existe une position remarquable, très appréciée au Probatoire : lorsque l'objet est placé à une distance $2f'$ devant la lentille.

\begin{propriete}[Position des points de Weierstrass]
Si $\overline{OA} = -2f'$, alors l'image est réelle, renversée, située à $\overline{OA'} = +2f'$, et \textbf{de même taille} que l'objet ($\gamma = -1$). Les points $A$ et $A'$, symétriques par rapport à $O$, sont appelés \emph{points de Weierstrass} (ou points antiprincipaux).
\end{propriete}

\textbf{Cas 2 : objet entre F et O — le fonctionnement de la loupe}\par

Rapprochons maintenant l'objet : plaçons-le \emph{entre} le foyer objet $F$ et le centre optique $O$. Le rayon parallèle à l'axe émerge toujours par $F'$, et le rayon par $O$ n'est pas dévié ; mais cette fois, ces deux rayons émergents \emph{divergent} : ils ne se croiseront jamais de l'autre côté. En revanche, leurs \emph{prolongements} (en pointillés, du côté de l'objet) se coupent en un point $B'$.

\begin{propriete}[Objet entre F et O : la loupe]
Si l'objet réel est placé entre le foyer objet et la lentille ($-f' < \overline{OA} < 0$), l'image est \textbf{virtuelle}, \textbf{droite} et \textbf{agrandie}. C'est le principe de la \cle{loupe} : on observe l'image en regardant à travers la lentille.
\end{propriete}

\begin{popfigure}
\begin{center}
\begin{tikzpicture}[scale=0.85,>=Stealth]
\draw[popmuted,thick,->] (-6,0) -- (6.5,0) node[below left] {axe optique};
\draw[popblue,very thick] (0,-2.8) -- (0,2.8);
\draw[popblue,->] (0,2.8) -- (0,2.4);
\draw[popblue,->] (0,-2.8) -- (0,-2.4);
\fill (0,0) circle (1.5pt) node[below right] {$O$};
\fill (-2,0) circle (1.5pt) node[below] {$F$};
\fill (2,0) circle (1.5pt) node[below] {$F'$};
\draw[poporange,very thick,->] (-1,0) node[below] {$A$} -- (-1,0.8) node[above] {$B$};
\draw[popgreen,thick] (-1,0.8) -- (0,0.8);
\draw[popgreen,thick,->] (0,0.8) -- (4.5,-1);
\draw[popgreen,thick,dashed] (0,0.8) -- (-2,1.6);
\draw[popgold,thick,->] (-1,0.8) -- (4.5,-3.52);
\draw[popgold,thick,dashed] (0,0) -- (-2,1.6);
\draw[popink,very thick,->] (-2,0) node[below] {$A'$} -- (-2,1.6) node[above] {$B'$};
\draw[popmuted,dashed] (-2,1.6) -- (-2,0);
\end{tikzpicture}
\end{center}
\legende{Objet entre $F$ et $O$ : les rayons émergents divergent ; leurs prolongements (pointillés) se coupent en $B'$. L'image est virtuelle, droite et agrandie : c'est la loupe.}
\end{popfigure}

\begin{savaistu}[La loupe, outil de précision]
Le bijoutier qui examine une pierre précieuse, la couturière qui pique un ourlet fin, le technicien de CAMTEL qui inspecte une soudure sur fibre optique : tous utilisent le même principe. En plaçant l'objet \emph{entre le foyer et la lentille}, ils obtiennent une image virtuelle droite et agrandie, confortable à observer. Plus la distance focale de la loupe est courte, plus le grossissement est important : c'est pourquoi les loupes d'horloger sont de petites lentilles très bombées, donc très convergentes.
\end{savaistu}

\textbf{Cas 3 : objet à l'infini et objet au foyer}\par

Deux situations limites achèvent le panorama :

\begin{propriete}[Positions limites]
\begin{itemize}
\item \textbf{Objet à l'infini} : les rayons arrivent parallèles entre eux ; l'image se forme dans le \cle{plan focal image} (en $F'$ si l'objet est sur l'axe). C'est le cas d'une étoile ou du Soleil observés avec une lunette.
\item \textbf{Objet au foyer objet} ($A$ confondu avec $F$) : les rayons émergent \emph{parallèles} entre eux ; l'image est \textbf{rejetée à l'infini}. C'est le principe du projecteur de diapositives et du phare de voiture : la source placée au foyer fournit un faisceau quasi parallèle.
\end{itemize}
\end{propriete}

\courssub{Cas de la lentille divergente}

Pour une lentille divergente, les foyers sont \emph{virtuels} : $F'$ est du côté de l'objet ($f' < 0$). Le rayon parallèle à l'axe émerge en \emph{s'écartant} de l'axe, comme s'il provenait de $F'$ ; le rayon dirigé vers $O$ n'est pas dévié. Les rayons émergents divergent toujours pour un objet réel.

\begin{propriete}[Image donnée par une lentille divergente]
Pour tout objet réel, une lentille divergente donne une image \textbf{toujours virtuelle}, \textbf{droite} et \textbf{réduite}, située entre $O$ et $F'$. C'est pourquoi les lentilles correctrices de la myopie (divergentes) donnent des images plus petites.
\end{propriete}

\begin{attention}[Divergente : prolonger du bon côté]
Avec une lentille divergente, les rayons émergents \emph{s'écartent} : il faut impérativement les \textbf{prolonger en pointillés du côté de l'objet} pour trouver $B'$. Erreur classique : chercher l'intersection des rayons réels après la lentille — elle n'existe pas ! Autre piège : le foyer image $F'$ d'une divergente est \emph{avant} la lentille ; vérifie toujours le signe de $f'$ avant de placer tes foyers.
\end{attention}

\begin{popfigure}
\begin{center}
\begin{tikzpicture}[scale=0.85,>=Stealth]
\draw[popmuted,thick,->] (-6,0) -- (6,0) node[below left] {axe optique};
\draw[popblue,very thick] (0,-2.4) -- (0,2.4);
\draw[popblue,->] (0,2.4) -- (0.35,2.75);
\draw[popblue,->] (0,2.4) -- (-0.35,2.75);
\draw[popblue,->] (0,-2.4) -- (0.35,-2.75);
\draw[popblue,->] (0,-2.4) -- (-0.35,-2.75);
\fill (0,0) circle (1.5pt) node[below right] {$O$};
\fill (-2,0) circle (1.5pt) node[below] {$F'$};
\fill (2,0) circle (1.5pt) node[below] {$F$};
\draw[poporange,very thick,->] (-3.5,0) node[below] {$A$} -- (-3.5,1.4) node[above] {$B$};
\draw[popgreen,thick] (-3.5,1.4) -- (0,1.4);
\draw[popgreen,thick,->] (0,1.4) -- (4.5,3.15);
\draw[popgreen,thick,dashed] (0,1.4) -- (-2,0);
\draw[popgold,thick,->] (-3.5,1.4) -- (4.5,-1.83);
\draw[popgold,thick,dashed] (0,0) -- (-1.2,0.48);
\draw[popink,very thick,->] (-1.2,0) node[below] {$A'$} -- (-1.2,0.48) node[above right] {$B'$};
\end{tikzpicture}
\end{center}
\legende{Lentille divergente : les rayons émergents divergent ; leurs prolongements (pointillés) se coupent en $B'$. L'image est virtuelle, droite et réduite.}
\end{popfigure}

\courssub{Image réelle ou virtuelle : comment les observer ?}

Récapitulons la différence pratique, car elle fait l'objet de nombreuses questions d'examen. Une \cle{image réelle} peut être \emph{projetée sur un écran} : place une feuille blanche à l'endroit où convergent les rayons, et tu y verras l'image se former (comme au cinéma ou avec le projecteur de la salle de classe). Une \cle{image virtuelle} ne peut pas être recueillie sur un écran : pour la voir, il faut placer l'œil \emph{derrière la lentille} et regarder à travers elle — l'œil reçoit les rayons divergents et les prolonge mentalement jusqu'à l'image (comme devant un miroir plan ou à travers une loupe).

\courssub{Construction à travers un système de deux lentilles}

De nombreux instruments (microscope, lunette astronomique, objectif d'appareil photo) associent plusieurs lentilles sur le même axe. La méthode de construction repose sur une idée simple et puissante.

\begin{methode}[Image à travers deux lentilles alignées]
\begin{enumerate}
\item Construire l'image $A_1B_1$ de l'objet $AB$ donnée par la \textbf{première lentille} $L_1$, \emph{comme si la seconde n'existait pas}.
\item Considérer $A_1B_1$ comme un \textbf{objet pour la deuxième lentille} $L_2$ (objet réel si $A_1B_1$ est avant $L_2$, virtuel sinon).
\item Construire l'image $A_2B_2$ de $A_1B_1$ à travers $L_2$ avec les rayons particuliers de $L_2$.
\item Conclure : $A_2B_2$ est l'image définitive de $AB$ à travers le système. Le grandissement total est le produit des grandissements : $\gamma = \gamma_1 \times \gamma_2$.
\end{enumerate}
\end{methode}

\begin{popfigure}
\begin{center}
\begin{tikzpicture}[scale=0.8,>=Stealth]
\draw[popmuted,thick,->] (-5.5,0) -- (9,0) node[below left] {axe optique};
\draw[popblue,very thick] (0,-2.4) -- (0,2.4);
\draw[popblue,->] (0,2.4) -- (0,2.05);
\draw[popblue,->] (0,-2.4) -- (0,-2.05);
\node[popblue] at (0,2.8) {$L_1$};
\fill (0,0) circle (1.5pt) node[below right] {$O_1$};
\fill (-1.5,0) circle (1pt) node[below] {$F_1$};
\fill (1.5,0) circle (1pt) node[below] {$F'_1$};
\draw[popblue,very thick] (5,-2.4) -- (5,2.4);
\draw[popblue,->] (5,2.4) -- (5,2.05);
\draw[popblue,->] (5,-2.4) -- (5,-2.05);
\node[popblue] at (5,2.8) {$L_2$};
\fill (5,0) circle (1.5pt) node[below right] {$O_2$};
\fill (4,0) circle (1pt) node[below] {$F_2$};
\fill (6,0) circle (1pt) node[below] {$F'_2$};
\draw[poporange,very thick,->] (-3,0) node[below] {$A$} -- (-3,1) node[above] {$B$};
\draw[popgreen,thick] (-3,1) -- (0,1);
\draw[popgreen,thick] (0,1) -- (3,-1);
\draw[popgold,thick] (-3,1) -- (3,-1);
\draw[popink,very thick,->] (3,0) node[above] {$A_1$} -- (3,-1) node[below] {$B_1$};
\draw[popgreen,thick] (3,-1) -- (5,-1);
\draw[popgreen,thick,->] (5,-1) -- (7,1);
\draw[popgold,thick,->] (3,-1) -- (7,1);
\draw[poppurple,very thick,->] (7,0) node[below] {$A_2$} -- (7,1) node[above] {$B_2$};
\draw[popmuted,dashed] (7,1) -- (7,0);
\end{tikzpicture}
\end{center}
\legende{Système de deux lentilles convergentes : $L_1$ donne de $AB$ l'image intermédiaire $A_1B_1$, qui sert d'objet pour $L_2$ ; l'image finale $A_2B_2$ est réelle et droite par rapport à l'objet initial.}
\end{popfigure}

Remarque sur la figure : l'image intermédiaire $A_1B_1$ est renversée par rapport à $AB$, puis $L_2$ la renverse à nouveau ; l'image finale $A_2B_2$ est donc \emph{droite} par rapport à l'objet initial. C'est exactement ce qui se passe dans une lunette terrestre, qui redresse l'image.

\courssub{Exemple chiffré de construction à l'échelle}

\begin{exemplebox}[Construction à l'échelle 1/5]
On place un objet $AB$ de hauteur $\SI{2}{cm}$ à $\SI{30}{cm}$ devant une lentille convergente de distance focale $f' = \SI{20}{cm}$. Construisons l'image à l'échelle $1/5$ : sur le papier, $\SI{5}{cm}$ réels sont représentés par $\SI{1}{cm}$.
\begin{itemize}
\item Sur le schéma : $OA = \SI{6}{cm}$ (pour $30$ cm réels), $OF' = \SI{4}{cm}$ (pour $20$ cm réels), $AB = \SI{0,4}{cm}$ (pour $2$ cm réels).
\item L'objet est entre $F$ et $2F$ ($20 < 30 < 40$ cm) : on prévoit une image réelle, renversée, agrandie, au-delà de $2f'$.
\item On trace le rayon issu de $B$ parallèle à l'axe, qui émerge par $F'$, et le rayon $BO$ non dévié : ils se coupent en $B'$.
\item On mesure sur le papier $OA' \approx \SI{12}{cm}$, soit $\overline{OA'} = +\SI{60}{cm}$ en réalité, et $A'B' \approx \SI{0,8}{cm}$, soit $\SI{4}{cm}$ réels.
\end{itemize}
L'image est réelle, renversée, deux fois plus grande que l'objet : la construction confirme la prévision. (La formule de conjugaison, étudiée dans la section suivante, donnera exactement $\overline{OA'} = +60$ cm et $\gamma = -2$.)
\end{exemplebox}

\begin{exoresolu}[Construction et caractérisation d'une image]
\textbf{Énoncé.} Un objet $AB$ de hauteur $\SI{3}{cm}$ est placé à $\SI{12}{cm}$ devant une lentille convergente de vergence $C = +10\ \delta$. 1) Calculer la distance focale. 2) Décrire la construction géométrique de l'image et prévoir ses caractéristiques. 3) Où faut-il placer un écran pour observer l'image nette ?
\tcblower
\textbf{Solution.}
1) La distance focale est l'inverse de la vergence :
\[ f' = \frac{1}{C} = \frac{1}{10} = \SI{0,10}{m} = \SI{10}{cm}. \]
2) L'objet est à $\SI{12}{cm}$ de $O$, donc entre $F$ ($\SI{10}{cm}$) et $2F$ ($\SI{20}{cm}$) : on a $-2f' < \overline{OA} < -f'$. Construction : on trace le rayon issu de $B$ parallèle à l'axe, qui émerge en passant par $F'$, puis le rayon $BO$ non dévié ; leur intersection réelle donne $B'$, et le pied de la perpendiculaire donne $A'$. L'image est donc \textbf{réelle}, \textbf{renversée} et \textbf{agrandie}, située au-delà de $2f'$.
3) L'image étant réelle, on peut la recueillir sur un écran placé à l'endroit où convergent les rayons, c'est-à-dire en $A'$, au-delà de $2f' = \SI{20}{cm}$ derrière la lentille (le calcul exact par la formule de conjugaison donnera $\overline{OA'} = +60$ cm).
\end{exoresolu}

\begin{aretenir}[L'essentiel de la construction géométrique]
\begin{itemize}
\item Deux rayons particuliers suffisent : le rayon parallèle à l'axe (qui émerge par $F'$) et le rayon passant par $O$ (non dévié) ; le rayon par $F$ (émergeant parallèle) sert de vérification.
\item Image \textbf{réelle} : les rayons convergent réellement, on la projette sur un écran. Image \textbf{virtuelle} : ce sont les prolongements (pointillés) qui se coupent, on la voit en regardant à travers la lentille.
\item Convergente : objet avant $F$ $\Rightarrow$ image réelle renversée ; objet entre $F$ et $O$ $\Rightarrow$ image virtuelle droite agrandie (loupe) ; objet à $2f'$ $\Rightarrow$ image à $2f'$, même taille, renversée.
\item Divergente : pour un objet réel, image toujours virtuelle, droite et réduite.
\item Deux lentilles : l'image $A_1B_1$ donnée par $L_1$ sert d'objet pour $L_2$ ; le grandissement total est $\gamma = \gamma_1 \gamma_2$.
\end{itemize}
\end{aretenir}

Ces constructions géométriques sont magnifiques, mais elles restent des mesures sur papier, entachées d'erreurs de tracé. Pour \emph{calculer} avec précision la position et la taille d'une image, il nous faut des relations algébriques : c'est l'objet de la section suivante, consacrée aux formules de position et de grandissement — dont tu pourras vérifier qu'elles confirment exactement tout ce que nos constructions viennent de révéler.

\coursec{Formules de position et de grandissement}

Dans la section précédente, tu as appris à construire géométriquement, règle et crayon en main, l'image d'un objet donnée par une lentille mince. Cette méthode graphique est précieuse : elle donne une intuition immédiate de la position, du sens et de la taille de l'image. Mais elle a ses limites : une construction à main levée ne fournit pas de valeurs précises, et elle devient fastidieuse lorsqu'il faut traiter plusieurs cas. Les physiciens ont donc cherché à traduire la géométrie en \cle{formules algébriques} : c'est l'objet de cette section. Nous allons établir rigoureusement les deux relations fondamentales des lentilles minces — la \cle{formule de conjugaison} et la \cle{formule du grandissement} — puis apprendre à les exploiter numériquement et à les vérifier expérimentalement au banc d'optique.

\courssub{Notations algébriques : le choix de l'origine au centre optique}

Pour écrire des formules, il faut d'abord \emph{algébriser} les positions. On munit l'axe optique d'un repère :

\begin{itemize}
\item l'\cle{origine} est le \textbf{centre optique} $O$ de la lentille ;
\item l'\cle{orientation} est celle du sens de propagation de la lumière (de gauche à droite sur nos schémas).
\end{itemize}

On note alors :

\begin{definition}[Positions algébriques objet et image]
La position de l'objet $A$ (pied de l'objet sur l'axe) est repérée par l'abscisse algébrique
\[ p = \overline{OA} \]
et la position de l'image $A'$ par
\[ p' = \overline{OA'}. \]
La distance focale image est $f' = \overline{OF'}$ et la distance focale objet $f = \overline{OF} = -f'$.
\end{definition}

Ces grandeurs sont \textbf{algébriques} : elles ont un signe.

\begin{itemize}
\item Un \textbf{objet réel} est placé \emph{avant} la lentille (du côté d'où vient la lumière) : $p = \overline{OA} < 0$.
\item Une \textbf{image réelle} se forme \emph{après} la lentille : $p' = \overline{OA'} > 0$.
\item Une \textbf{image virtuelle} se forme \emph{avant} la lentille : $p' < 0$.
\item Pour une lentille \textbf{convergente} : $f' > 0$ ; pour une lentille \textbf{divergente} : $f' < 0$.
\end{itemize}

\begin{attention}[Le piège du signe de $\overline{OA}$]
L'erreur la plus fréquente au Probatoire est d'oublier que $\overline{OA}$ est \textbf{négatif} pour un objet réel. Si l'énoncé dit « l'objet est placé à $30\ \mathrm{cm}$ devant la lentille », tu dois écrire $p = \overline{OA} = -30\ \mathrm{cm} = \num{-0,30}\ \mathrm{m}$, et non $+30\ \mathrm{cm}$. Deuxième piège : la vergence s'exprime en dioptries ($\delta$) avec $f'$ en \textbf{mètres} : convertis toujours toutes les distances en mètres avant d'appliquer $C = \dfrac{1}{f'}$.
\end{attention}

\courssub{Formule de conjugaison (formule de position)}

\courssub{Énoncé et intuition}

La construction géométrique montre que la position $p'$ de l'image dépend de la position $p$ de l'objet et de la distance focale $f'$ de la lentille. La formule de conjugaison exprime ce lien de manière remarquablement simple avec les \emph{inverses} de ces grandeurs.

\begin{propriete}[Formule de conjugaison des lentilles minces — démonstration exigible]
Pour une lentille mince utilisée dans les conditions de Gauss, un objet $AB$ et son image $A'B'$ vérifient :
\[ \frac{1}{\overline{OA'}} - \frac{1}{\overline{OA}} = \frac{1}{\overline{OF'}} \qquad \text{soit} \qquad \frac{1}{p'} - \frac{1}{p} = \frac{1}{f'}. \]
Cette relation est aussi appelée \cle{relation de Descartes}. Sa démonstration, à partir des triangles semblables de la construction géométrique, est exigible au Probatoire.
\end{propriete}

\courssub{Démonstration guidée par les triangles semblables}

\begin{experience}[Démonstration de la formule de conjugaison]
Considérons une lentille convergente, un objet réel $AB$ ($A$ sur l'axe, $AB \perp$ axe) donnant une image réelle $A'B'$. Traçons les deux rayons particuliers issus de $B$ :
\begin{itemize}
\item le rayon passant par le centre optique $O$, non dévié, qui rencontre l'image en $B'$ ;
\item le rayon parallèle à l'axe, qui émerge en passant par le foyer image $F'$ ; notons $J$ le point où ce rayon incident rencontre la lentille.
\end{itemize}

\textbf{Étape 1 : Triangles $A'B'F'$ et $OJF'$ (côté image).} Ces triangles sont semblables (angles droits en $A'$ et $J$, angle commun en $F'$). Comme le rayon incident $BJ$ est parallèle à l'axe, $OJ = AB$. Le théorème de Thalès donne :
\[ \frac{\overline{F'A'}}{\overline{F'O}} = \frac{\overline{A'B'}}{\overline{OJ}} = \frac{\overline{A'B'}}{\overline{AB}}. \tag{1} \]

\textbf{Étape 2 : Triangles $OAB$ et $OA'B'$ (rayon non dévié).} Les points $B$, $O$, $B'$ sont alignés. Les triangles sont semblables (angle commun en $O$, angles droits en $A$ et $A'$) :
\[ \frac{\overline{A'B'}}{\overline{AB}} = \frac{\overline{OA'}}{\overline{OA}}. \tag{2} \]

\textbf{Étape 3 : Égalisation des deux expressions du rapport $\dfrac{\overline{A'B'}}{\overline{AB}}$.} D'après (1) et (2) :
\[ \frac{\overline{F'A'}}{\overline{F'O}} = \frac{\overline{OA'}}{\overline{OA}}. \]
Or $\overline{F'A'} = \overline{OA'} - \overline{OF'}$ et $\overline{F'O} = -\overline{OF'}$. En substituant :
\[ \frac{\overline{OA'} - \overline{OF'}}{-\overline{OF'}} = \frac{\overline{OA'}}{\overline{OA}}. \]

\textbf{Étape 4 : Division par $\overline{OA'}$ (non nul).}
\[ \frac{1}{-\overline{OF'}} + \frac{1}{\overline{OA'}} = \frac{1}{\overline{OA}} \quad\Longrightarrow\quad \frac{1}{\overline{OA'}} - \frac{1}{\overline{OA}} = \frac{1}{\overline{OF'}}. \]
La formule est établie. Elle reste valable pour \emph{toutes} les situations (objet réel ou virtuel, image réelle ou virtuelle, lentille convergente ou divergente), à condition de respecter les signes algébriques.
\end{experience}

\begin{popfigure}
\begin{center}
\begin{tikzpicture}[scale=0.85,>=Stealth]
% Axe optique
\draw[->,thick,popdark] (-5.5,0) -- (7.5,0) node[below left] {axe optique};
% Lentille
\draw[thick,popblue] (0,-2.6) -- (0,2.6);
\draw[->,thick,popblue] (0,2.6) -- (-0.25,2.35);
\draw[->,thick,popblue] (0,2.6) -- (0.25,2.35);
\draw[->,thick,popblue] (0,-2.6) -- (-0.25,-2.35);
\draw[->,thick,popblue] (0,-2.6) -- (0.25,-2.35);
\node[popblue,above] at (0.1,2.7) {$L$};
% Points remarquables
\fill[popdark] (0,0) circle (1.5pt) node[below right] {$O$};
\fill[poporange] (2,0) circle (1.5pt) node[below] {$F'$};
\fill[poporange] (-2,0) circle (1.5pt) node[below] {$F$};
\fill[popink] (-3,0) circle (1.5pt) node[below] {$A$};
\fill[popgreen] (6,0) circle (1.5pt) node[below] {$A'$};
% Objet et image
\draw[->,very thick,popink] (-3,0) -- (-3,1.5) node[above] {$B$};
\draw[->,very thick,popgreen] (6,0) -- (6,-3) node[below] {$B'$};
% Point J sur la lentille
\fill[popblue] (0,1.5) circle (1.5pt) node[above right] {$J$};
% Rayons
\draw[thick,poporange] (-3,1.5) -- (0,1.5);
\draw[thick,poporange,->] (-3,1.5) -- (-1.2,1.5);
\draw[thick,poporange] (0,1.5) -- (6,-3);
\draw[thick,poporange,->] (0,1.5) -- (3.6,-1.1);
\draw[thick,poppurple] (-3,1.5) -- (6,-3);
\draw[thick,poppurple,->] (-3,1.5) -- (0,-0.75);
\draw[thick,poppurple,->] (0,-0.75) -- (3.6,-2.25);
% Cotes
\draw[<->,popmuted] (-3,-3.6) -- (0,-3.6) node[midway,below] {$p=\overline{OA}<0$};
\draw[<->,popmuted] (0,-3.6) -- (6,-3.6) node[midway,below] {$p'=\overline{OA'}>0$};
\draw[<->,popmuted] (0,3.0) -- (2,3.0) node[midway,above] {$f'=\overline{OF'}$};
\end{tikzpicture}
\end{center}
\legende{Construction servant à la démonstration : les triangles $OAB$ et $OA'B'$ (rayon non dévié par $O$, en violet) et les triangles $F'OJ$ et $F'A'B'$ (rayon parallèle à l'axe émergeant par $F'$, en orange) fournissent les deux expressions du rapport $\overline{A'B'}/\overline{AB}$.}
\end{popfigure}

\begin{savaistu}[La relation de Descartes]
La formule $\dfrac{1}{p'} - \dfrac{1}{p} = \dfrac{1}{f'}$ est souvent appelée \textbf{relation de Descartes} pour les lentilles minces, en hommage au philosophe et mathématicien français René Descartes (1596--1650), auteur de la \emph{Dioptrique} (1637), l'un des premiers traités scientifiques sur la formation des images. Descartes est aussi l'inventeur des coordonnées « cartésiennes » que tu utilises chaque jour en mathématiques : les abscisses algébriques $\overline{OA}$ et $\overline{OA'}$ de cette leçon sont précisément des coordonnées cartésiennes !
\end{savaistu}

\courssub{Formule du grandissement}

\begin{definition}[Grandissement transversal]
Le \cle{grandissement} (transversal) d'une lentille est le rapport algébrique
\[ \gamma = \frac{\overline{A'B'}}{\overline{AB}}. \]
C'est un nombre \textbf{sans dimension} (rapport de deux longueurs).
\end{definition}

\begin{propriete}[Expression pratique du grandissement — démonstration exigible]
Le grandissement s'exprime à partir des positions de l'objet et de l'image :
\[ \gamma = \frac{\overline{A'B'}}{\overline{AB}} = \frac{\overline{OA'}}{\overline{OA}} = \frac{p'}{p}. \]
\end{propriete}

\begin{experience}[Démonstration immédiate]
Le rayon issu de $B$ et passant par le centre optique $O$ n'est pas dévié : les points $B$, $O$ et $B'$ sont alignés. Les triangles $OAB$ et $OA'B'$ sont donc semblables (angle commun en $O$, angles droits en $A$ et $A'$). Le théorème de Thalès donne directement, en grandeurs algébriques :
\[ \frac{\overline{A'B'}}{\overline{AB}} = \frac{\overline{OA'}}{\overline{OA}}. \]
C'est la démonstration la plus courte du chapitre : une seule droite, celle du rayon non dévié !
\end{experience}

\courssub{Interprétation du grandissement}

Le signe et la valeur absolue de $\gamma$ renseignent entièrement sur l'image :

\begin{center}
\begin{tabularx}{\linewidth}{|l|X|X|}
\hline
\rowcolor{popblueL} \textbf{Cas} & \textbf{Signification} & \textbf{Exemple} \\
\hline
$\gamma > 0$ & image \textbf{droite} (même sens que l'objet) & loupe : $\gamma = +3$ \\
\hline
$\gamma < 0$ & image \textbf{renversée} (sens opposé) & projecteur : $\gamma = -2$ \\
\hline
$|\gamma| > 1$ & image \textbf{agrandie} & $\gamma = -2$ : deux fois plus grande \\
\hline
$|\gamma| < 1$ & image \textbf{réduite} & $\gamma = -0{,}5$ : deux fois plus petite \\
\hline
$|\gamma| = 1$ & image de même taille & $\gamma = -1$ (objet à $2f'$) \\
\hline
\end{tabularx}
\end{center}

\begin{aretenir}[Les deux formules à connaître par cœur]
\[ \boxed{\frac{1}{\overline{OA'}} - \frac{1}{\overline{OA}} = \frac{1}{\overline{OF'}}} \qquad \boxed{\gamma = \frac{\overline{A'B'}}{\overline{AB}} = \frac{\overline{OA'}}{\overline{OA}}} \]
Origine au centre optique $O$, axe orienté dans le sens de la lumière, distances en mètres pour la vergence $C = \dfrac{1}{f'}$.
\end{aretenir}

\courssub{Exploitation numérique : la méthode en 5 étapes}

\begin{methode}[Résoudre un problème de lentille en 5 étapes]
\begin{enumerate}
\item \textbf{Schématiser} : placer la lentille, $O$, $F$, $F'$, l'objet $AB$ ; orienter l'axe dans le sens de la lumière.
\item \textbf{Algébriser les données} : écrire $p = \overline{OA}$, $f' = \overline{OF'}$ avec leurs signes, et convertir en mètres.
\item \textbf{Appliquer la formule de position} : isoler $\dfrac{1}{p'} = \dfrac{1}{f'} + \dfrac{1}{p}$, calculer, puis inverser pour obtenir $p'$.
\item \textbf{Calculer le grandissement} : $\gamma = \dfrac{p'}{p}$, puis $\overline{A'B'} = \gamma \times \overline{AB}$.
\item \textbf{Conclure en français} : nature de l'image (réelle si $p'>0$, virtuelle si $p'<0$), sens (droite si $\gamma>0$, renversée si $\gamma<0$), taille (comparer $|\gamma|$ à 1). Vérifier la cohérence avec une construction géométrique rapide.
\end{enumerate}
\end{methode}

\begin{exemplebox}[Objet à 30 cm d'une lentille de $f' = +20$ cm]
Un objet $AB$ de hauteur $2\ \mathrm{cm}$ est placé à $30\ \mathrm{cm}$ devant une lentille convergente de distance focale $f' = +20\ \mathrm{cm}$.

\textbf{Données algébriques :} $p = \overline{OA} = \num{-0,30}\ \mathrm{m}$ ; $f' = \num{0,20}\ \mathrm{m}$.

\textbf{Formule de position :}
\[ \frac{1}{p'} = \frac{1}{f'} + \frac{1}{p} = \frac{1}{\num{0,20}} + \frac{1}{\num{-0,30}} = 5 - \frac{10}{3} = \frac{5}{3}\ \mathrm{m^{-1}} \approx 1{,}67\ \mathrm{m^{-1}} \]
\[ p' = \frac{3}{5} = \num{0,60}\ \mathrm{m} = +60\ \mathrm{cm}. \]

\textbf{Grandissement :}
\[ \gamma = \frac{p'}{p} = \frac{+60}{-30} = -2 \quad\Longrightarrow\quad \overline{A'B'} = -2 \times 2 = -4\ \mathrm{cm}. \]

\textbf{Conclusion :} l'image est \textbf{réelle} ($p' > 0$), située à $60\ \mathrm{cm}$ derrière la lentille, \textbf{renversée} ($\gamma < 0$) et \textbf{deux fois plus grande} que l'objet ($|\gamma| = 2$). C'est cohérent avec la construction géométrique : l'objet est entre $F$ et $2F$ (car $f' < |p| < 2f'$), donc l'image est au-delà de $2F'$, réelle, renversée et agrandie.
\end{exemplebox}

\begin{exoresolu}[Lentille utilisée en loupe]
Un timbre de hauteur $AB = 1{,}0\ \mathrm{cm}$ est placé devant une lentille convergente de vergence $C = +20\ \delta$. On étudie deux positions de l'objet.
\tcblower
\textbf{Données communes :} $f' = \dfrac{1}{C} = \dfrac{1}{20} = \num{0,050}\ \mathrm{m} = 5{,}0\ \mathrm{cm}$.

\textbf{Cas 1 : Objet à $p = -8{,}0\ \mathrm{cm}$ (au-delà du foyer).}
$p = \num{-0,080}\ \mathrm{m}$.
\[ \frac{1}{p'} = \frac{1}{f'} + \frac{1}{p} = 20 + \frac{1}{\num{-0,080}} = 20 - 12{,}5 = 7{,}5\ \mathrm{m^{-1}} \]
\[ p' = \frac{1}{7{,}5} \approx \num{0,133}\ \mathrm{m} = +13{,}3\ \mathrm{cm}. \]
$\gamma = \dfrac{p'}{p} = \dfrac{+13{,}3}{-8{,}0} \approx -1{,}67$. Image \textbf{réelle}, \textbf{renversée}, agrandie. Ce n'est \textbf{pas} le fonctionnement d'une loupe (l'objet est hors du foyer).

\textbf{Cas 2 : Objet à $p = -4{,}0\ \mathrm{cm}$ (entre $O$ et $F$, régime loupe).}
$p = \num{-0,040}\ \mathrm{m}$.
\[ \frac{1}{p'} = 20 + \frac{1}{\num{-0,040}} = 20 - 25 = -5{,}0\ \mathrm{m^{-1}} \quad\Longrightarrow\quad p' = \num{-0,20}\ \mathrm{m} = -20\ \mathrm{cm}. \]
\[ \gamma = \frac{p'}{p} = \frac{-20}{-4} = +5 \quad\Longrightarrow\quad \overline{A'B'} = 5 \times 1{,}0 = +5{,}0\ \mathrm{cm}. \]

\textbf{Conclusion :} L'image est \textbf{virtuelle} ($p' < 0$, elle se forme du même côté que l'objet), \textbf{droite} ($\gamma > 0$) et \textbf{cinq fois plus grande} : c'est bien le fonctionnement d'une \textbf{loupe}. La cohérence avec la construction géométrique est parfaite : objet entre $O$ et $F$, les rayons émergents divergent et leurs prolongements se coupent en $B'$ avant la lentille.
\end{exoresolu}

\courssub{Vérification expérimentale de la formule de position au banc d'optique}

\begin{experience}[Vérification de la relation de conjugaison]
\textbf{Matériel :} banc d'optique gradué, lanterne avec objet lumineux (lettre F découpée), lentille convergente de vergence inconnue, écran translucide.

\textbf{Protocole :}
\begin{enumerate}
\item Réaliser les conditions de Gauss : objet, lentille et écran centrés sur l'axe optique, plans parallèles.
\item Placer l'objet à une position connue ; mesurer l'abscisse $x_A$ de l'objet et $x_O$ de la lentille ; en déduire $p = \overline{OA} = x_A - x_O$.
\item Déplacer l'écran jusqu'à obtenir une image \emph{nette} ; mesurer $x_{A'}$ et calculer $p' = \overline{OA'} = x_{A'} - x_O$.
\item Recommencer pour au moins cinq positions de l'objet.
\item Pour chaque mesure, calculer $\dfrac{1}{p}$ et $\dfrac{1}{p'}$ (en $\mathrm{m^{-1}}$).
\end{enumerate}

\textbf{Exploitation :} on trace $\dfrac{1}{p'}$ en fonction de $\dfrac{1}{p}$. D'après la formule de conjugaison :
\[ \frac{1}{p'} = \frac{1}{p} + \frac{1}{f'} \]
les points doivent s'aligner sur une \textbf{droite de pente $+1$} dont l'ordonnée à l'origine vaut $\dfrac{1}{f'} = C$. On vérifie ainsi la formule et on détermine simultanément la vergence de la lentille.
\end{experience}

\begin{popfigure}
\begin{center}
\begin{tikzpicture}
\begin{axis}[
width=0.8\linewidth, height=6.5cm,
xlabel={$\dfrac{1}{p}$ en $\mathrm{m^{-1}}$},
ylabel={$\dfrac{1}{p'}$ en $\mathrm{m^{-1}}$},
xmin=-12, xmax=-2, ymin=2, ymax=12,
grid=both, grid style={popline!40},
axis lines=left, thick,
legend style={at={(0.03,0.97)}, anchor=north west, font=\small}
]
\addplot[popcurve, domain=-11:-3, samples=50] {x + 10};
\addlegendentry{$\dfrac{1}{p'} = \dfrac{1}{p} + \dfrac{1}{f'}$}
\addplot[only marks, mark=*, mark size=2pt, poporange] coordinates {(-10,0.4) (-8,2.1) (-6.67,3.4) (-5,5.2) (-4,6.1) (-3.33,6.8)};
\addlegendentry{mesures expérimentales}
\addplot[only marks, mark=square*, mark size=2.5pt, popgreen] coordinates {(0,10)};
\node[popgreen, anchor=west, font=\small] at (axis cs:-11.5,10.3) {ordonnée à l'origine $= \dfrac{1}{f'} = C$};
\end{axis}
\end{tikzpicture}
\end{center}
\legende{Vérification expérimentale : les points $\left(\frac{1}{p}, \frac{1}{p'}\right)$ s'alignent sur une droite de pente $+1$ ; son ordonnée à l'origine donne la vergence $C = \frac{1}{f'}$ (ici $C \approx 10\ \delta$, soit $f' \approx 10\ \mathrm{cm}$).}
\end{popfigure}

\courssub{Détermination expérimentale de la vergence d'une lentille}

\courssub{Méthode de l'objet à l'infini (méthode rapide)}

\begin{methode}[Méthode de l'objet lointain]
\begin{enumerate}
\item Viser avec la lentille un objet très éloigné (un bâtiment, un arbre dans la cour du lycée, la flamme d'une lampe placée au fond de la salle) : on considère $p \to -\infty$, donc $\dfrac{1}{p} \to 0$.
\item La formule de conjugaison donne alors $\dfrac{1}{p'} = \dfrac{1}{f'}$, soit $p' = f'$ : l'image se forme \textbf{dans le plan focal image}.
\item Recueillir l'image nette sur un écran (ou une feuille blanche) et mesurer directement la distance lentille--écran : c'est $f'$.
\item Calculer la vergence : $C = \dfrac{1}{f'}$ avec $f'$ en mètres.
\end{enumerate}
\textbf{Exemple :} l'image du portail du lycée se forme nette à $12{,}5\ \mathrm{cm}$ derrière la lentille : $f' = \num{0,125}\ \mathrm{m}$ et $C = \dfrac{1}{\num{0,125}} = +8{,}0\ \delta$.
\end{methode}

\courssub{Méthode par autocollimation (complément au programme)}

\begin{experience}[Autocollimation : la méthode de précision]
\textbf{Principe :} on place un miroir plan derrière la lentille, perpendiculairement à l'axe optique. On déplace l'objet lumineux $AB$ jusqu'à ce que son image, après un aller-retour à travers la lentille, se forme \emph{nette dans le plan de l'objet lui-même}, renversée.

\textbf{Interprétation :} cela n'est possible que si les rayons frappent le miroir perpendiculairement, c'est-à-dire s'ils sont parallèles à l'axe entre la lentille et le miroir. Or des rayons émergent parallèles à l'axe si et seulement si l'objet est \textbf{dans le plan focal objet} : alors $A$ coïncide avec $F$.

\textbf{Mesure :} la distance objet--lentille est exactement $f'$, indépendante de la position du miroir. On mesure $f'$, puis $C = \dfrac{1}{f'}$. Cette méthode, plus précise que la précédente, est celle utilisée par les opticiens pour vérifier la vergence des verres correcteurs.
\end{experience}

\begin{attention}[Erreurs fréquentes lors des mesures au banc d'optique]
\begin{itemize}
\item Mesurer les distances à partir de l'\textbf{écran ou de l'objet} au lieu du \textbf{centre optique $O$} : toutes les distances $p$, $p'$, $f'$ ont leur origine en $O$.
\item Oublier de convertir en mètres : $f' = 20\ \mathrm{cm}$ donne $C = \dfrac{1}{\num{0,20}} = 5\ \delta$, et non $C = \dfrac{1}{20}$ !
\item Conclure « image réelle » parce que $p'$ est grand : c'est le \textbf{signe} de $p'$ qui décide, pas sa valeur.
\item Confondre $\gamma$ et $|\gamma|$ : un grandissement de $-2$ signifie « renversée et deux fois plus grande », pas « deux fois plus petite ».
\end{itemize}
\end{attention}

\begin{exercice}[Pour t'entraîner]
Une lentille convergente a pour vergence $C = +4{,}0\ \delta$. Un objet réel de hauteur $3{,}0\ \mathrm{cm}$ est placé à $50\ \mathrm{cm}$ devant la lentille.
\begin{enumerate}
\item Calculer la distance focale $f'$.
\item Déterminer la position $p'$ de l'image.
\item Calculer le grandissement et la taille de l'image ; préciser sa nature et son sens.
\item Vérifier la cohérence du résultat avec la position de l'objet par rapport à $F$ et $2F$.
\end{enumerate}
\end{exercice}

\begin{corrige}[Exercice n°1]
\textbf{1.} $f' = \dfrac{1}{C} = \dfrac{1}{4{,}0} = \num{0,25}\ \mathrm{m} = 25\ \mathrm{cm}$.

\textbf{2.} $p = \num{-0,50}\ \mathrm{m}$ :
\[ \frac{1}{p'} = \frac{1}{f'} + \frac{1}{p} = 4{,}0 + \frac{1}{\num{-0,50}} = 4{,}0 - 2{,}0 = 2{,}0\ \mathrm{m^{-1}} \quad\Longrightarrow\quad p' = +\num{0,50}\ \mathrm{m} = +50\ \mathrm{cm}. \]

\textbf{3.} $\gamma = \dfrac{p'}{p} = \dfrac{+50}{-50} = -1$ ; donc $\overline{A'B'} = -1 \times 3{,}0 = -3{,}0\ \mathrm{cm}$. L'image est \textbf{réelle}, \textbf{renversée} et de \textbf{même taille} que l'objet.

\textbf{4.} L'objet est à $50\ \mathrm{cm} = 2f'$ de la lentille : c'est le cas remarquable où l'image se forme à $2f'$ de l'autre côté, avec $\gamma = -1$. Le résultat est parfaitement cohérent avec la construction géométrique (configuration dite « en $2f$ », utilisée en photographie pour la reproduction à l'échelle 1).
\end{corrige}

En résumé, les deux formules $\dfrac{1}{p'} - \dfrac{1}{p} = \dfrac{1}{f'}$ et $\gamma = \dfrac{p'}{p}$, démontrées à partir de la construction géométrique, permettent de résoudre \emph{tous} les problèmes de lentille mince sans tracer un seul rayon — à condition de respecter scrupuleusement la convention des signes. Nous verrons dans la section suivante comment ces formules se généralisent élégamment aux associations de lentilles grâce au théorème des vergences.

\coursec{Association de lentilles : théorème des vergences}

Dans la section précédente, nous avons établi les formules de position et de grandissement pour une lentille mince unique : $\dfrac{1}{\overline{OA'}} - \dfrac{1}{\overline{OA}} = \dfrac{1}{f'} = C$. Ces formules suffisent à étudier une loupe ou un objectif simple. Mais dans la réalité, les instruments d'optique -- lunettes, appareils photo, microscopes -- combinent presque toujours \emph{plusieurs} lentilles. Chez l'opticien de Yaoundé ou de Douala, un verre correcteur complexe est souvent obtenu en associant deux verres. La question qui guide cette section est donc naturelle : \emph{à quelle lentille unique un système de plusieurs lentilles est-il équivalent ?} Nous allons voir que la réponse est d'une élégance remarquable lorsque les lentilles sont \cle{accolées} : c'est le \cle{théorème des vergences}.

\courssub{Système de deux lentilles minces accolées}

\begin{definition}[Lentilles accolées]
Deux lentilles minces $L_1$ et $L_2$ sont dites \cle{accolées} lorsqu'elles ont le \emph{même axe optique principal} et que leurs centres optiques $O_1$ et $O_2$ sont \emph{pratiquement confondus} : $O_1 \approx O_2 = O$. On peut alors les considérer comme une seule lentille mince de centre $O$.
\end{definition}

L'idée physique est simple : chaque lentille dévie les rayons lumineux. Lorsque les deux lentilles sont au même endroit, leurs déviations s'ajoutent. Or la déviation produite par une lentille est d'autant plus grande que sa vergence est élevée. On devine donc que les vergences vont s'additionner. Vérifions-le rigoureusement.

\begin{popfigure}
\begin{tikzpicture}[scale=1.0, >=Stealth]
  % axe optique
  \draw[->, popdark] (-5,0) -- (5.6,0) node[below left]{$\Delta$};
  % Lentille L1
  \draw[popblue, very thick] (-0.35,-1.8) .. controls (-0.75,0) .. (-0.35,1.8);
  \draw[popblue, very thick] (-0.35,-1.8) .. controls (0.05,0) .. (-0.35,1.8);
  \node[popblue] at (-1.1,2.15){$L_1$};
  \node[popblue] at (-1.1,-2.15){$C_1$};
  % Lentille L2
  \draw[poporange, very thick] (0.35,-1.8) .. controls (0.05,0) .. (0.35,1.8);
  \draw[poporange, very thick] (0.35,-1.8) .. controls (0.75,0) .. (0.35,1.8);
  \node[poporange] at (1.1,2.15){$L_2$};
  \node[poporange] at (1.1,-2.15){$C_2$};
  % centre optique commun
  \fill[popdark] (0,0) circle (0.06) node[below right=1pt]{$O_1 \approx O_2 = O$};
  % lentille équivalente
  \draw[popgreen, very thick] (3.6,-1.8) .. controls (3.2,0) .. (3.6,1.8);
  \draw[popgreen, very thick] (3.6,-1.8) .. controls (4.0,0) .. (3.6,1.8);
  \node[popgreen] at (3.6,2.15){$L$};
  \node[popgreen] at (3.6,-2.15){$C = C_1 + C_2$};
  \node at (2.0,0.6){\Large $\equiv$};
\end{tikzpicture}
\legende{Deux lentilles minces accolées $L_1$ (vergence $C_1$) et $L_2$ (vergence $C_2$), de centres optiques confondus en $O$, sont équivalentes à une lentille unique $L$ de vergence $C = C_1 + C_2$.}
\end{popfigure}

\courssub{Théorème des vergences : énoncé et démonstration guidée}

\begin{propriete}[Théorème des vergences]
Un ensemble de deux lentilles minces \emph{accolées}, de vergences $C_1$ et $C_2$, est équivalent à une lentille mince unique de même centre optique, dont la vergence est la somme algébrique des vergences :
\[ C = C_1 + C_2 \qquad \text{soit} \qquad \frac{1}{f'} = \frac{1}{f'_1} + \frac{1}{f'_2}. \]
La démonstration, fondée sur deux applications successives de la formule de conjugaison, est exigible au Probatoire.
\end{propriete}

\begin{experience}[Démonstration guidée du théorème des vergences]
\textbf{Hypothèses.} Deux lentilles minces $L_1$ et $L_2$ accolées, de centres $O_1 \approx O_2 = O$, de distances focales images $f'_1$ et $f'_2$. Un objet $AB$ est placé devant le système, dans les conditions de Gauss.

\textbf{Étape 1 : rôle de la première lentille.} $L_1$ donne de l'objet $AB$ une image $A_1B_1$. La formule de conjugaison appliquée à $L_1$, avec l'origine en $O$, s'écrit :
\[ \frac{1}{\overline{OA_1}} - \frac{1}{\overline{OA}} = \frac{1}{f'_1} = C_1. \tag{1} \]

\textbf{Étape 2 : rôle de la deuxième lentille.} L'image $A_1B_1$ joue le rôle d'\emph{objet} pour $L_2$, qui en donne l'image définitive $A'B'$. Comme $O_1 \approx O_2 = O$, la formule de conjugaison appliquée à $L_2$ s'écrit avec la même origine $O$ :
\[ \frac{1}{\overline{OA'}} - \frac{1}{\overline{OA_1}} = \frac{1}{f'_2} = C_2. \tag{2} \]

\textbf{Étape 3 : addition membre à membre.} Additionnons les relations (1) et (2) :
\[ \frac{1}{\overline{OA_1}} - \frac{1}{\overline{OA}} + \frac{1}{\overline{OA'}} - \frac{1}{\overline{OA_1}} = C_1 + C_2. \]
Le terme $\dfrac{1}{\overline{OA_1}}$ s'élimine (c'est le rôle charnière de l'image intermédiaire) :
\[ \frac{1}{\overline{OA'}} - \frac{1}{\overline{OA}} = C_1 + C_2. \]

\textbf{Étape 4 : conclusion.} La relation obtenue est exactement la formule de conjugaison d'une lentille unique de centre $O$ et de vergence $C = C_1 + C_2$, donnant de $AB$ l'image $A'B'$. Le système accolé se comporte donc comme une lentille unique de vergence $C = C_1 + C_2$. \hfill $\blacksquare$
\end{experience}

\begin{popfigure}
\begin{tikzpicture}[scale=1.0, >=Stealth]
  % axe
  \draw[->, popdark] (-5.5,0) -- (6.5,0);
  % lentilles accolées
  \draw[popblue, very thick] (-0.3,-2) .. controls (-0.7,0) .. (-0.3,2);
  \draw[popblue, very thick] (-0.3,-2) .. controls (0.1,0) .. (-0.3,2);
  \draw[poporange, very thick] (0.3,-2) .. controls (-0.1,0) .. (0.3,2);
  \draw[poporange, very thick] (0.3,-2) .. controls (0.7,0) .. (0.3,2);
  \fill[popdark] (0,0) circle (0.06) node[below left=1pt]{$O$};
  % foyer image équivalent
  \fill[popgreen] (3,0) circle (0.07) node[below]{$F'$};
  % rayon incident parallèle
  \draw[very thick, poppurple, ->] (-5,1.4) -- (0,1.4);
  % rayon dévié vers F'
  \draw[very thick, poppurple, ->] (0,1.4) -- (3,0);
  \draw[very thick, poppurple] (3,0) -- (4.6,-0.85);
  % prolongement sans lentilles (pointillés)
  \draw[dashed, popmuted] (0,1.4) -- (5.2,1.4);
  \node[popmuted] at (4.6,1.7){sans lentilles};
  % trajet après L1 seule
  \draw[popblue, dashed] (0,1.4) -- (2.6,0.35);
  \node[popblue] at (2.5,1.05){après $L_1$ seule};
  % annotations
  \node[poppurple] at (-3.4,1.75){rayon parallèle à l'axe};
  \node[popgreen] at (4.4,-0.4){émerge en passant par $F'$};
\end{tikzpicture}
\legende{Trajet d'un rayon parallèle à l'axe à travers deux lentilles accolées : chaque lentille ajoute sa déviation ; le rayon émerge en passant par le foyer image $F'$ du système équivalent, avec $\dfrac{1}{\overline{OF'}} = C_1 + C_2$.}
\end{popfigure}

\begin{propriete}[Généralisation à $n$ lentilles accolées]
Pour $n$ lentilles minces accolées sur le même axe, de vergences $C_1, C_2, \ldots, C_n$, la vergence du système équivalent est :
\[ C = \sum_{i=1}^{n} C_i = C_1 + C_2 + \cdots + C_n. \]
On le démontre par récurrence en accolant les lentilles une à une.
\end{propriete}

\begin{attention}[Algébricité des vergences]
Les vergences s'additionnent \emph{algébriquement} : une lentille convergente a $C > 0$, une lentille divergente a $C < 0$. Ne jamais additionner les distances focales : c'est faux en général ! On additionne les \emph{inverses} des distances focales : $\dfrac{1}{f'} = \dfrac{1}{f'_1} + \dfrac{1}{f'_2}$. Par exemple, deux lentilles de $f'_1 = \SI{20}{cm}$ et $f'_2 = \SI{20}{cm}$ accolées donnent $\dfrac{1}{f'} = \dfrac{1}{20} + \dfrac{1}{20} = \dfrac{1}{10}$, soit $f' = \SI{10}{cm}$, et non $\SI{40}{cm}$.
\end{attention}

\courssub{Conséquences : système convergent, divergent ou afocal}

Puisque les vergences s'ajoutent algébriquement, l'association d'une lentille convergente ($C_1 > 0$) et d'une lentille divergente ($C_2 < 0$) peut donner trois cas :

\begin{itemize}
\item si $|C_1| > |C_2|$ : $C = C_1 + C_2 > 0$, le système est \cle{convergent} ;
\item si $|C_1| < |C_2|$ : $C < 0$, le système est \cle{divergent} ;
\item si $C_1 + C_2 = 0$ : $C = 0$, c'est-à-dire $f' \to \infty$ : le système est \cle{afocal}. Un faisceau parallèle incident en ressort parallèle (mais déplacé ou élargi). C'est le principe de certaines lunettes astronomiques.
\end{itemize}

\begin{exemplebox}[Association convergent + divergent]
On accole une lentille $L_1$ de vergence $C_1 = +3\ \delta$ et une lentille $L_2$ de vergence $C_2 = -1\ \delta$. La vergence du système est :
\[ C = C_1 + C_2 = 3 - 1 = +2\ \delta, \qquad f' = \frac{1}{C} = \frac{1}{2} = \SI{0,50}{m}. \]
Le système est \emph{convergent}, de distance focale image $\SI{50}{cm}$ : la lentille convergente « l'emporte ». Si l'on avait pris $C_2 = -5\ \delta$, on aurait obtenu $C = -2\ \delta$ : système divergent. Avec $C_2 = -3\ \delta$ : système afocal.
\end{exemplebox}

\begin{exoresolu}[Vérification de la formule de position pour un système accolé]
\textbf{Énoncé.} Deux lentilles minces accolées $L_1$ ($f'_1 = \SI{25}{cm}$) et $L_2$ ($f'_2 = -\SI{50}{cm}$) forment un système $S$. Un objet $AB$ de hauteur $\SI{2}{cm}$ est placé à $\SI{100}{cm}$ devant le centre optique commun $O$. Déterminer la position, la nature et la taille de l'image $A'B'$ donnée par le système.
\tcblower
\textbf{Solution.} \emph{Étape 1 : vergences.} $C_1 = \dfrac{1}{f'_1} = \dfrac{1}{0,25} = +4\ \delta$ et $C_2 = \dfrac{1}{f'_2} = \dfrac{1}{-0,50} = -2\ \delta$.

\emph{Étape 2 : vergence équivalente.} Les lentilles étant accolées :
\[ C = C_1 + C_2 = 4 - 2 = +2\ \delta \quad \Rightarrow \quad f' = \frac{1}{C} = \SI{0,50}{m} = \SI{50}{cm}. \]
Le système est convergent.

\emph{Étape 3 : position de l'image.} Avec $\overline{OA} = -\SI{100}{cm} = -\SI{1}{m}$, la formule de conjugaison donne :
\[ \frac{1}{\overline{OA'}} = C + \frac{1}{\overline{OA}} = 2 + \frac{1}{-1} = 1 \quad \Rightarrow \quad \overline{OA'} = +\SI{1}{m} = +\SI{100}{cm}. \]
L'image est \emph{réelle} ($\overline{OA'} > 0$), située $\SI{100}{cm}$ derrière le système.

\emph{Étape 4 : grandissement et taille.}
\[ \gamma = \frac{\overline{OA'}}{\overline{OA}} = \frac{+1}{-1} = -1 \quad \Rightarrow \quad \overline{A'B'} = \gamma \cdot \overline{AB} = -\SI{2}{cm}. \]
L'image est renversée, de même taille que l'objet ($\SI{2}{cm}$).

\emph{Vérification de cohérence.} L'objet est à $2f'$ du système ($\SI{100}{cm} = 2 \times \SI{50}{cm}$) : on retrouve le cas bien connu où l'image se forme à $2f'$ derrière la lentille, renversée et de même taille ($\gamma = -1$). Le résultat est cohérent.
\end{exoresolu}

\courssub{Lentilles non accolées : méthode de l'image intermédiaire}

\begin{attention}[Le théorème des vergences exige des lentilles accolées]
Si les lentilles sont \emph{séparées} par une distance $e = \overline{O_1O_2}$ non négligeable, le théorème des vergences \textbf{ne s'applique plus} : $C \neq C_1 + C_2$ en général. En effet, dans la démonstration, c'est l'hypothèse $O_1 \approx O_2$ qui permet d'écrire les deux formules de conjugaison avec la \emph{même origine} et d'éliminer le terme intermédiaire. Pour des lentilles séparées, on revient systématiquement à la méthode de l'\cle{image intermédiaire}.
\end{attention}

\begin{methode}[Étudier un système de deux lentilles séparées]
\begin{enumerate}
\item \textbf{Première lentille.} Appliquer la formule de conjugaison à $L_1$ avec l'origine en $O_1$ : on obtient $\overline{O_1A_1}$, position de l'image intermédiaire $A_1B_1$, et le grandissement $\gamma_1 = \dfrac{\overline{O_1A_1}}{\overline{O_1A}}$.
\item \textbf{Changement d'origine.} $A_1B_1$ devient l'objet pour $L_2$. Calculer sa position par rapport à $O_2$ grâce à la relation de Chasles sur l'axe : $\overline{O_2A_1} = \overline{O_2O_1} + \overline{O_1A_1}$. Attention aux signes : $\overline{O_2O_1} = -\overline{O_1O_2} = -e$.
\item \textbf{Deuxième lentille.} Appliquer la formule de conjugaison à $L_2$ avec l'origine en $O_2$ : on obtient $\overline{O_2A'}$, position de l'image définitive, et $\gamma_2 = \dfrac{\overline{O_2A'}}{\overline{O_2A_1}}$.
\item \textbf{Grandissement total.} $\gamma = \gamma_1 \times \gamma_2$, d'où la taille de l'image définitive $\overline{A'B'} = \gamma \cdot \overline{AB}$.
\item \textbf{Construction géométrique.} Construire d'abord $A_1B_1$ à travers $L_1$ (deux rayons particuliers), puis prolonger les rayons utiles jusqu'à $L_2$ et construire l'image définitive à travers $L_2$.
\end{enumerate}
\end{methode}

\begin{exoresolu}[Deux lentilles séparées]
\textbf{Énoncé.} Une lentille convergente $L_1$ ($f'_1 = \SI{20}{cm}$, centre $O_1$) et une lentille convergente $L_2$ ($f'_2 = \SI{10}{cm}$, centre $O_2$) ont le même axe et sont distantes de $\overline{O_1O_2} = +\SI{60}{cm}$. Un objet $AB$ de $\SI{4}{cm}$ est placé à $\SI{30}{cm}$ devant $O_1$. Déterminer la position et la taille de l'image définitive.
\tcblower
\textbf{Solution.} \emph{Étape 1 : image par $L_1$.} $\overline{O_1A} = -\SI{30}{cm} = -\SI{0,30}{m}$, $C_1 = \dfrac{1}{0,20} = +5\ \delta$.
\[ \frac{1}{\overline{O_1A_1}} = C_1 + \frac{1}{\overline{O_1A}} = 5 - \frac{1}{0,30} = 5 - \frac{10}{3} = \frac{5}{3} \ \Rightarrow \ \overline{O_1A_1} = +\SI{0,60}{m} = +\SI{60}{cm}. \]
$\gamma_1 = \dfrac{\overline{O_1A_1}}{\overline{O_1A}} = \dfrac{60}{-30} = -2$, donc $\overline{A_1B_1} = \gamma_1 \cdot \overline{AB} = -\SI{8}{cm}$ (image renversée, deux fois plus grande).

\emph{Étape 2 : changement d'origine.} $\overline{O_2A_1} = \overline{O_2O_1} + \overline{O_1A_1} = -60 + 60 = \SI{0}{cm}$ : l'image intermédiaire se forme \emph{exactement au centre} $O_2$ de la deuxième lentille.

\emph{Étape 3 : image par $L_2$.} Un objet placé au centre optique d'une lentille n'est pas dévié (tout rayon passant par $O_2$ traverse la lentille sans déviation) : l'image définitive se forme donc en $O_2$, avec $\gamma_2 = +1$.

\emph{Étape 4 : grandissement total.} $\gamma = \gamma_1 \gamma_2 = -2 \times 1 = -2$, donc $\overline{A'B'} = -\SI{8}{cm}$ : image réelle (on peut la recueillir sur un écran placé en $O_2$), renversée, de $\SI{8}{cm}$.
\end{exoresolu}

\courssub{Application : la prescription de l'opticien}

Lorsqu'un ophtalmologue prescrit une correction, il indique une \emph{vergence} (en dioptries) : par exemple $-2\ \delta$ pour un myope, $+1,5\ \delta$ pour un hypermétrope. Le théorème des vergences est alors l'outil de travail de l'opticien : pour obtenir une vergence prescrite $C$, il peut accoler des verres de vergences connues, par exemple un verre sphérique $C_1$ et un verre cylindrique correcteur d'astigmatisme $C_2$, de sorte que $C_1 + C_2 = C$. C'est aussi le principe des \emph{verres à double foyer} (bifocaux) : la partie supérieure du verre et la partie inférieure n'ont pas la même vergence, ce qui permet au presbyte de voir de loin et de près avec la même paire de lunettes.

\begin{savaistu}[Des verres composés aux doubles foyers]
C'est Benjamin Franklin, au XVIII\textsuperscript{e} siècle, qui aurait inventé les lunettes à double foyer, fatigué de changer sans cesse de paires pour lire et pour regarder au loin. Aujourd'hui, les verres \emph{progressifs} modernes réalisent une variation continue de la vergence du haut vers le bas du verre -- une prouesse de fabrication qui repose toujours, en arrière-plan, sur l'addition des vergences. Au Cameroun, les campagnes de santé visuelle organisées dans les écoles et les hôpitaux (par exemple à l'Hôpital Central de Yaoundé) reposent sur la mesure de la vergence correctrice à l'aide de \emph{lunettes d'essai} : le médecin accole des verres de vergences connues devant l'œil du patient jusqu'à obtention d'une vision nette -- une application directe et quotidienne du théorème des vergences !
\end{savaistu}

\begin{exercice}[Pour s'entraîner]
On accole trois lentilles minces de distances focales images $f'_1 = \SI{20}{cm}$, $f'_2 = \SI{50}{cm}$ et $f'_3 = -\SI{25}{cm}$.
\begin{enumerate}
\item Calculer la vergence du système équivalent et sa distance focale image. Le système est-il convergent ou divergent ?
\item Un objet réel est placé à $\SI{40}{cm}$ devant le centre optique commun. Déterminer la position et la nature de l'image.
\end{enumerate}
\end{exercice}

\begin{corrige}[Exercice : association de trois lentilles]
\begin{enumerate}
\item Vergences : $C_1 = \dfrac{1}{0,20} = +5\ \delta$, $C_2 = \dfrac{1}{0,50} = +2\ \delta$, $C_3 = \dfrac{1}{-0,25} = -4\ \delta$. D'où $C = 5 + 2 - 4 = +3\ \delta$ et $f' = \dfrac{1}{3} \approx \SI{0,33}{m}$. Le système est \emph{convergent}.
\item $\overline{OA} = -\SI{0,40}{m}$ : $\dfrac{1}{\overline{OA'}} = C + \dfrac{1}{\overline{OA}} = 3 - \dfrac{1}{0,40} = 3 - 2,5 = 0,5$, d'où $\overline{OA'} = +\SI{2}{m}$. L'image est \emph{réelle}, située $\SI{2}{m}$ derrière le système.
\end{enumerate}
\end{corrige}

\begin{aretenir}[L'essentiel de la section]
\begin{itemize}
\item Deux lentilles minces \textbf{accolées} (même axe, centres confondus) équivalent à une lentille unique de vergence $C = C_1 + C_2$ ; en général, $C = \sum_i C_i$.
\item Les vergences s'additionnent \emph{algébriquement} ; on n'additionne jamais les distances focales.
\item Selon le signe de $C_1 + C_2$, le système est convergent, divergent ou afocal ($C = 0$).
\item Pour des lentilles \textbf{séparées}, le théorème ne s'applique pas : on utilise la méthode de l'image intermédiaire, avec un changement d'origine soigneux entre $O_1$ et $O_2$ (relation de Chasles), et $\gamma = \gamma_1 \gamma_2$.
\item Application : l'opticien combine des verres pour réaliser exactement la vergence prescrite.
\end{itemize}
\end{aretenir}

\coursec{Applications des lentilles minces}

Dans la section précédente, nous avons vu que l'association de deux lentilles minces accolées se comporte comme une lentille unique dont la vergence est la somme des vergences : $C = C_1 + C_2$. Ce résultat, le théorème des vergences, n'est pas qu'un exercice d'algèbre : c'est précisément le principe qui permet à un opticien de corriger la vision d'un élève de Première en posant un verre correcteur devant son œil, ou à un fabricant d'appareils photo de composer un objectif performant. Toutes les applications des lentilles minces — loupe, œil, appareil photographique, lunette, microscope — reposent sur les mêmes outils que tu maîtrises maintenant : les conditions de Gauss, la construction géométrique des images, la formule de position et le grandissement. Cette section les met en situation.

\courssub{La loupe : une lentille convergente utilisée en loupe}

\textbf{Observation.} Un horloger de Yaoundé qui répare une montre, un élève qui observe les nervures d'une feuille au laboratoire : tous deux rapprochent une petite lentille épaisse de l'objet et voient celui-ci \emph{agrandi et à l'endroit}. Pourquoi l'image est-elle droite alors que la lentille convergente donne d'habitude des images renversées ?

\begin{definition}[La loupe]
Une \cle{loupe} est une lentille mince \cle{convergente} de courte distance focale (typiquement $f' \in [1\ \text{cm}\ ;\ 10\ \text{cm}]$, donc de forte vergence $C = \dfrac{1}{f'}$ comprise entre $10\ \delta$ et $100\ \delta$). On l'utilise en plaçant l'objet \emph{entre le foyer objet $F$ et le centre optique $O$}, c'est-à-dire en \emph{fonctionnement en loupe}.
\end{definition}

\begin{propriete}[Fonctionnement de la loupe]
Lorsque l'objet réel $AB$ est placé entre $F$ et $O$ ($-f' < \overline{OA} < 0$), la lentille convergente donne une image $A'B'$ :
\begin{itemize}
\item \cle{virtuelle} (les rayons émergents divergent, seuls leurs prolongements se croisent) ;
\item \cle{droite} (de même sens que l'objet, $\gamma > 0$) ;
\item \cle{agrandie} ($|\gamma| > 1$).
\end{itemize}
L'œil placé derrière la loupe reçoit les rayons émergents comme s'ils provenaient de $A'B'$ : c'est cette image virtuelle agrandie que l'on « voit ». La démonstration s'obtient immédiatement par la construction géométrique : le rayon parallèle à l'axe émerge en passant par $F'$, le rayon passant par $O$ n'est pas dévié ; comme ces deux rayons émergents divergent, ce sont leurs \emph{prolongements} (tracés en pointillés) qui se coupent du côté de l'objet, en $B'$.
\end{propriete}

\begin{exemplebox}[Calcul avec une loupe]
Une loupe de vergence $C = +25\ \delta$ ($f' = 4{,}0\ \text{cm}$) est placée à $3{,}0\ \text{cm}$ d'un timbre. Donc $\overline{OA} = -3{,}0\ \text{cm} = -0{,}030\ \text{m}$. La formule de position donne :
\[ \frac{1}{\overline{OA'}} = \frac{1}{\overline{OA}} + C = \frac{1}{-0{,}030} + 25 = -33{,}3 + 25 = -8{,}3\ \text{m}^{-1} \]
d'où $\overline{OA'} = -0{,}12\ \text{m} = -12\ \text{cm}$ : image virtuelle (signe négatif), située du même côté que l'objet. Le grandissement vaut $\gamma = \dfrac{\overline{OA'}}{\overline{OA}} = \dfrac{-12}{-3{,}0} = +4$ : l'image est droite et quatre fois plus grande.
\end{exemplebox}

\begin{attention}[Piège classique]
Si l'objet est placé \emph{au-delà} du foyer objet ($\overline{OA} < -f'$), la lentille donne une image \emph{réelle et renversée} : elle ne joue plus son rôle de loupe. C'est pourquoi, en éloignant progressivement une loupe d'un texte, on voit d'abord le texte agrandi et droit, puis flou (objet au voisinage de $F$, image rejetée à l'infini), puis renversé.
\end{attention}

\courssub{L'œil humain : un instrument d'optique vivant}

\textbf{Observation.} Tu lis ces lignes à $30\ \text{cm}$, puis tu regardes le tableau au fond de la salle, puis le sommet du mont Cameroun par la fenêtre : dans les trois cas, tu vois \emph{net}, alors que la distance œil-objet varie énormément. L'œil est un système optique à \emph{distance focale variable}.

\begin{definition}[Modèle de l'œil réduit]
L'œil est modélisé par l'\cle{œil réduit} : le \cle{cristallin} (associé à la cornée et aux humeurs) est assimilé à une lentille mince convergente de vergence \emph{variable} $C$, et la \cle{rétine} joue le rôle d'écran fixe, située à une distance $d \approx 17\ \text{mm}$ (soit $1{,}7\ \text{cm}$) derrière le centre optique. L'image d'un objet regardé doit se former \emph{exactement sur la rétine} : elle y est \cle{réelle} et \cle{renversée} (c'est le cerveau qui la « remet à l'endroit »).
\end{definition}

L'\cle{accommodation} est la modification de la courbure du cristallin (donc de sa vergence) par les muscles ciliaires :
\begin{itemize}
\item objet \emph{lointain} : cristallin au repos, peu bombé, vergence \emph{minimale} $C_{\min}$ ;
\item objet \emph{proche} : cristallin bombé, vergence \emph{maximale} $C_{\max}$.
\end{itemize}
L'amplitude d'accommodation est $A = C_{\max} - C_{\min}$ (en dioptries) ; elle diminue avec l'âge (presbytie).

\begin{exemplebox}[L'œil emmétrope]
Un œil \cle{emmétrope} (normal) voit net un objet à l'infini \emph{sans accommoder} : l'image se forme exactement sur la rétine, donc $f'_{\max} = d$. Avec $d = 1{,}7\ \text{cm}$ : $C_{\min} = \dfrac{1}{0{,}017} \approx 59\ \delta$. Pour lire à $25\ \text{cm}$ (punctum proximum), la formule de position $\dfrac{1}{\overline{OA'}} = \dfrac{1}{\overline{OA}} + C$ avec $\overline{OA'} = +0{,}017\ \text{m}$ et $\overline{OA} = -0{,}25\ \text{m}$ donne :
\[ C_{\max} = \frac{1}{0{,}017} - \frac{1}{-0{,}25} = 58{,}8 + 4{,}0 \approx 63\ \delta. \]
L'amplitude d'accommodation vaut donc environ $A = 63 - 59 = 4\ \delta$.
\end{exemplebox}

\begin{popfigure}
\begin{center}
\begin{tikzpicture}[scale=0.9, >=Stealth]
% Axe optique
\draw[->, popline] (-5.5,0) -- (4.5,0) node[below left] {\small axe optique};
% Cristallin (lentille convergente)
\draw[popblue, very thick] (2,-1.6) -- (2,1.6);
\draw[popblue, very thick, ->] (2,1.6) -- (1.8,1.35);
\draw[popblue, very thick, ->] (2,1.6) -- (2.2,1.35);
\draw[popblue, very thick, ->] (2,-1.6) -- (1.8,-1.35);
\draw[popblue, very thick, ->] (2,-1.6) -- (2.2,-1.35);
\node[popblue, below] at (2,-1.8) {\small cristallin $O$};
% Rétine (écran)
\draw[popdark, very thick] (3.7,-1.4) -- (3.7,1.4);
\draw[pattern=north east lines, pattern color=popdark!50] (3.7,-1.4) rectangle (3.95,1.4);
\node[popdark, below] at (3.8,-1.6) {\small rétine};
% Objet
\draw[->, poporange, very thick] (-4.5,0) -- (-4.5,1.8) node[above] {\small $B$};
\node[poporange, below] at (-4.5,0) {\small $A$};
% Rayons
\draw[popgreen, thick] (-4.5,1.8) -- (2,1.8);
\draw[popgreen, thick, ->] (2,1.8) -- (3.7,-0.6);
\draw[poppurple, thick, ->] (-4.5,1.8) -- (3.7,-0.6);
% Image renversée
\draw[->, poporange, very thick] (3.7,0) -- (3.7,-0.6) node[right] {\small $B'$};
\node[poporange, left] at (3.55,-0.1) {\small $A'$};
\node[popmuted] at (0.2,2.2) {\small rayon parallèle puis par $F'$};
\node[popmuted] at (-1.5,-1.1) {\small rayon non dévié par $O$};
\end{tikzpicture}
\end{center}
\legende{Modèle de l'œil réduit : le cristallin (lentille convergente) forme sur la rétine une image réelle, renversée et réduite de l'objet regardé.}
\end{popfigure}

\courssub{Défauts de l'œil et leur correction}

Un œil est dit \emph{amétrope} lorsque, au repos, l'image d'un objet lointain ne se forme pas sur la rétine. La correction consiste à placer devant l'œil une lentille (verre correcteur) qui, associée au cristallin par le théorème des vergences, ramène l'image sur la rétine.

\begin{definition}[Myopie et hypermétropie]
\begin{itemize}
\item L'œil \cle{myope} est \emph{trop convergent} (globe oculaire trop long ou cristallin trop bombé) : l'image d'un objet lointain se forme \emph{en avant} de la rétine. Le myope voit flou de loin, net de près. Correction : lentille \cle{divergente} ($C < 0$) qui diminue la vergence totale.
\item L'œil \cle{hypermétrope} est \emph{pas assez convergent} : l'image d'un objet lointain se formerait \emph{derrière} la rétine. L'hypermétrope accommode en permanence et se fatigue. Correction : lentille \cle{convergente} ($C > 0$) qui augmente la vergence totale.
\end{itemize}
\end{definition}

\begin{attention}[Ne pas inverser les corrections]
C'est LA confusion classique au Probatoire. Retiens la logique : on corrige toujours par le défaut \emph{opposé}.
\begin{center}
\begin{tabularx}{\linewidth}{|l|X|X|}\hline
\rowcolor{popblueL} \textbf{Défaut} & \textbf{Image d'un objet lointain} & \textbf{Correction} \\ \hline
Myopie & en avant de la rétine (œil trop convergent) & lentille \textbf{divergente}, $C<0$ \\ \hline
Hypermétropie & derrière la rétine (œil pas assez convergent) & lentille \textbf{convergente}, $C>0$ \\ \hline
\end{tabularx}
\end{center}
Moyen mnémotechnique : « \emph{myope} $\rightarrow$ \emph{moins} » : on enlève de la vergence.
\end{attention}

\begin{popfigure}
\begin{center}
\begin{tikzpicture}[scale=0.85, >=Stealth]
% ---- Oeil myope ----
\draw[->, popline] (-5.2,0) -- (1.2,0);
\draw[popblue, very thick] (-0.6,-1.1) -- (-0.6,1.1);
\draw[popdark, very thick] (0.8,-1.1) -- (0.8,1.1);
\node[popblue, below] at (-0.6,-1.25) {\small cristallin};
\node[popdark, below] at (0.8,-1.25) {\small rétine};
\draw[popgreen, thick] (-5,0.9) -- (-0.6,0.9) -- (0.1,0);
\draw[popgreen, thick] (-5,0.45) -- (-0.6,0.45) -- (0.1,0);
\draw[popgreen, thick, dashed] (0.1,0) -- (0.8,0.42);
\draw[popgreen, thick, dashed] (0.1,0) -- (0.8,-0.42);
\fill[poporange] (0.1,0) circle (2pt);
\node[poporange, above] at (0.1,0.15) {\small foyer};
\node at (-2,1.6) {\small \textbf{Œil myope} : convergence en avant de la rétine};
% ---- Oeil corrigé ----
\begin{scope}[yshift=-4.2cm]
\draw[->, popline] (-5.2,0) -- (1.2,0);
% lentille divergente correctrice
\draw[poppurple, very thick] (-2.6,-1.1) -- (-2.6,1.1);
\draw[poppurple, very thick, ->] (-2.6,1.1) -- (-2.38,0.85);
\draw[poppurple, very thick, ->] (-2.6,1.1) -- (-2.82,0.85);
\draw[poppurple, very thick, ->] (-2.6,-1.1) -- (-2.38,-0.85);
\draw[poppurple, very thick, ->] (-2.6,-1.1) -- (-2.82,-0.85);
\node[poppurple, below] at (-2.6,-1.25) {\small verre divergent};
\draw[popblue, very thick] (-0.6,-1.1) -- (-0.6,1.1);
\draw[popdark, very thick] (0.8,-1.1) -- (0.8,1.1);
\node[popdark, below] at (0.8,-1.25) {\small rétine};
\draw[popgreen, thick] (-5,0.9) -- (-2.6,0.9) -- (-0.6,0.55) -- (0.8,0);
\draw[popgreen, thick] (-5,0.45) -- (-2.6,0.45) -- (-0.6,0.28) -- (0.8,0);
\fill[poporange] (0.8,0) circle (2pt);
\node at (-2,1.6) {\small \textbf{Œil corrigé} : l'image se forme sur la rétine};
\end{scope}
\end{tikzpicture}
\end{center}
\legende{Œil myope : les rayons parallèles convergent en avant de la rétine. Une lentille divergente écarte légèrement les rayons et ramène le point de convergence exactement sur la rétine.}
\end{popfigure}

\begin{savaistu}[Les campagnes de dons de lunettes au Cameroun]
Chaque année, des associations et des ONG organisent à Douala, Yaoundé et dans les zones rurales des campagnes de dépistage visuel et de distribution de lunettes. La vergence de chaque verre correcteur (par exemple $-2\ \delta$ pour un myope, $+1{,}5\ \delta$ pour un hypermétrope) est déterminée par l'ophtalmologiste puis choisie grâce aux formules que tu viens d'étudier : la formule de position et le théorème des vergences $C_{\text{total}} = C_{\text{œil}} + C_{\text{verre}}$. Quelques dioptries bien calculées suffisent à rendre la lecture possible à un élève — la physique au service de la santé publique.
\end{savaistu}

\courssub{L'appareil photographique}

\begin{definition}[Modèle de l'appareil photographique]
L'\cle{objectif} d'un appareil photographique est assimilé à une lentille mince convergente de distance focale $f'$ fixe (par exemple $f' = 50\ \text{mm}$). Le \cle{capteur} (ou la pellicule) joue le rôle d'écran. L'objet photographié est en général très éloigné devant $f'$ ($|\overline{OA}| \gg f'$) : l'image est \cle{réelle}, \cle{renversée} et \cle{réduite}, et elle se forme pratiquement dans le plan focal image.
\end{definition}

Comme $f'$ est fixe mais que la distance objet varie, la distance objectif-capteur doit varier : c'est la \cle{mise au point}, réalisée par déplacement de l'objectif (bague de mise au point, moteur autofocus). Pour un objet à l'infini, le capteur est placé en $F'$ ($\overline{OA'} = f'$) ; plus l'objet se rapproche, plus l'image recule : il faut \emph{éloigner} l'objectif du capteur.

\begin{exemplebox}[Mise au point d'un objectif de $50\ \text{mm}$]
Pour un objet à l'infini : $\overline{OA'} = f' = 50\ \text{mm}$. Pour un portrait à $\overline{OA} = -1{,}0\ \text{m}$ :
\[ \frac{1}{\overline{OA'}} = \frac{1}{-1{,}0} + \frac{1}{0{,}050} = -1 + 20 = 19\ \text{m}^{-1} \quad\Rightarrow\quad \overline{OA'} = 52{,}6\ \text{mm}. \]
L'objectif doit donc s'éloigner du capteur de $2{,}6\ \text{mm}$ par rapport à sa position « infini ». Le grandissement vaut $\gamma = \dfrac{\overline{OA'}}{\overline{OA}} = \dfrac{0{,}0526}{-1{,}0} = -0{,}053$ : un visage de $25\ \text{cm}$ donne une image renversée de $25 \times 0{,}053 \approx 1{,}3\ \text{cm}$ sur le capteur.
\end{exemplebox}

\courssub{Lunette astronomique et microscope : associations de lentilles}

Ces deux instruments, dont l'étude détaillée est au programme de Terminale, illustrent le théorème des vergences et l'enchaînement des images :
\begin{itemize}
\item la \cle{lunette astronomique} associe un \emph{objectif} convergent de grande distance focale (qui forme de l'astre lointain une image réelle réduite dans son plan focal image) et un \emph{oculaire} convergent de courte focale qui joue le rôle de \emph{loupe} pour examiner cette image intermédiaire ;
\item le \cle{microscope} associe un objectif de très courte focale (image réelle très agrandie d'un objet placé juste devant son foyer) et un oculaire-loupe. C'est l'instrument des laboratoires de biologie des lycées et des hôpitaux du Cameroun, par exemple pour le diagnostic du paludisme par goutte épaisse.
\end{itemize}
Dans les deux cas, l'image définitive est \emph{virtuelle} et observée à travers l'oculaire : l'instrument se modélise comme une chaîne « objet $\rightarrow$ image$_1$ (par l'objectif) $\rightarrow$ image$_2$ (par l'oculaire) », où l'image$_1$ sert d'objet pour l'oculaire.

\begin{methode}[Résoudre tout exercice d'application des lentilles]
\begin{enumerate}
\item \textbf{Identifier le rôle de la lentille} dans l'énoncé : objectif (image réelle sur écran/capteur/rétine), loupe ou oculaire (image virtuelle droite agrandie), cristallin (vergence variable), verre correcteur (association de vergences).
\item \textbf{Modéliser} : dessiner le schéma équivalent (lentille mince, axe, foyers, écran éventuel) et repérer les grandeurs connues en \emph{valeurs algébriques}.
\item \textbf{Choisir l'outil} : formule de position $\frac{1}{\overline{OA'}} - \frac{1}{\overline{OA}} = \frac{1}{f'}$, grandissement $\gamma = \frac{\overline{OA'}}{\overline{OA}}$, ou théorème des vergences $C = C_1 + C_2$.
\item \textbf{Calculer} en convertissant toutes les distances en mètres si la vergence est en dioptries.
\item \textbf{Interpréter physiquement} le résultat : signe de $\overline{OA'}$ (réelle si $>0$, virtuelle si $<0$), signe et valeur de $\gamma$ (droite/renversée, agrandie/réduite), cohérence avec le rôle identifié à l'étape 1.
\end{enumerate}
\end{methode}

\begin{exoresolu}[Correction d'un œil myope]
Un œil myope est modélisé par une lentille convergente (cristallin) de vergence au repos $C_{\text{œil}} = 60\ \delta$, la rétine étant à $d = 1{,}7\ \text{cm}$ du centre optique. 1) Montrer que cet œil est défectueux pour la vision de loin. 2) Déterminer la vergence $C_c$ du verre correcteur, accolé au cristallin, qui permet de voir net un objet à l'infini.
\tcblower
\textbf{1)} Pour un objet à l'infini, l'image se forme dans le plan focal image : $f'_{\text{œil}} = \dfrac{1}{C_{\text{œil}}} = \dfrac{1}{60} = 1{,}67\ \text{cm} < d = 1{,}7\ \text{cm}$. L'image se forme \emph{en avant} de la rétine : l'œil est bien myope (trop convergent).
\textbf{2)} Il faut que le système \{verre + cristallin\} ait une distance focale image égale à $d$, donc une vergence totale $C_{\text{tot}} = \dfrac{1}{d} = \dfrac{1}{0{,}017} = 58{,}8\ \delta$. Par le théorème des vergences (lentilles accolées) :
\[ C_c = C_{\text{tot}} - C_{\text{œil}} = 58{,}8 - 60 = -1{,}2\ \delta. \]
Le verre correcteur est bien \emph{divergent} ($C_c < 0$), de vergence environ $-1{,}2\ \delta$, ce qui correspond à une distance focale $f'_c = \dfrac{1}{-1{,}2} \approx -83\ \text{cm}$.
\end{exoresolu}

\courssub{Bilan : tableau récapitulatif des images}

\begin{aretenir}[Les cinq cas de la lentille convergente]
Pour une lentille \textbf{convergente}, la nature de l'image dépend de la position de l'objet réel par rapport à $F$ et au point $\Gamma$ tel que $\overline{O\Gamma} = 2\,\overline{OF}$ :
\begin{center}
\begin{tabularx}{\linewidth}{|l|X|X|X|}\hline
\rowcolor{popblueL} \textbf{Position de l'objet} & \textbf{Nature} & \textbf{Sens} & \textbf{Taille} \\ \hline
Objet à l'infini & réelle (en $F'$) & renversée & très réduite (point) \\ \hline
$-\infty < \overline{OA} < 2\,\overline{OF}$ & réelle & renversée & réduite \\ \hline
$\overline{OA} = 2\,\overline{OF}$ & réelle (en $\Gamma'$) & renversée & égale ($\gamma=-1$) \\ \hline
$2\,\overline{OF} < \overline{OA} < \overline{OF}$ & réelle & renversée & agrandie \\ \hline
$\overline{OA} = \overline{OF}$ & image à l'infini & --- & --- \\ \hline
$\overline{OF} < \overline{OA} < 0$ (loupe) & \textbf{virtuelle} & \textbf{droite} & \textbf{agrandie} \\ \hline
\end{tabularx}
\end{center}
Pour une lentille \textbf{divergente}, un objet réel donne \emph{toujours} une image virtuelle, droite et réduite : c'est le cas du verre correcteur du myope et du judas optique des portes.
\end{aretenir}

\begin{popfigure}
\begin{center}
\begin{tikzpicture}[scale=0.78, >=Stealth, every node/.style={font=\small}]
% Ligne 1 : objet au-delà de 2F -> image réelle réduite
\begin{scope}
\draw[popline] (-4.6,0) -- (4.6,0);
\draw[popblue, very thick] (0,-1.2) -- (0,1.2);
\fill (1.5,0) circle (1.5pt) node[below] {$F'$}; \fill (-1.5,0) circle (1.5pt) node[below] {$F$};
\fill (3,0) circle (1pt) node[below] {$\Gamma'$}; \fill (-3,0) circle (1pt) node[below] {$\Gamma$};
\draw[->, poporange, very thick] (-3.8,0) -- (-3.8,1.1);
\draw[popgreen] (-3.8,1.1) -- (0,1.1) -- (2.48,-0.72);
\draw[poppurple] (-3.8,1.1) -- (2.48,-0.72);
\draw[->, poporange, very thick] (2.48,0) -- (2.48,-0.72);
\node[popmuted] at (0,1.7) {objet au-delà de $\Gamma$ : réelle, renversée, réduite};
\end{scope}
% Ligne 2 : objet entre 2F et F -> image réelle agrandie
\begin{scope}[yshift=-3.6cm]
\draw[popline] (-4.6,0) -- (5.2,0);
\draw[popblue, very thick] (0,-1.2) -- (0,1.2);
\fill (1.5,0) circle (1.5pt) node[below] {$F'$}; \fill (-1.5,0) circle (1.5pt) node[below] {$F$};
\draw[->, poporange, very thick] (-2.2,0) -- (-2.2,0.9);
\draw[popgreen] (-2.2,0.9) -- (0,0.9) -- (4.71,-1.93);
\draw[poppurple] (-2.2,0.9) -- (4.71,-1.93);
\draw[->, poporange, very thick] (4.71,0) -- (4.71,-1.93);
\node[popmuted] at (0,1.7) {objet entre $\Gamma$ et $F$ : réelle, renversée, agrandie};
\end{scope}
% Ligne 3 : objet entre F et O -> image virtuelle droite agrandie (loupe)
\begin{scope}[yshift=-7.6cm]
\draw[popline] (-4.6,0) -- (4.6,0);
\draw[popblue, very thick] (0,-1.2) -- (0,1.2);
\fill (1.5,0) circle (1.5pt) node[below] {$F'$}; \fill (-1.5,0) circle (1.5pt) node[below] {$F$};
\draw[->, poporange, very thick] (-0.8,0) -- (-0.8,0.7);
\draw[popgreen] (-0.8,0.7) -- (0,0.7) -- (2.2,-0.33);
\draw[popgreen, dashed] (0,0.7) -- (-1.71,1.5);
\draw[poppurple] (-0.8,0.7) -- (1.6,-1.4);
\draw[poppurple, dashed] (-0.8,0.7) -- (-1.71,1.5);
\draw[->, poporange, very thick, dashed] (-1.71,0) -- (-1.71,1.5);
\node[popmuted] at (0.4,2.0) {objet entre $F$ et $O$ (loupe) : virtuelle, droite, agrandie};
\end{scope}
\end{tikzpicture}
\end{center}
\legende{Les trois régimes essentiels de la lentille convergente : image réelle réduite, image réelle agrandie, et fonctionnement en loupe (image virtuelle droite agrandie, obtenue avec les prolongements en pointillés).}
\end{popfigure}

\begin{exercice}[Pour t'entraîner]
1) Un professeur hypermétrope a un cristallin de vergence au repos $C_{\text{œil}} = 57\ \delta$ ; sa rétine est à $1{,}7\ \text{cm}$ du centre optique. Quelle doit être la vergence du verre correcteur pour voir net à l'infini ? Conclure sur la nature du verre.
2) Une loupe de distance focale $f' = 5{,}0\ \text{cm}$ est placée à $4{,}0\ \text{cm}$ d'un objet de $2\ \text{mm}$. Déterminer la position, la nature et la taille de l'image.
3) Un appareil photo ($f' = 35\ \text{mm}$) photographie un objet à l'infini puis à $70\ \text{cm}$. De quelle distance faut-il déplacer l'objectif pour faire la mise au point ?
\end{exercice}

\begin{corrige}[Exercice : applications des lentilles]
\textbf{1)} Il faut $C_{\text{tot}} = \dfrac{1}{d} = \dfrac{1}{0{,}017} = 58{,}8\ \delta$. Donc $C_c = 58{,}8 - 57 = +1{,}8\ \delta > 0$ : verre \emph{convergent}, ce qui confirme l'hypermétropie (œil pas assez convergent).
\textbf{2)} $\dfrac{1}{\overline{OA'}} = \dfrac{1}{-0{,}040} + \dfrac{1}{0{,}050} = -25 + 20 = -5\ \text{m}^{-1}$, donc $\overline{OA'} = -20\ \text{cm}$ : image \emph{virtuelle}. $\gamma = \dfrac{-20}{-4{,}0} = +5$ : image droite, de taille $5 \times 2 = 10\ \text{mm}$.
\textbf{3)} À l'infini : $\overline{OA'}_1 = 35\ \text{mm}$. À $70\ \text{cm}$ : $\dfrac{1}{\overline{OA'}_2} = \dfrac{1}{-0{,}70} + \dfrac{1}{0{,}035} = -1{,}43 + 28{,}6 = 27{,}1\ \text{m}^{-1}$, soit $\overline{OA'}_2 = 36{,}8\ \text{mm}$. Déplacement : $36{,}8 - 35 = 1{,}8\ \text{mm}$ vers l'avant (on éloigne l'objectif du capteur).
\end{corrige}

Tu disposes désormais de tous les outils du chapitre : schématisation, construction géométrique, formules de position et de grandissement, théorème des vergences. Les applications étudiées ici — loupe, œil, appareil photo — sont les situations d'examen les plus fréquentes au Probatoire : entraîne-toi à toujours commencer par identifier le rôle de la lentille, et le reste suivra naturellement.

\newpage

\coursec{Activité d'intégration}

\courssub{Objectifs de l'activité}

À l'issue de cette activité d'intégration, tu seras capable de :

\begin{itemize}
\item modéliser une situation concrète de projection par le schéma objet -- lentille -- écran ;
\item utiliser la \cle{formule de position} (relation de conjugaison) et la formule du \cle{grandissement} pour déterminer la distance focale d'une lentille et la position d'un objet ;
\item vérifier un résultat de calcul par une \cle{construction géométrique} à l'échelle ;
\item choisir une lentille dans un catalogue à partir de sa \cle{vergence} ;
\item discuter l'effet d'une modification du dispositif (recul de l'écran) et proposer une nouvelle mise au point ;
\item rédiger une fiche technique claire, comme le ferait un technicien ou un ingénieur.
\end{itemize}

\courssub{Situation}

Le lycée de Ngaoundéré vient de recevoir, grâce à un projet d'équipement scientifique, un lot de matériel d'optique. Monsieur Hamadou, enseignant de SVT, souhaite monter un \textbf{projecteur de transparents} artisanal pour sa salle de classe : il veut projeter sur un mur blanc l'image de schémas (cellule, appareil digestif) dessinés sur des transparents de hauteur $AB = \SI{5,0}{cm}$.

La salle impose une contrainte : le mur-écran se trouve à une distance fixe $D = \SI{3,00}{m}$ de la lentille de projection. Pour que toute la classe voie les détails, l'image doit être agrandie \textbf{20 fois} (elle sera renversée, ce qui n'est pas gênant : il suffira de retourner le transparent).

Le club scientifique du lycée dispose d'un catalogue de lentilles convergentes de vergences : $C_1 = 2\,\delta$ ; $C_2 = 5\,\delta$ ; $C_3 = 8\,\delta$ ; $C_4 = 10\,\delta$.

\textbf{Mission du groupe d'élèves de Première C :} déterminer la lentille à utiliser, la position exacte du transparent, vérifier par une construction géométrique, étudier le cas où le mur serait reculé à $\SI{4,00}{m}$, puis rédiger la fiche technique du projecteur.

\courssub{Consignes}

\begin{enumerate}
\item Faire un schéma du dispositif : transparent $AB$ (objet réel), lentille convergente de centre optique $O$, écran à la distance $\overline{OA'} = +\SI{3,00}{m}$. Préciser le sens positif choisi et le signe du grandissement attendu.
\item À l'aide de la formule du grandissement $\gamma = \dfrac{\overline{A'B'}}{\overline{AB}} = \dfrac{\overline{OA'}}{\overline{OA}}$, déterminer la position algébrique $\overline{OA}$ du transparent.
\item En déduire, par la formule de position $\dfrac{1}{\overline{OA'}} - \dfrac{1}{\overline{OA}} = \dfrac{1}{f'}$, la distance focale $f'$ de la lentille nécessaire, puis sa vergence $C$.
\item Vérifier ce résultat par une construction géométrique soignée à l'échelle $1/10$ ($\SI{1}{cm}$ sur le papier pour $\SI{10}{cm}$ réels) : tracer les trois rayons particuliers issus de $B$ et mesurer $f'$ sur le schéma.
\item Choisir dans le catalogue la lentille la mieux adaptée. Justifier le choix par la vergence, puis recalculer le grandissement réellement obtenu avec cette lentille, l'écran restant à $\SI{3,00}{m}$. Commenter.
\item Le proviseur fait déplacer le mur-écran à $\SI{4,00}{m}$ de la lentille. En conservant la lentille de la question 3, déterminer la nouvelle position du transparent et le nouveau grandissement. Que conclure sur la « mise au point » ?
\item Rédiger la fiche technique du projecteur (lentille retenue, distances, grandissement, mode d'emploi).
\end{enumerate}

\courssub{Ressources à mobiliser}

\begin{aretenir}[Formulaire utile]
Pour une lentille mince utilisée dans les conditions de Gauss, avec la convention des grandeurs algébriques (sens positif = sens de propagation de la lumière) :
\begin{itemize}
\item \textbf{Formule de position (conjugaison)} : $\dfrac{1}{\overline{OA'}} - \dfrac{1}{\overline{OA}} = \dfrac{1}{f'} = C$ ;
\item \textbf{Grandissement} : $\gamma = \dfrac{\overline{A'B'}}{\overline{AB}} = \dfrac{\overline{OA'}}{\overline{OA}}$ ;
\item \textbf{Vergence} : $C = \dfrac{1}{f'}$, avec $f'$ en mètres et $C$ en dioptries ($\delta$) ;
\item image réelle (projetable sur un écran) : $\overline{OA'} > 0$, obtenue avec un objet réel placé au-delà du foyer objet ($\overline{OA} < -f'$) ; elle est alors \textbf{renversée} ($\gamma < 0$) ;
\item construction géométrique : rayon par le centre optique non dévié ; rayon parallèle à l'axe émergeant par $F'$ ; rayon passant par $F$ émergeant parallèle à l'axe.
\end{itemize}
\end{aretenir}

\courssub{Démarche et corrigé commenté}

\begin{exoresolu}[Fiche technique du projecteur de transparents]
\textbf{Données :} $\overline{AB} = +\SI{5,0}{cm}$ ; écran à $\overline{OA'} = +\SI{3,00}{m} = \SI{300}{cm}$ ; grandissement souhaité en valeur absolue $|\gamma| = 20$ ; catalogue : $2\,\delta$, $5\,\delta$, $8\,\delta$, $10\,\delta$.
\tcblower
\textbf{1. Schéma et signe du grandissement.} L'objet est réel (devant la lentille, $\overline{OA}<0$) et l'image doit se former \emph{sur} l'écran : c'est une image réelle, donc renversée. On attend $\gamma = -20$.

\textbf{2. Position du transparent.} De $\gamma = \dfrac{\overline{OA'}}{\overline{OA}} = -20$, on tire :
\[ \overline{OA} = \frac{\overline{OA'}}{\gamma} = \frac{+300}{-20} = \SI{-15,0}{cm}. \]
Le transparent doit être placé à $\SI{15,0}{cm}$ devant la lentille (et retourné, puisque l'image est inversée).

\textbf{3. Distance focale et vergence.} La formule de position donne, avec les distances en centimètres :
\[ \frac{1}{f'} = \frac{1}{\overline{OA'}} - \frac{1}{\overline{OA}} = \frac{1}{300} - \frac{1}{-15} = \frac{1}{300} + \frac{20}{300} = \frac{21}{300}, \]
d'où :
\[ f' = \frac{300}{21} \approx \SI{14,3}{cm} = \SI{0,143}{m} \qquad\text{et}\qquad C = \frac{1}{f'} = \frac{1}{\num{0,143}} \approx 7,0\,\delta. \]
Vérification dimensionnelle : $C$ est l'inverse d'une longueur en mètres, donc s'exprime bien en dioptries.

\textbf{4. Construction géométrique à l'échelle 1/10.} Sur la feuille : $OA = \SI{1,5}{cm}$, $OA' = \SI{30}{cm}$, $AB = \SI{0,5}{cm}$, $A'B' = \SI{10}{cm}$ (image renversée). Depuis $B$, on trace : le rayon passant par $O$ (non dévié), le rayon parallèle à l'axe qui émerge en passant par $F'$, et le rayon passant par $F$ qui émerge parallèle à l'axe. Les trois rayons convergent en $B'$, à $\SI{10}{cm}$ sous l'axe, sur l'écran. On mesure $OF' \approx \SI{1,4}{cm}$ sur le papier, soit $f' \approx \SI{14}{cm}$ réels : le calcul est confirmé.

\textbf{5. Choix de la lentille.} La vergence requise est $C \approx 7\,\delta$. Le catalogue propose $2\,\delta$ ($f' = \SI{50}{cm}$), $5\,\delta$ ($f' = \SI{20}{cm}$), $8\,\delta$ ($f' = \SI{12,5}{cm}$), $10\,\delta$ ($f' = \SI{10}{cm}$). La lentille la plus proche est $C_3 = 8\,\delta$. Avec elle, l'écran restant à $\overline{OA'} = +\SI{300}{cm}$ :
\[ \frac{1}{\overline{OA}} = \frac{1}{\overline{OA'}} - C = \frac{1}{3} - 8 = -\frac{23}{3} \;\Rightarrow\; \overline{OA} = -\frac{3}{23} \approx \SI{-0,130}{m} = \SI{-13,0}{cm}. \]
Le grandissement réel vaut $\gamma = \dfrac{300}{-13,0} \approx -23$ : l'image est agrandie 23 fois au lieu de 20, donc légèrement plus grande que prévu (hauteur $\approx \SI{115}{cm}$). C'est acceptable ; on peut aussi rapprocher très légèrement l'écran pour retrouver exactement $\gamma = -20$.

\textbf{6. Mur reculé à \SI{4,00}{m}.} Avec la lentille idéale $f' = \SI{14,3}{cm}$ ($C = 7\,\delta$) et $\overline{OA'} = +\SI{400}{cm}$ :
\[ \frac{1}{\overline{OA}} = \frac{1}{4} - 7 = -\frac{27}{4} \;\Rightarrow\; \overline{OA} = -\frac{4}{27} \approx \SI{-0,148}{m} = \SI{-14,8}{cm}, \]
et $\gamma = \dfrac{400}{-14,8} \approx -27$. Conclusion : quand l'écran recule, le transparent doit se rapprocher légèrement du foyer objet (de $15,0$ à $\SI{14,8}{cm}$) et l'image devient plus grande (27 fois) mais moins lumineuse (la lumière se répartit sur une surface plus grande). C'est exactement le réglage de \textbf{mise au point} d'un projecteur réel : on déplace l'objet (ou la lentille) quand on change la distance de projection.

\textbf{7. Fiche technique (extrait).} Lentille convergente $C = 8\,\delta$ ($f' = \SI{12,5}{cm}$) ; transparent à $\SI{13,0}{cm}$ de la lentille, retourné ; écran à $\SI{3,00}{m}$ ; image renversée, agrandie environ 23 fois ; mise au point par déplacement du transparent.
\end{exoresolu}

\courssub{Bilan de l'activité}

Cette mission a permis de mettre en évidence plusieurs idées essentielles de la leçon :

\begin{itemize}
\item une image \textbf{réelle et projetable} exige une lentille convergente et un objet placé au-delà du foyer objet ; elle est toujours \textbf{renversée} ;
\item les formules de position et de grandissement sont des outils de \textbf{prévision} : deux données (distance de projection, grandissement) suffisent pour déterminer entièrement le dispositif ($\overline{OA}$ et $f'$) ;
\item la \textbf{vergence} est la grandeur qui caractérise une lentille en pratique : c'est elle qui figure dans les catalogues, et le choix d'une lentille réelle impose souvent un compromis que l'on évalue par le calcul ;
\item la construction géométrique et le calcul se \textbf{confirment mutuellement} : c'est la démarche du physicien, qui croise modélisation mathématique et schéma expérimental ;
\item toute modification du dispositif (recul de l'écran) exige une nouvelle \textbf{mise au point} : objet et image sont liés par la relation de conjugaison, on ne peut pas déplacer l'un sans déplacer l'autre.
\end{itemize}

Tu as ainsi mobilisé, dans une situation authentique proche du métier de technicien, l'ensemble des savoir-faire de la leçon : schématiser, construire, calculer, choisir et argumenter.

\newpage

\coursec{Méthodes et exercices résolus}

\coursec{Les lentilles minces}

\courssub{Définition, classification et éléments caractéristiques}

\begin{definition}[Lentille mince]
Une \cle{lentille mince} est un système optique transparent, limité par deux surfaces sphériques (ou une plane et une sphérique), dont l'épaisseur au centre est négligeable devant les rayons de courbure de ses faces. On assimile la lentille à son plan principal (plan de la lentille) et on définit un \cle{centre optique} $O$ tel que tout rayon lumineux le traversant n'est pas dévié (sauf translation négligeable).
\end{definition}

\begin{propriete}[Classification des lentilles]
On distingue deux familles selon la forme des bords (épaisseur au centre $e_c$ vs épaisseur au bord $e_b$) :
\begin{itemize}
\item \textbf{Lentilles convergentes (bords minces) :} $e_c > e_b$. Elles font converger un faisceau incident parallèle à l'axe optique vers un foyer réel $F'$. Symboles : biconvexe, plan-convexe, ménisque convergent. Représentation schématique : trait vertical avec flèches vers l'extérieur ($\uparrow\!\downarrow$).
\item \textbf{Lentilles divergentes (bords épais) :} $e_c < e_b$. Elles font diverger un faisceau incident parallèle comme s'il provenait d'un foyer virtuel $F'$. Symboles : biconcave, plan-concave, ménisque divergent. Représentation : trait vertical avec flèches vers l'intérieur ($\downarrow\!\uparrow$).
\end{itemize}
\end{propriete}

\begin{methode}[Reconnaître et classifier une lentille par son aspect]
Pour identifier une lentille inconnue et déterminer sa nature :
\begin{enumerate}
\item \textbf{Palper la lentille} (ou observer son profil) : comparer l'épaisseur au centre $e_c$ et l'épaisseur au bord $e_b$.
\item \textbf{Conclure sur la forme} : si $e_c > e_b$, la lentille est à \cle{bords minces} ; si $e_c < e_b$, elle est à \cle{bords épais}.
\item \textbf{Conclure sur la nature optique} : une lentille à bords minces est \cle{convergente} ; une lentille à bords épais est \cle{divergente}.
\item \textbf{Vérifier par un test optique} : placer la lentille près d'un texte. Si les lettres paraissent grossies, la lentille est convergente ; si elles paraissent rapetissées, elle est divergente.
\item \textbf{Nommer la forme exacte} : biconvexe, plan-convexe ou ménisque convergent (convergentes) ; biconcave, plan-concave ou ménisque divergent (divergentes).
\end{enumerate}
\end{methode}

\begin{popfigure}
\begin{tikzpicture}[scale=0.9, every node/.style={font=\small}]
% Convergentes
\node[align=center] at (0,2.5) {\textbf{Lentilles convergentes (bords minces)}};
% Biconvexe
\draw[thick, popblue] (-2.5,1.5) arc (90:-90:0.5 and 1) -- (-1.5,-1.5) arc (-90:90:0.5 and 1) -- cycle;
\node at (-2, -2) {Biconvexe};
% Plan-convexe
\draw[thick, popblue] (0.5,1.5) -- (1.5,1.5) arc (90:-90:0.5 and 1) -- (0.5,-1.5) -- cycle;
\node at (1, -2) {Plan-convexe};
% Ménisque convergent
\draw[thick, popblue] (3.5,1.5) arc (90:-90:0.5 and 1) -- (4.5,-1.5) arc (-90:90:0.3 and 0.8) -- cycle;
\node at (4, -2) {Ménisque convergent};

% Divergentes
\node[align=center] at (0,-0.5) {\textbf{Lentilles divergentes (bords épais)}};
% Biconcave
\draw[thick, poporange] (-2.5,-3.5) arc (90:270:0.5 and 1) -- (-1.5,-5.5) arc (270:90:0.5 and 1) -- cycle;
\node at (-2, -6) {Biconcave};
% Plan-concave
\draw[thick, poporange] (0.5,-3.5) -- (1.5,-3.5) arc (90:270:0.5 and 1) -- (0.5,-5.5) -- cycle;
\node at (1, -6) {Plan-concave};
% Ménisque divergent
\draw[thick, poporange] (3.5,-3.5) arc (90:270:0.5 and 1) -- (4.5,-5.5) arc (270:90:0.3 and 0.8) -- cycle;
\node at (4, -6) {Ménisque divergent};
\end{tikzpicture}
\legende{Formes courantes de lentilles minces. Les flèches sur les schémas de cours indiquent l'épaississement (convergente) ou l'amincissement (divergente) vers les bords.}
\end{popfigure}

\begin{exoresolu}[Classification de trois lentilles]
On dispose de trois lentilles $L_1$, $L_2$, $L_3$ dont les épaisseurs mesurées sont :
\begin{center}
\begin{tabular}{|l|c|c|c|}
\hline
\rowcolor{popblueL} Lentille & $e_c$ (centre) & $e_b$ (bord) & Aspect observé sur un texte \\
\hline
$L_1$ & \SI{4,0}{mm} & \SI{1,5}{mm} & lettres grossies \\
\hline
$L_2$ & \SI{1,0}{mm} & \SI{3,5}{mm} & lettres rapetissées \\
\hline
$L_3$ & \SI{5,0}{mm} & \SI{2,0}{mm} & lettres grossies \\
\hline
\end{tabular}
\end{center}
Classer chaque lentille (nature et forme possible) en justifiant.
\tcblower
\textbf{Solution.}
\begin{enumerate}
\item \textbf{Lentille $L_1$} : $e_c = \SI{4,0}{mm} > e_b = \SI{1,5}{mm}$, donc $L_1$ est à \textbf{bords minces} : c'est une lentille \textbf{convergente}. Le test optique (lettres grossies) confirme ce résultat. Sa forme est biconvexe, plan-convexe ou ménisque convergent.
\item \textbf{Lentille $L_2$} : $e_c = \SI{1,0}{mm} < e_b = \SI{3,5}{mm}$, donc $L_2$ est à \textbf{bords épais} : c'est une lentille \textbf{divergente}. Le test (lettres rapetissées) confirme. Sa forme est biconcave, plan-concave ou ménisque divergent.
\item \textbf{Lentille $L_3$} : $e_c = \SI{5,0}{mm} > e_b = \SI{2,0}{mm}$, donc $L_3$ est à bords minces, donc \textbf{convergente} (confirmé par le grossissement du texte).
\end{enumerate}
\textbf{Conclusion} : $L_1$ et $L_3$ sont convergentes (bords minces) ; $L_2$ est divergente (bords épais). Le test optique sur un texte est toujours cohérent avec la mesure des épaisseurs.
\end{exoresolu}

\begin{aretenir}[Éléments caractéristiques d'une lentille mince]
Pour une lentille mince de centre optique $O$ :
\begin{itemize}
\item \textbf{Foyer objet $F$} : point de l'axe optique tel qu'un rayon passant par $F$ émerge parallèle à l'axe. $\overline{OF} = -f'$.
\item \textbf{Foyer image $F'$} : point de l'axe optique où convergent les rayons incidents parallèles à l'axe (convergente) ou d'où ils semblent provenir (divergente). $\overline{OF'} = f'$.
\item \textbf{Distance focale image $f'$} : distance algébrique $\overline{OF'}$. $f' > 0$ pour une convergente, $f' < 0$ pour une divergente.
\item \textbf{Vergence $C$} : $C = \dfrac{1}{f'}$ (en dioptries, $\delta$ ou $\text{m}^{-1}$). $C > 0$ convergente, $C < 0$ divergente.
\item \textbf{Plans focaux} : plans perpendiculaires à l'axe optique passant par $F$ (plan focal objet) et $F'$ (plan focal image).
\end{itemize}
\end{aretenir}

\begin{savaistu}[Optique au Cameroun : de la lunette au satellite]
La correction de la vision (myopie, hypermétropie, presbytie) repose sur le choix de lentilles ophtalmiques (divergentes pour la myopie, convergentes pour l'hypermétropie/presbytie) fabriquées ou adaptées dans les laboratoires d'optique de Yaoundé ou Douala. Les projecteurs de cinéma et de salle de classe (ENEO/SONATREL pour l'alimentation) utilisent une lentille convergente pour former une image réelle, renversée et agrandie sur l'écran. Les objectifs des caméras de surveillance (sécurité routière, banques) et des smartphones (MTN, Orange, CAMTEL) sont des systèmes de lentilles accolées corrigées (achromatiques). Le télescope spatial et les satellites d'observation (télédétection pour l'agriculture, la gestion des barrages comme Lom Pangar ou Nachtigal) emploient des optiques complexes dérivées des lentilles minces.
\end{savaistu}

\courssub{Conditions de Gauss et propriétés optiques}

\begin{definition}[Conditions de Gauss (stigmatisme approché)]
Une lentille mince donne une image nette (stigmatisme approché) d'un objet ponctuel si les rayons lumineux considérés :
\begin{enumerate}
\item sont \textbf{peu inclinés} par rapport à l'axe optique (rayons \cle{paraxiaux}) ;
\item passent \textbf{près du centre optique} $O$ (on diaphragme la lentille pour limiter son ouverture).
\end{enumerate}
Ces conditions garantissent que les approximations $\sin \alpha \approx \tan \alpha \approx \alpha$ (en radians) sont valides.
\end{definition}

\begin{propriete}[Propriétés optiques du centre optique et des foyers]
Sous conditions de Gauss :
\begin{itemize}
\item \textbf{Centre optique $O$ :} Tout rayon passant par $O$ n'est \textbf{pas dévié} (il poursuit sa route en ligne droite).
\item \textbf{Foyer image $F'$ :} Tout rayon incident \textbf{parallèle à l'axe optique} émerge en passant par $F'$ (convergente) ou semble en provenir (divergente).
\item \textbf{Foyer objet $F$ :} Tout rayon incident passant par $F$ (ou dirigé vers $F$) émerge \textbf{parallèle à l'axe optique}.
\end{itemize}
Ces trois propriétés définissent les \cle{trois rayons particuliers} pour la construction géométrique.
\end{propriete}

\begin{popfigure}
\begin{tikzpicture}[scale=1.1, >=Stealth]
% Axes
\draw[popline, thick, <->] (-5,0) -- (5,0) node[right] {Axe optique};
% Lentille convergente
\draw[thick, popblue] (0,-1.5) -- (0,1.5);
\draw[popblue, thick] (-0.1,1.5) -- (0.1,1.5) (-0.1,-1.5) -- (0.1,-1.5);
\draw[popblue, thick, <->] (-0.15,1.5) -- (-0.15,-1.5) node[midway, left] {$e_c$};
% Points
\node[below left] at (0,0) {$O$};
\node[below] at (-2,0) {$F$};
\node[below] at (2,0) {$F'$};
\draw[popblue, thick] (-2,0.1) -- (-2,-0.1) (2,0.1) -- (2,-0.1);
% Objet
\draw[popdark, thick, ->] (-3,0) -- (-3,1) node[above] {$B$};
\node[below] at (-3,0) {$A$};
% Rayon 1 : par O
\draw[popgreen, thick, ->] (-3,1) -- (0,0) -- (3.5,-1.16);
% Rayon 2 : parallèle -> F'
\draw[poporange, thick, ->] (-3,1) -- (0,1) -- (2,0) -- (4.5,-1.25);
% Rayon 3 : par F -> parallèle
\draw[poppurple, thick, ->] (-3,1) -- (-2,0) -- (4,1);
% Image
\draw[popdark, thick, ->] (3.5,-1.16) -- (3.5,0) node[below] {$A'$};
\node[below right] at (3.5,-1.16) {$B'$};
% Legendes rayons
\node[popgreen] at (-1.5,0.8) {Rayon 1};
\node[poporange] at (1,0.8) {Rayon 2};
\node[poppurple] at (-1,0.8) {Rayon 3};
\end{tikzpicture}
\legende{Les trois rayons particuliers pour une lentille convergente (objet réel $AB$ avant $F$). L'image $A'B'$ est réelle, renversée.}
\end{popfigure}

\begin{methode}[Construction géométrique de l'image $A'B'$ d'un objet $AB$]
Soit un objet $AB$ perpendiculaire à l'axe optique, $A$ sur l'axe.
\begin{enumerate}
\item \textbf{Tracer le schéma de base} : axe optique horizontal, lentille (segment fléché), centre $O$, foyers $F$ et $F'$ à l'échelle ($\overline{OF'}=f'$).
\item \textbf{Tracer deux rayons particuliers} issus de $B$ parmi :
\begin{itemize}
\item Rayon passant par $O$ : non dévié.
\item Rayon parallèle à l'axe : émerge par $F'$.
\item Rayon passant par $F$ : émerge parallèle à l'axe.
\end{itemize}
\item \textbf{Intersection} : les rayons émergents (ou leurs prolongements en \textbf{pointillés}) se coupent en $B'$. $A'$ est le pied de la perpendiculaire de $B'$ sur l'axe.
\item \textbf{Caractériser l'image} : Réelle (rayons convergents réels, $\overline{OA'}>0$) ou virtuelle (prolongements, $\overline{OA'}<0$) ; Droite ($\gamma>0$) ou renversée ($\gamma<0$) ; Agrandie ($|\gamma|>1$) ou réduite ($|\gamma|<1$).
\end{enumerate}
\textbf{Réflexe} : toujours dessiner les prolongements virtuels en pointillés et indiquer le sens de propagation par des flèches sur les rayons.
\end{methode}

\begin{exoresolu}[Construction et caractérisation d'une image]
Un objet $AB$ de hauteur \SI{2,0}{cm} est placé à \SI{6,0}{cm} devant une lentille convergente de distance focale $f' = \SI{4,0}{cm}$, $A$ sur l'axe optique.
\begin{enumerate}
\item Construire l'image $A'B'$ à l'échelle $1/2$.
\item Déterminer graphiquement la position et la taille de l'image, puis la caractériser.
\end{enumerate}
\tcblower
\textbf{Solution.}
\begin{enumerate}
\item \textbf{Construction} (échelle $1/2$ : \SI{1}{cm} papier = \SI{2}{cm} réel) : on place $O$, $F$ et $F'$ à \SI{2}{cm} de $O$ sur le papier, $A$ à \SI{3}{cm} à gauche de $O$, $B$ à \SI{1}{cm} au-dessus. De $B$, on trace le rayon par $O$ et le rayon parallèle à l'axe émergent par $F'$. Leur intersection donne $B'$.
\item \textbf{Lecture graphique} : intersection à \SI{6}{cm} de $O$ sur le papier $\Rightarrow$ \SI{12}{cm} réel. $B'$ à \SI{2}{cm} sous l'axe sur le papier $\Rightarrow$ \SI{4,0}{cm} réel.
\end{enumerate}
\textbf{Vérification par le calcul} (formule de conjugaison) :
\[ \frac{1}{\overline{OA'}} - \frac{1}{\overline{OA}} = \frac{1}{f'} \quad\Rightarrow\quad \frac{1}{\overline{OA'}} = \frac{1}{0,040} + \frac{1}{-0,060} = 25 - 16,7 = 8,33 \ \text{m}^{-1} \]
\[ \overline{OA'} = \SI{0,12}{m} = \SI{12}{cm} \qquad \gamma = \frac{\overline{OA'}}{\overline{OA}} = \frac{12}{-6} = -2 \]
donc $\overline{A'B'} = \gamma \cdot \overline{AB} = -2 \times 2,0 = \SI{-4,0}{cm}$.
\textbf{Conclusion} : l'image est à \SI{12}{cm} derrière la lentille ; \textbf{réelle} ($\overline{OA'} > 0$), \textbf{renversée} ($\gamma < 0$), \textbf{deux fois plus grande} ($|\gamma| = 2$). Graphique et calcul concordent.
\end{exoresolu}

\courssub{Formules de position et de grandissement}

\begin{propriete}[Formule de conjugaison (Descartes) et grandissement]
Dans le repère orienté (origine $O$, sens de la lumière), pour un objet $AB$ et son image $A'B'$ :
\begin{align*}
\textbf{Formule de position :} \quad & \frac{1}{\overline{OA'}} - \frac{1}{\overline{OA}} = \frac{1}{f'} = C \quad (C \text{ en } \delta = \text{m}^{-1}) \\
\textbf{Grandissement transverse :} \quad & \gamma = \frac{\overline{A'B'}}{\overline{AB}} = \frac{\overline{OA'}}{\overline{OA}}
\end{align*}
\textbf{Convention de signes :} $\overline{OA} < 0$ (objet réel devant la lentille), $\overline{OA'} > 0$ (image réelle derrière), $\overline{OA'} < 0$ (image virtuelle devant). $f'>0$ (convergente), $f'<0$ (divergente).
\end{propriete}

\begin{aretenir}[Analyse dimensionnelle et unités]
La vergence $C = 1/f'$ s'exprime en dioptries ($\delta$) si $f'$ est en \textbf{mètres} ($\text{m}^{-1}$).
\begin{center}
\begin{tabular}{|c|c|c|}
\hline
\rowcolor{popblueL} Grandeur & Symbole & Unité SI \\
\hline
Distance focale & $f'$ & mètre (m) \\
\hline
Vergence & $C$ & dioptrie ($\delta$) = $\text{m}^{-1}$ \\
\hline
Distances algébriques & $\overline{OA}, \overline{OA'}$ & mètre (m) \\
\hline
Grandissement & $\gamma$ & sans unité \\
\hline
\end{tabular}
\end{center}
\textbf{Attention :} Convertir \textbf{systématiquement} les cm en m avant tout calcul de $C$.
\end{aretenir}

\begin{methode}[Résolution de problèmes avec les formules des lentilles]
\begin{enumerate}
\item \textbf{Schématiser} et adopter la convention : axe orienté dans le sens de la lumière, origine $O$.
\item \textbf{Traduire les données} : $\overline{OA} < 0$ (objet réel) ; $f' > 0$ ou $<0$.
\item \textbf{Convertir en mètres} si $C$ intervient.
\item \textbf{Appliquer la formule} : $\dfrac{1}{\overline{OA'}} = C + \dfrac{1}{\overline{OA}}$.
\item \textbf{Calculer $\gamma$} : $\gamma = \dfrac{\overline{OA'}}{\overline{OA}}$.
\item \textbf{Conclure par une phrase complète} : position (valeur algébrique + distance), nature (réelle/virtuelle), sens (droite/renversée), taille ($|\gamma|$ ou $\overline{A'B'}$).
\end{enumerate}
\end{methode}

\begin{exoresolu}[Projection d'un diaporama]
Pour projeter l'image d'un objet sur un écran, on utilise une lentille convergente de vergence $C = +12,5\ \delta$. L'objet $AB$, de hauteur \SI{3,0}{cm}, est placé à \SI{12,0}{cm} devant la lentille.
\begin{enumerate}
\item Calculer la distance focale $f'$ de la lentille.
\item Déterminer la position de l'image $A'B'$.
\item Calculer le grandissement et la taille de l'image. Caractériser l'image.
\item À quelle distance de la lentille faut-il placer l'écran ?
\end{enumerate}
\tcblower
\textbf{Solution.}
\begin{enumerate}
\item \textbf{Distance focale} : $f' = \dfrac{1}{C} = \dfrac{1}{12,5} = \SI{0,080}{m} = \SI{8,0}{cm}$.
\item \textbf{Position de l'image} : $\overline{OA} = \SI{-0,120}{m}$.
\begin{align*}
\frac{1}{\overline{OA'}} &= C + \frac{1}{\overline{OA}} = 12,5 + \frac{1}{-0,120} = 12,5 - 8,33 = 4,17 \ \text{m}^{-1} \\
\overline{OA'} &= \frac{1}{4,17} = \SI{0,240}{m} = \SI{24,0}{cm}
\end{align*}
\item \textbf{Grandissement} : $\gamma = \dfrac{\overline{OA'}}{\overline{OA}} = \dfrac{0,240}{-0,120} = -2,0$. \\
Taille : $\overline{A'B'} = \gamma \cdot \overline{AB} = -2,0 \times 3,0 = \SI{-6,0}{cm}$ (valeur absolue \SI{6,0}{cm}).
\item \textbf{Position de l'écran} : image réelle ($\overline{OA'} > 0$), écran à \SI{24,0}{cm} derrière la lentille.
\end{enumerate}
\textbf{Conclusion} : image à \SI{24,0}{cm}, \textbf{réelle}, \textbf{renversée}, \textbf{deux fois plus grande}. Principe du projecteur de diapositives.
\end{exoresolu}

\courssub{Association de lentilles : Théorème des vergences}

\begin{propriete}[Théorème des vergences (lentilles accolées)]
Deux lentilles minces $L_1$ et $L_2$ en contact (accolées) se comportent comme une lentille unique de vergence :
\[ C = C_1 + C_2 \]
La distance focale équivalente est $f' = \dfrac{1}{C} = \dfrac{1}{C_1 + C_2} = \dfrac{f'_1 f'_2}{f'_1 + f'_2}$.
\textbf{Attention :} Ce sont les \textbf{vergences} qui s'ajoutent, jamais les distances focales ($f' \neq f'_1 + f'_2$).
\end{propriete}

\begin{propriete}[Système de deux lentilles séparées]
Soit deux lentilles $L_1$ (centre $O_1$) et $L_2$ (centre $O_2$) séparées de $e = \overline{O_1O_2} > 0$.
\begin{enumerate}
\item Déterminer l'image $A_1B_1$ de $AB$ par $L_1$ (formules ou construction).
\item $A_1B_1$ devient l'\textbf{objet} pour $L_2$. Son abscisse par rapport à $O_2$ est $\overline{O_2A_1} = \overline{O_2O_1} + \overline{O_1A_1} = -e + \overline{O_1A_1}$.
\item Déterminer l'image finale $A'B'$ par $L_2$ avec origine $O_2$.
\item Grandissement total : $\gamma = \gamma_1 \times \gamma_2$.
\end{enumerate}
\end{propriete}

\begin{methode}[Étudier un système de deux lentilles]
\textbf{(a) Lentilles accolées :}
\begin{enumerate}
\item Appliquer $C = C_1 + C_2$.
\item En déduire $f' = 1/C$ et la nature (convergente si $C>0$, divergente si $C<0$).
\end{enumerate}
\textbf{(b) Lentilles séparées (distance $e$) :}
\begin{enumerate}
\item Calculer/construire image $A_1B_1$ par $L_1$ (origine $O_1$).
\item Calculer $\overline{O_2A_1} = -e + \overline{O_1A_1}$. C'est l'objet pour $L_2$.
\item Calculer/construire image finale $A'B'$ par $L_2$ (origine $O_2$).
\item $\gamma_{\text{tot}} = \gamma_1 \times \gamma_2$.
\end{enumerate}
\end{methode}

\begin{popfigure}
\begin{tikzpicture}[scale=0.9, >=Stealth]
% Lentille 1
\draw[thick, popblue] (0,-1.2) -- (0,1.2);
\draw[popblue, thick] (-0.1,1.2) -- (0.1,1.2) (-0.1,-1.2) -- (0.1,-1.2);
\node[below] at (0,-1.4) {$O_1$};
\node[below] at (-1.5,-1.4) {$F_1$};
\node[below] at (1.5,-1.4) {$F'_1$};
\draw[popblue, thick] (-1.5,0.1) -- (-1.5,-0.1) (1.5,0.1) -- (1.5,-0.1);
% Lentille 2
\draw[thick, poporange] (4,-1.2) -- (4,1.2);
\draw[poporange, thick] (3.9,1.2) -- (4.1,1.2) (3.9,-1.2) -- (4.1,-1.2);
\node[below] at (4,-1.4) {$O_2$};
\node[below] at (2.5,-1.4) {$F_2$};
\node[below] at (5.5,-1.4) {$F'_2$};
\draw[poporange, thick] (2.5,0.1) -- (2.5,-0.1) (5.5,0.1) -- (5.5,-0.1);
% Distance e
\draw[<->, popmuted] (0,-1.8) -- (4,-1.8) node[midway, below] {$e = \overline{O_1O_2}$};
% Objet
\draw[popdark, thick, ->] (-2,0) -- (-2,0.8) node[above] {$B$};
\node[below] at (-2,0) {$A$};
% Rayons L1 (simplifiés)
\draw[popgreen, dashed, ->] (-2,0.8) -- (0,0) -- (4, -0.8);
\draw[popgreen, dashed, ->] (-2,0.8) -- (0,0.8) -- (1.5,0) -- (4, -1.6);
% Image intermediaire
\draw[popdark, thick, ->] (3, -1.2) -- (3,0) node[below] {$A_1$};
\node[below right] at (3, -1.2) {$B_1$};
% Rayons L2
\draw[poppurple, ->] (3, -1.2) -- (4, -1.2) -- (5.5,0);
\draw[poppurple, ->] (3, -1.2) -- (4, -0.8) -- (4, -0.8); % approx
% Image finale
\draw[popdark, thick, ->] (6.5, 0.5) -- (6.5,0) node[below] {$A'$};
\node[above right] at (6.5, 0.5) {$B'$};
% Labels
\node[popgreen] at (-1, 1.1) {$\gamma_1$};
\node[poppurple] at (5, -1.5) {$\gamma_2$};
\end{tikzpicture}
\legende{Schéma de principe pour deux lentilles séparées. L'image intermédiaire $A_1B_1$ (réelle ici) sert d'objet à $L_2$.}
\end{popfigure}

\begin{exoresolu}[Lentilles accolées et système de deux lentilles]
\textbf{Partie A.} On accole une lentille convergente $L_1$ ($C_1 = +8,0\ \delta$) et une lentille divergente $L_2$ ($C_2 = -3,0\ \delta$).
\begin{enumerate}
\item Calculer la vergence et la distance focale du système accolé. Quelle est sa nature ?
\end{enumerate}
\textbf{Partie B.} On sépare les deux lentilles de $e = \overline{O_1O_2} = \SI{30}{cm}$. Un objet $AB$ de \SI{2,0}{cm} est placé à \SI{25}{cm} devant $L_1$.
\begin{enumerate}
\setcounter{enumi}{1}
\item Déterminer la position et la taille de l'image intermédiaire $A_1B_1$ donnée par $L_1$.
\item Déterminer la position de l'image finale $A'B'$ donnée par $L_2$ et le grandissement total.
\end{enumerate}
\tcblower
\textbf{Solution.}
\textbf{Partie A.}
\begin{enumerate}
\item Théorème des vergences : $C = C_1 + C_2 = 8,0 - 3,0 = +5,0\ \delta$. \\
Distance focale équivalente : $f' = \dfrac{1}{C} = \dfrac{1}{5,0} = \SI{0,20}{m} = \SI{20}{cm}$. \\
Comme $C > 0$, le système accolé est \textbf{convergent}.
\end{enumerate}
\textbf{Partie B.}
\begin{enumerate}
\setcounter{enumi}{1}
\item \textbf{Image par $L_1$} : $f'_1 = \dfrac{1}{C_1} = \SI{0,125}{m} = \SI{12,5}{cm}$ ; $\overline{O_1A} = \SI{-0,25}{m}$.
\[ \frac{1}{\overline{O_1A_1}} = C_1 + \frac{1}{\overline{O_1A}} = 8,0 - \frac{1}{0,25} = 8,0 - 4,0 = 4,0 \ \text{m}^{-1} \]
\[ \overline{O_1A_1} = \SI{0,25}{m} = \SI{25}{cm} \qquad \gamma_1 = \frac{\overline{O_1A_1}}{\overline{O_1A}} = \frac{0,25}{-0,25} = -1,0 \]
L'image $A_1B_1$ est réelle, renversée, de même taille : $\overline{A_1B_1} = \SI{-2,0}{cm}$.
\item \textbf{Image par $L_2$} : $A_1$ se trouve à $\overline{O_2A_1} = \overline{O_2O_1} + \overline{O_1A_1} = -30 + 25 = \SI{-5}{cm}$ de $L_2$ : objet réel pour $L_2$. $f'_2 = \dfrac{1}{C_2} = \SI{-0,333}{m}$.
\[ \frac{1}{\overline{O_2A'}} = C_2 + \frac{1}{\overline{O_2A_1}} = -3,0 + \frac{1}{-0,05} = -3,0 - 20 = -23 \ \text{m}^{-1} \]
\[ \overline{O_2A'} = \SI{-0,0435}{m} \approx \SI{-4,3}{cm} \qquad \gamma_2 = \frac{\overline{O_2A'}}{\overline{O_2A_1}} = \frac{-4,3}{-5} = +0,87 \]
Grandissement total : $\gamma = \gamma_1 \times \gamma_2 = -1,0 \times 0,87 = -0,87$.
\end{enumerate}
\textbf{Conclusion} : accolées, les deux lentilles forment un système convergent ($C=+5,0\ \delta$, $f'=\SI{20}{cm}$). Séparées de \SI{30}{cm}, l'image finale est \textbf{virtuelle} ($\overline{O_2A'} < 0$), à \SI{4,3}{cm} devant $L_2$, \textbf{renversée} et \textbf{légèrement réduite} ($\gamma \approx -0,87$, taille $\approx \SI{1,7}{cm}$).
\end{exoresolu}

\courssub{Applications des lentilles minces}

\begin{experience}[Détermination expérimentale de la vergence $C$ (Banc d'optique)]
\textbf{Matériel :} Banc d'optique gradué, source lumineuse (objet fente ou grille), lentille mince convergente sur porte-lentille, écran sur chariot, règle.
\textbf{Protocole (Méthode de la formule de conjugaison) :}
\begin{enumerate}
\item Aligner source, lentille et écran à même hauteur (conditions de Gauss : centres alignés, axe horizontal).
\item Placer l'objet en $A_1$ (ex: $\overline{OA_1} = -30\ \text{cm}$). Déplacer l'écran jusqu'à obtenir une image nette $A'_1$. Relever $\overline{OA_1}$ et $\overline{OA'_1}$.
\item Répéter pour au moins 3 positions de l'objet (ex: $-20\ \text{cm}$, $-15\ \text{cm}$) en veillant à $|\overline{OA}| > f'$ pour avoir une image réelle.
\item Convertir distances en mètres. Calculer $C_i = \dfrac{1}{\overline{OA'_i}} - \dfrac{1}{\overline{OA_i}}$ pour chaque mesure.
\item Calculer la moyenne $\bar{C}$ et l'écart-type. $f' = 1/\bar{C}$.
\end{enumerate}
\textbf{Interprétation :} La constance de $C_i$ vérifie expérimentalement la formule de conjugaison.
\textbf{Méthode rapide (objet à l'infini) :} Viser une fenêtre/arbre lointain ; image nette sur écran $\Rightarrow f' \approx \overline{OA'}$.
\end{experience}

\begin{exoresolu}[TP : Mesure de la vergence d'une lentille]
Au laboratoire, un groupe d'élèves réalise les mesures suivantes avec une lentille convergente :
\begin{center}
\begin{tabular}{|l|c|c|c|}
\hline
\rowcolor{popblueL} Mesure & $\overline{OA}$ (cm) & $\overline{OA'}$ (cm) & $C$ calculée ($\delta$) \\
\hline
1 & $-30,0$ & $+15,0$ & \\
\hline
2 & $-20,0$ & $+20,0$ & \\
\hline
3 & $-15,0$ & $+30,0$ & \\
\hline
\end{tabular}
\end{center}
\begin{enumerate}
\item Compléter la troisième colonne.
\item En déduire la vergence moyenne et la distance focale.
\item Pourquoi une mesure avec $\overline{OA} = \SI{-10}{cm}$ n'aurait-elle donné aucune image sur l'écran ?
\end{enumerate}
\tcblower
\textbf{Solution.}
\begin{enumerate}
\item \textbf{Calculs} ($C = \dfrac{1}{\overline{OA'}} - \dfrac{1}{\overline{OA}}$, distances en m) :
\begin{itemize}
\item Mesure 1 : $C_1 = \dfrac{1}{0,150} - \dfrac{1}{-0,300} = 6,67 + 3,33 = 10,0\ \delta$.
\item Mesure 2 : $C_2 = \dfrac{1}{0,200} - \dfrac{1}{-0,200} = 5,00 + 5,00 = 10,0\ \delta$.
\item Mesure 3 : $C_3 = \dfrac{1}{0,300} - \dfrac{1}{-0,150} = 3,33 + 6,67 = 10,0\ \delta$.
\end{itemize}
\item \textbf{Vergence moyenne} : $\bar{C} = 10,0\ \delta$, d'où $f' = \dfrac{1}{\bar{C}} = \SI{0,100}{m} = \SI{10,0}{cm}$. La constance de $C$ vérifie la formule.
\item \textbf{Cas $\overline{OA} = \SI{-10}{cm}$} : l'objet est au foyer objet ($|\overline{OA}| = f'$). La formule donne $\dfrac{1}{\overline{OA'}} = 10 - \dfrac{1}{0,10} = 0 \Rightarrow \overline{OA'} \to +\infty$ : l'image est à l'infini, impossible à recevoir sur un écran. Si $|\overline{OA}| < f'$, l'image est virtuelle ($\overline{OA'} < 0$), non recevable sur écran.
\end{enumerate}
\textbf{Conclusion} : $C = 10,0\ \delta$, $f' = \SI{10,0}{cm}$. Pour une image réelle sur écran, il faut $|\overline{OA}| > f'$.
\end{exoresolu}

\begin{methode}[Conditions de Gauss et fonctionnement de la loupe]
\begin{enumerate}
\item \textbf{Rappel conditions de Gauss} : rayons paraxiaux, diaphragme pour netteté.
\item \textbf{Loupe} : lentille convergente utilisée avec l'objet \textbf{entre $F$ et $O$} ($-f' < \overline{OA} < 0$).
\item L'image est \textbf{virtuelle} ($\overline{OA'} < 0$), \textbf{droite} ($\gamma > 0$) et \textbf{agrandie} ($|\gamma| > 1$).
\item Elle ne peut être reçue sur un écran ; on l'observe directement à travers la lentille.
\item Grossissement maximal quand l'objet $\to F$ ($\gamma \to +\infty$). Grossissement commercial souvent défini pour image à $25\ \text{cm}$ (distance de vision distincte) : $G = \dfrac{25}{f'(\text{cm})} + 1$.
\end{enumerate}
\end{methode}

\begin{popfigure}
\begin{tikzpicture}[scale=1.1, >=Stealth]
% Lentille
\draw[thick, popblue] (0,-1.5) -- (0,1.5);
\draw[popblue, thick] (-0.1,1.5) -- (0.1,1.5) (-0.1,-1.5) -- (0.1,-1.5);
\node[below left] at (0,0) {$O$};
\node[below] at (-3,0) {$F$};
\node[below] at (3,0) {$F'$};
\draw[popblue, thick] (-3,0.1) -- (-3,-0.1) (3,0.1) -- (3,-0.1);
% Objet entre F et O
\draw[popdark, thick, ->] (-2,0) -- (-2,0.8) node[above] {$B$};
\node[below] at (-2,0) {$A$};
% Rayons
\draw[popgreen, thick, ->] (-2,0.8) -- (0,0) -- (4,-1.6);
\draw[poporange, thick, ->] (-2,0.8) -- (0,0.8) -- (-3,0) -- (4,2.4);
% Image virtuelle (pointillés)
\draw[popdark, thick, ->, dashed] (-6,2.4) -- (-6,0) node[below] {$A'$};
\node[above left] at (-6,2.4) {$B'$};
% Oeil
\draw[poppurple, thick] (4.5,0) circle (0.3);
\draw[poppurple, thick] (4.5,0.3) -- (4.7,0.5);
\draw[poppurple, thick] (4.5,-0.3) -- (4.7,-0.5);
\node[above] at (4.5,0.6) {Œil};
% Flèches virtuelles
\draw[popmuted, <->, dashed] (-2,1.2) -- (-6,1.2) node[midway, above] {Prolongements};
\end{tikzpicture}
\legende{Principe de la loupe : objet $AB$ entre $F$ et $O$. Les rayons émergents divergent ; leurs prolongements (pointillés) se coupent en $B'$ donnant une image virtuelle, droite, agrandie. L'œil placé derrière la lentille perçoit cette image.}
\end{popfigure}

\begin{exoresolu}[Observation d'un insecte à la loupe]
Un élève observe un insecte de \SI{4,0}{mm} avec une loupe de vergence $C = +20\ \delta$. Il place l'insecte à \SI{4,0}{cm} de la lentille.
\begin{enumerate}
\item Calculer la distance focale et vérifier les conditions d'utilisation.
\item Déterminer la position, la nature et la taille de l'image observée.
\item L'élève pourrait-il projeter cette image sur un écran ? Justifier.
\end{enumerate}
\tcblower
\textbf{Solution.}
\begin{enumerate}
\item $f' = \dfrac{1}{C} = \dfrac{1}{20} = \SI{0,050}{m} = \SI{5,0}{cm}$. Objet à \SI{4,0}{cm} : $|\overline{OA}| = 4,0 < f' = 5,0$ cm. L'objet est bien entre $F$ et $O$ : conditions remplies.
\item \textbf{Position} : $\overline{OA} = \SI{-0,040}{m}$.
\[ \frac{1}{\overline{OA'}} = C + \frac{1}{\overline{OA}} = 20 - \frac{1}{0,040} = 20 - 25 = -5,0 \ \text{m}^{-1} \]
\[ \overline{OA'} = \SI{-0,20}{m} = \SI{-20}{cm} \qquad \gamma = \frac{\overline{OA'}}{\overline{OA}} = \frac{-0,20}{-0,040} = +5,0 \]
Taille : $\overline{A'B'} = 5,0 \times 4,0 = \SI{20}{mm} = \SI{2,0}{cm}$.
\item L'image est \textbf{virtuelle} ($\overline{OA'} < 0$) : formée par prolongements, du même côté que l'objet. Aucun rayon ne passe en $A'B'$ $\Rightarrow$ \textbf{impossible} à recevoir sur un écran. On l'observe en regardant à travers la loupe.
\end{enumerate}
\textbf{Conclusion} : image virtuelle, droite, 5 fois agrandie (\SI{20}{mm}), à \SI{20}{cm} devant la loupe. Non projetable.
\end{exoresolu}

\begin{savaistu}[L'œil humain : une lentille vivante]
L'œil fonctionne comme un système optique à lentille variable (cristallin). La cornée (vergence $\approx +43\ \delta$) et le cristallin ($\approx +19\ \delta$ au repos, jusqu'à $+33\ \delta$ en accommodation) forment une image réelle, renversée, réduite sur la rétine (écran). La myopie (œil trop long ou trop convergent) se corrige par une lentille divergente ; l'hypermétropie (œil trop court) par une convergente. La presbytie (perte d'accommodation avec l'âge) nécessite une addition convergente pour la vision de près. Au Cameroun, les campagnes de dépistage visuel dans les écoles et les centres de santé (MINSANTÉ, ONG) permettent de distribuer des lunettes adaptées, améliorant la scolarité et la productivité.
\end{savaistu}

\begin{attention}[Les 5 erreurs qui coûtent des points au Probatoire]
\begin{enumerate}
\item \textbf{Oublier les signes algébriques.} Dans $\dfrac{1}{\overline{OA'}} - \dfrac{1}{\overline{OA}} = \dfrac{1}{f'}$, un objet réel a $\overline{OA} < 0$. Écrire $\dfrac{1}{OA'} + \dfrac{1}{OA} = \dfrac{1}{f'}$ avec distances positives sans justifier le changement de convention est une erreur sanctionnée.
\item \textbf{Mélanger centimètres et dioptries.} $C = 1/f'$ (en $\delta$) exige $f'$ en \textbf{mètres} : $f' = \SI{10}{cm} \Rightarrow C = 10\ \delta$, pas $0,1\ \delta$. Convertis \textbf{systématiquement} avant calcul.
\item \textbf{Additionner les distances focales au lieu des vergences.} Pour deux lentilles accolées : $C = C_1 + C_2 \Rightarrow \dfrac{1}{f'} = \dfrac{1}{f'_1} + \dfrac{1}{f'_2}$, mais $f' \neq f'_1 + f'_2$ (sauf cas particuliers).
\item \textbf{Prétendre recevoir une image virtuelle sur un écran.} Seule une image réelle ($\overline{OA'} > 0$) est recevable. Pour une convergente, si $|\overline{OA}| \leq f'$, aucune image nette sur écran.
\item \textbf{Conclure sans caractériser complètement l'image.} Une réponse complète exige : \textbf{position} (valeur algébrique + distance), \textbf{nature} (réelle/virtuelle), \textbf{sens} (droite/renversée via signe de $\gamma$), \textbf{taille} ($|\gamma|$ ou $\overline{A'B'}$). Sur schéma : flèches sur rayons, pointillés pour prolongements virtuels.
\end{enumerate}
\end{attention}

\newpage

\coursec{Exercices}

\courssub{Je vérifie mes connaissances}

\begin{exercice}[QCM : Reconnaissance et classification]
Parmi les propositions suivantes, une seule est correcte pour chaque question. Recopie le numéro de la question et la lettre de la bonne réponse.
\begin{enumerate}
    \item Une lentille mince convergente a une vergence $C = +5\,\delta$. Sa distance focale image $f'$ vaut :
    \begin{itemize}
        \item[A)] $0,20\,\text{m}$
        \item[B)] $5\,\text{m}$
        \item[C)] $0,02\,\text{m}$
        \item[D)] $0,05\,\text{m}$
    \end{itemize}
    \item Un rayon incident passant par le centre optique $O$ d'une lentille mince :
    \begin{itemize}
        \item[A)] est dévié vers le foyer image $F'$
        \item[B)] n'est pas dévié (poursuit sa route en ligne droite)
        \item[C)] ressort parallèle à l'axe optique
        \item[D)] est réfléchi symétriquement
    \end{itemize}
    \item Une lentille divergente de vergence $C = -4\,\delta$ reçoit un faisceau incident parallèle à l'axe optique. Le faisceau émergent :
    \begin{itemize}
        \item[A)] converge vers le foyer image $F'$ (réel)
        \item[B)] diverge comme s'il provenait du foyer objet $F$ (réel)
        \item[C)] diverge comme s'il provenait du foyer image $F'$ (virtuel)
        \item[D)] reste parallèle à l'axe optique
    \end{itemize}
    \item L'image d'un objet réel donnée par une lentille convergente est réelle, inversée et de même taille que l'objet. L'objet se trouve nécessairement :
    \begin{itemize}
        \item[A)] au foyer objet $F$
        \item[B)] au centre optique $O$
        \item[C)] à $2f'$ du centre optique (au centre de symétrie)
        \item[D)] entre $F$ et $O$
    \end{itemize}
\end{enumerate}
\end{exercice}

\begin{exercice}[Vrai ou Faux ? Justifie chaque réponse]
Pour chaque affirmation, indique si elle est \textbf{Vraie} ou \textbf{Fausse} et justifie brièvement ton choix par une loi, une définition ou un contre-exemple.
\begin{enumerate}
    \item La vergence d'une lentille mince s'exprime en dioptries ($\delta$) et sa distance focale en mètres (m) ; la relation $C = \dfrac{1}{f'}$ est toujours valable si $f'$ est en mètres.
    \item Pour une lentille mince, le centre optique $O$ est le milieu du segment $FF'$ reliant les deux foyers.
    \item Si un objet est placé entre le foyer objet $F$ et le centre optique $O$ d'une lentille convergente, l'image est virtuelle, droite et plus grande que l'objet (loupe).
    \item Le grandissement $g = \dfrac{A'B'}{AB}$ est toujours négatif pour une lentille convergente.
    \item Deux lentilles minces accolées de vergences $C_1 = +3\,\delta$ et $C_2 = -5\,\delta$ se comportent comme une lentille divergente de vergence $-2\,\delta$.
\end{enumerate}
\end{exercice}

\begin{exercice}[Définitions et formules essentielles]
Complète les phrases suivantes ou donne les formules demandées.
\begin{enumerate}
    \item On appelle \textbf{vergence} d'une lentille mince la grandeur algébrique notée $C$ qui caractérise son pouvoir convergent ou divergent. Son unité SI est le \ldots\ldots\ldots\ldots\ldots Elle est liée à la distance focale image $f'$ par la relation \ldots\ldots\ldots\ldots\ldots
    \item Énonce les \textbf{trois conditions de Gauss} pour une lentille mince convergente (décris le sort des rayons incidents particuliers).
    \item Donne la \textbf{formule de position} (relation de Descartes) liant la distance algébrique objet $p = \overline{OA}$, la distance algébrique image $p' = \overline{OA'}$ et la distance focale image $f' = \overline{OF'}$.
    \item Donne la formule du \textbf{grandissement transverse} $g$ en fonction de $p$ et $p'$. Que signifie le signe de $g$ sur le sens de l'image par rapport à l'objet ?
    \item Énonce le \textbf{théorème des vergences} pour deux lentilles minces accolées.
\end{enumerate}
\end{exercice}

\begin{exercice}[Calculs directs d'application]
\begin{enumerate}
    \item Une lentille convergente a une distance focale $f' = 25\,\text{cm}$. Calcule sa vergence $C$ en dioptries.
    \item Une lentille divergente de vergence $C = -2,5\,\delta$ reçoit un objet réel placé à $p = -40\,\text{cm}$ (objet à gauche de la lentille). Calcule la position de l'image $p'$ (en cm) et le grandissement $g$.
    \item Un objet $AB$ de hauteur $2\,\text{cm}$ est placé à $30\,\text{cm}$ d'une lentille convergente de focale $f' = 10\,\text{cm}$. Détermine par le calcul la position $p'$, la nature (réelle/virtuelle), le sens (droite/inversée) et la taille de l'image $A'B'$.
    \item Deux lentilles minces accolées ont des vergences $C_1 = +6\,\delta$ et $C_2 = -1,5\,\delta$. Calcule la vergence $C$ et la distance focale $f'$ du système équivalent. Précise sa nature (convergente ou divergente).
\end{enumerate}
\end{exercice}

\courssub{Je m'entraîne}

\begin{exercice}[Construction géométrique et analyse : La loupe]
Un collectionneur observe un timbre-poste $AB$ de hauteur $h = 0,5\,\text{cm}$ à l'aide d'une loupe (lentille convergente) de distance focale $f' = 5\,\text{cm}$. Il place le timbre à une distance $p = -4\,\text{cm}$ du centre optique $O$ de la lentille (objet entre $F$ et $O$).
\begin{enumerate}
    \item Sur une figure à l'échelle (par exemple $1\,\text{cm}$ pour $1\,\text{cm}$ réel), trace la lentille, ses foyers $F$ et $F'$, l'objet $AB$ et construis l'image $A'B'$ en utilisant au moins deux rayons particuliers.
    \item Par la construction, détermine graphiquement la position de l'image $p'$ (en cm) et sa hauteur $h'$ (en cm). Précise la nature et le sens de l'image.
    \item Vérifie tes résultats graphiques par le calcul (formules de Descartes et de grandissement).
    \item Calcule le grandissement $g$. L'image est-elle agrandie ? L'observateur voit-il l'image du même côté que l'objet ou de l'autre côté de la lentille ?
\end{enumerate}
\end{exercice}

\begin{popfigure}
\begin{tikzpicture}[scale=0.8, >=Stealth]
    % Axes
    \draw[popline, thick] (-6,0) -- (6,0) node[right] {Axe optique};
    % Lentille
    \draw[popblue, very thick] (0,-1.5) -- (0,1.5) node[above] {$L$};
    \draw[popblue, thick] (-0.1,-1.5) -- (0.1,-1.5) (-0.1,1.5) -- (0.1,1.5);
    % Foyers
    \foreach \x/\lab in {-5/F, 5/F'} {
        \draw[poporange, thick] (\x,-0.2) -- (\x,0.2);
        \node[poporange, below] at (\x,-0.3) {$\lab$};
    }
    % Centre optique
    \draw[popdark, thick] (0,-0.2) -- (0,0.2);
    \node[popdark, below] at (0,-0.5) {$O$};
    % Objet AB
    \draw[popgreen, very thick] (-4,0) -- (-4,0.5);
    \node[popgreen, left] at (-4,0.5) {$B$};
    \node[popgreen, left] at (-4,0) {$A$};
    % Rayon 1 : Parallèle à l'axe -> passe par F'
    \draw[popgreen, dashed] (-4,0.5) -- (0,0.5);
    \draw[popgreen, thick, ->] (0,0.5) -- (5,0);
    % Rayon 2 : Passe par O -> non dévié
    \draw[popgreen, thick, ->] (-4,0.5) -- (6, -0.75);
    % Rayon 3 : Passe par F -> ressort parallèle
    \draw[popgreen, thick, ->] (-4,0.5) -- (-5,0);
    \draw[popgreen, dashed, ->] (0,-0.5) -- (6, -0.5);
    % Image virtuelle A'B' (intersection des prolongements en pointillés)
    \draw[poppurple, very thick] (-20,0) -- (-20,2.5);
    \node[poppurple, left] at (-20,2.5) {$B'$};
    \node[poppurple, left] at (-20,0) {$A'$};
    % Flèches pour distances
    \draw[<->, popmuted] (-4,-1) -- (0,-1) node[midway, below] {$|p|=4\,\text{cm}$};
    \draw[<->, popmuted] (0,-1.5) -- (5,-1.5) node[midway, below] {$f'=5\,\text{cm}$};
    \draw[<->, popmuted] (-20,-1) -- (0,-1) node[midway, below] {$|p'|=20\,\text{cm}$};
\end{tikzpicture}
\legende{Construction de l'image par une loupe (lentille convergente, objet entre F et O). L'image $A'B'$ est virtuelle, droite et agrandie.}
\end{popfigure}

\begin{exercice}[Détermination expérimentale de la vergence : Banc d'optique]
Au laboratoire, on utilise un banc d'optique muni d'une lentille mince convergente inconnue, d'un objet lumineux (fente) et d'un écran. On déplace l'objet et l'écran pour obtenir une image nette sur l'écran. On relève les positions sur la règle graduée (en cm) :
\begin{center}
\begin{tabular}{|c|c|c|c|c|c|}\hline
Position objet $x_A$ & $30,0$ & $25,0$ & $20,0$ & $15,0$ & $10,0$ \\ \hline
Position lentille $x_O$ & $50,0$ & $50,0$ & $50,0$ & $50,0$ & $50,0$ \\ \hline
Position écran (image) $x_{A'}$ & $70,0$ & $66,7$ & $65,0$ & $64,0$ & $63,3$ \\ \hline
\end{tabular}
\end{center}
\begin{enumerate}
    \item Pour chaque mesure, calcule les distances algébriques $p = \overline{OA}$ et $p' = \overline{OA'}$ (en cm). Rappelle la convention de signes (objet réel à gauche de la lentille, lumière incidente de gauche vers droite).
    \item Calcule les inverses $\dfrac{1}{p'}$ et $\dfrac{1}{p}$ (en $\text{cm}^{-1}$).
    \item Trace sur papier millimétré (ou représente le nuage de points) le graphique de $\dfrac{1}{p'}$ en fonction de $\dfrac{1}{p}$. Quelle allure a ce graphique ? Justifie par la formule de Descartes.
    \item Déduis graphiquement la distance focale image $f'$ (en cm) et la vergence $C$ (en $\delta$) de la lentille.
    \item Quelle est l'incertitude approximative sur $f'$ si la lecture sur la règle se fait au millimètre près ?
\end{enumerate}
\end{exercice}

\begin{popfigure}
\begin{tikzpicture}
\begin{axis}[
    width=\linewidth, height=8cm,
    xlabel={$1/p$ (cm$^{-1}$)}, ylabel={$1/p'$ (cm$^{-1}$)},
    xmin=-0.11, xmax=0.01, ymin=-0.01, ymax=0.11,
    grid=major, minor tick num=4,
    axis lines=middle,
    tick label style={font=\small},
    label style={font=\small},
]
% Points théoriques pour f'=10cm -> 1/f'=0.1 cm^-1
% 1/p' = 1/f' + 1/p = 0.1 + 1/p
\addplot[popcurve, domain=-0.1:0, samples=50] {0.1 + x};
\addplot[only marks, mark=*, popblue, mark size=2] coordinates {
    (-0.05, 0.05)   % p=-20, p'=20
    (-0.0667, 0.0333) % p=-15, p'=30
    (-0.0833, 0.0167) % p=-12, p'=60
    (-0.1, 0)       % p=-10, p'=inf
};
% Points expérimentaux (avec légère dispersion)
\addplot[only marks, mark=square*, poporange, mark size=2, error bars/.cd, y dir=both, y explicit, error bar style={poporange}] coordinates {
    (-0.05, 0.050)   % p=-20, p'=20
    (-0.0667, 0.033) % p=-15, p'=30.3
    (-0.0833, 0.017) % p=-12, p'=58.8
    (-0.1, 0.001)    % p=-10, p'=1000 (approx)
};
\end{axis}
\end{tikzpicture}
\legende{Graphique $1/p' = f(1/p)$ pour une lentille convergente $f'=10\,\text{cm}$. La droite a pour équation $y = 0,1 + x$. L'ordonnée à l'origine donne $1/f'$.}
\end{popfigure}

\begin{exercice}[Association de deux lentilles accolées : Correction d'amétropie]
On dispose de deux lentilles minces accolées : $L_1$ convergente de vergence $C_1 = +8\,\delta$ et $L_2$ divergente de vergence $C_2 = -3\,\delta$.
\begin{enumerate}
    \item Calcule la vergence $C$ et la distance focale $f'$ du système équivalent $L$. Précise sa nature.
    \item Un objet réel $AB$ de taille $1\,\text{cm}$ est placé à $50\,\text{cm}$ devant le système $L$. Détermine par le calcul la position $p'$, la nature, le sens et la taille de l'image $A'B'$.
    \item On souhaite obtenir un système équivalent de vergence $C = +2\,\delta$ en accolant une troisième lentille $L_3$ aux deux premières. Quelle doit être la vergence $C_3$ de $L_3$ ? De quelle nature est $L_3$ ?
    \item \textbf{Contexte ophtalmologique :} Un œil myope a une puissance de convergence trop forte. Pour le corriger, on place devant lui un verre correcteur (lentille mince). Si l'œil a une vergence excédentaire de $3\,\delta$ par rapport à la normale, quelle vergence doit avoir le verre correcteur ? Est-ce un verre convergent ou divergent ?
\end{enumerate}
\end{exercice}

\begin{exercice}[Système de deux lentilles séparées : Projecteur]
Un système optique est constitué de deux lentilles minces convergentes séparées d'une distance $d = 60\,\text{cm}$.
La lentille $L_1$ (objectif) a une distance focale $f'_1 = 20\,\text{cm}$.
La lentille $L_2$ (oculaire/relai) a une distance focale $f'_2 = 10\,\text{cm}$.
Un objet $AB$ de hauteur $2\,\text{cm}$ est placé à $25\,\text{cm}$ devant $L_1$ ($p_1 = -25\,\text{cm}$).
\begin{enumerate}
    \item \textbf{Image par $L_1$ :} Calcule la position $p'_1$ de l'image $A_1B_1$ donnée par $L_1$. Précise sa nature, son sens et sa hauteur $h_1$. Cette image est-elle réelle ou virtuelle par rapport à $L_1$ ?
    \item \textbf{Objet pour $L_2$ :} L'image $A_1B_1$ sert d'objet pour $L_2$. Calcule la distance algébrique objet $p_2 = \overline{O_2A_1}$ pour la lentille $L_2$. (Attention au signe : $A_1$ est-il à gauche ou à droite de $O_2$ ?)
    \item \textbf{Image finale par $L_2$ :} Calcule la position $p'_2$ de l'image finale $A'B'$ donnée par $L_2$. Précise sa nature, son sens et sa hauteur $h'$ finale.
    \item Calcule le grandissement total $G = g_1 \times g_2$ du système. Vérifie que $h' = G \times h$.
    \item Trace schématiquement le système (pas nécessairement à l'échelle) en reportant les positions de $O_1, O_2, F_1, F'_1, F_2, F'_2, A, B, A_1, B_1, A', B'$.
\end{enumerate}
\end{exercice}

\begin{popfigure}
\begin{tikzpicture}[scale=0.5, >=Stealth]
    % Axes
    \draw[popline, thick] (-5,0) -- (25,0) node[right] {Axe optique};
    % Lentille L1
    \draw[popblue, very thick] (0,-2) -- (0,2) node[above] {$L_1$};
    \node[popblue, below] at (0,-2.5) {$O_1$};
    % Foyers L1
    \foreach \x/\lab in {-20/F_1, 20/F'_1} {
        \draw[popblue, thick] (\x,-0.3) -- (\x,0.3);
        \node[popblue, below] at (\x,-0.6) {$\lab$};
    }
    % Lentille L2
    \draw[poporange, very thick] (60,-2) -- (60,2) node[above] {$L_2$};
    \node[poporange, below] at (60,-2.5) {$O_2$};
    % Foyers L2
    \foreach \x/\lab in {50/F_2, 70/F'_2} {
        \draw[poporange, thick] (\x,-0.3) -- (\x,0.3);
        \node[poporange, below] at (\x,-0.6) {$\lab$};
    }
    % Objet AB
    \draw[popgreen, very thick] (-25,0) -- (-25,2);
    \node[popgreen, left] at (-25,2) {$B$};
    \node[popgreen, left] at (-25,0) {$A$};
    % Rayons L1
    \draw[popgreen, dashed] (-25,2) -- (0,2);
    \draw[popgreen, thick, ->] (0,2) -- (20,0);
    \draw[popgreen, thick, ->] (-25,2) -- (0,0) -- (25, -1.25);
    % Image intermédiaire A1B1 (réelle, inversée)
    \draw[poppurple, very thick] (100,-4) -- (100,0);
    \node[poppurple, right] at (100,-4) {$B_1$};
    \node[poppurple, right] at (100,0) {$A_1$};
    % Rayons L2 (objet virtuel pour L2 en A1B1)
    % Rayon passant par O2
    \draw[poppurple, thick, ->] (100,-4) -- (60,0) -- (50, 2.5);
    % Rayon parallèle à l'axe -> F'2
    \draw[poppurple, dashed] (100,-4) -- (60,-4);
    \draw[poppurple, thick, ->] (60,-4) -- (70,0);
    % Image finale A'B' (réelle)
    \draw[poppink, very thick] (80,0) -- (80,1.33);
    \node[poppink, right] at (80,1.33) {$B'$};
    \node[poppink, right] at (80,0) {$A'$};
    % Distances
    \draw[<->, popmuted] (-25,-3) -- (0,-3) node[midway, below] {$p_1=-25\,\text{cm}$};
    \draw[<->, popmuted] (0,-3.5) -- (100,-3.5) node[midway, below] {$p'_1=100\,\text{cm}$};
    \draw[<->, popmuted] (60,-4) -- (100,-4) node[midway, below] {$p_2=+40\,\text{cm}$};
    \draw[<->, popmuted] (60,-4.5) -- (80,-4.5) node[midway, below] {$p'_2\approx13,3\,\text{cm}$};
    \draw[<->, popmuted] (0,-5) -- (60,-5) node[midway, below] {$d=60\,\text{cm}$};
\end{tikzpicture}
\legende{Schéma du système à deux lentilles séparées (projecteur). $L_1$ forme une image réelle $A_1B_1$ qui sert d'objet virtuel pour $L_2$. L'image finale $A'B'$ est réelle.}
\end{popfigure}

\courssub{Je me prépare au Probatoire}

\begin{exercice}[Sujet type Probatoire : Projecteur diapositive -- 12 points]
On souhaite réaliser un projecteur de diapositives. La diapositive (objet $AB$) a une hauteur $h = 24\,\text{mm}$. Elle est placée devant une lentille convergente $L$ (objectif) de distance focale $f' = 50\,\text{mm}$. L'écran de projection est fixe, situé à $3,00\,\text{m}$ de la diapositive.
\begin{enumerate}
    \item On veut obtenir une image nette $A'B'$ sur l'écran. La lentille $L$ peut coulisser entre la diapositive et l'écran. Détermine par le calcul les \textbf{deux positions possibles} du centre optique $O$ de la lentille par rapport à la diapositive (distance $p = \overline{OA}$) pour obtenir une image nette sur l'écran. Donne les résultats en cm.
    \item Pour chacune de ces deux positions, calcule le grandissement $g$ et la hauteur $h'$ de l'image sur l'écran. Précise le sens de l'image (droite ou inversée) par rapport à l'objet.
    \item Laquelle des deux positions choisir pour obtenir une image la plus grande possible sur l'écran ? Justifie.
    \item On remplace la lentille $L$ par une lentille de distance focale $f' = 100\,\text{mm}$. Est-il encore possible d'obtenir une image nette sur l'écran situé à $3,00\,\text{m}$ ? Justifie par le calcul (discussion de l'équation du second degré).
    \item \textbf{Analyse dimensionnelle :} Vérifie l'homogénéité de la relation de Descartes $\dfrac{1}{p'} - \dfrac{1}{p} = \dfrac{1}{f'}$.
\end{enumerate}
\end{exercice}

\begin{exercice}[Sujet type Probatoire : Microscope composé simplifié -- 14 points]
Un microscope optique composé est modélisé par deux lentilles minces convergentes accolées (on néglige leur séparation pour simplifier).
L'objectif $L_1$ a une vergence $C_1 = +50\,\delta$. L'oculaire $L_2$ a une vergence $C_2 = +10\,\delta$.
Un objet $AB$ (une cellule) de hauteur $h = 0,10\,\text{mm}$ est placé à $2,2\,\text{cm}$ devant l'objectif $L_1$.
\begin{enumerate}
    \item Calcule la vergence $C$ et la distance focale $f'$ du système objectif-oculaire équivalent $L_{eq}$.
    \item Détermine par le calcul la position $p'$, la nature, le sens et la hauteur $h'$ de l'image finale $A'B'$ donnée par le système équivalent $L_{eq}$.
    \item Reprenons l'étude étape par étape (système non accolé, mais on garde l'hypothèse accolée pour le calcul global). Calcule le grandissement total $G$.
    \item \textbf{Réalisme :} En réalité, l'objectif et l'oculaire sont séparés par une distance $d = 16\,\text{cm}$ (longueur du tube). L'image intermédiaire $A_1B_1$ formée par l'objectif se trouve au foyer \textbf{objet} de l'oculaire ($F_2$) pour que l'image finale soit à l'infini (observation au repos de l'œil).
    \begin{enumerate}
        \item Calcule la position $p'_1$ de l'image intermédiaire $A_1B_1$ donnée par l'objectif seul ($p_1 = -2,2\,\text{cm}$).
        \item Quelle doit être la distance $O_1O_2$ entre les deux lentilles pour que $A_1B_1$ soit au foyer objet de l'oculaire $F_2$ ? Comparer à $d = 16\,\text{cm}$.
        \item Si l'image finale est à l'infini, quel est le grandissement angulaire (pouvoir grossissant) $G_v$ du microscope ? (Rappel : $G_v = \dfrac{\Delta}{f'_1} \times \dfrac{25\,\text{cm}}{f'_2}$ où $\Delta = \overline{F'_1F_2}$ est l'intervalle optique). Calcule $\Delta$ et $G_v$.
    \end{enumerate}
\end{enumerate}
\end{exercice}

\begin{popfigure}
\begin{tikzpicture}[scale=0.7, >=Stealth]
    % Axes
    \draw[popline, thick] (-3,0) -- (18,0) node[right] {Axe optique};
    % Objectif L1
    \draw[popblue, very thick] (0,-1.5) -- (0,1.5) node[above] {$L_1$ (Obj.)};
    \node[popblue, below] at (0,-2) {$O_1$};
    \foreach \x/\lab in {-2/F_1, 2/F'_1} {
        \draw[popblue, thick] (\x,-0.2) -- (\x,0.2);
        \node[popblue, below] at (\x,-0.5) {$\lab$};
    }
    % Oculaire L2
    \draw[poporange, very thick] (16,-1.5) -- (16,1.5) node[above] {$L_2$ (Ocul.)};
    \node[poporange, below] at (16,-2) {$O_2$};
    \foreach \x/\lab in {6/F_2, 26/F'_2} {
        \draw[poporange, thick] (\x,-0.2) -- (\x,0.2);
        \node[poporange, below] at (\x,-0.5) {$\lab$};
    }
    % Objet
    \draw[popgreen, very thick] (-2.2,0) -- (-2.2,0.5);
    \node[popgreen, left] at (-2.2,0.5) {$B$};
    \node[popgreen, left] at (-2.2,0) {$A$};
    % Image intermédiaire A1B1 au foyer objet F2 de l'oculaire (x=6)
    \draw[poppurple, very thick] (6,0) -- (6,-2.5);
    \node[poppurple, right] at (6,-2.5) {$B_1$};
    \node[poppurple, right] at (6,0) {$A_1$};
    % Rayons objectif
    \draw[popgreen, dashed] (-2.2,0.5) -- (0,0.5);
    \draw[popgreen, thick, ->] (0,0.5) -- (2,0) -- (6,-2.5);
    \draw[popgreen, thick, ->] (-2.2,0.5) -- (0,0) -- (16, 1.5);
    % Rayons oculaire (image à l'infini -> rayons émergents parallèles)
    \draw[poppurple, thick, ->] (6,-2.5) -- (16,0) -- (20,0);
    \draw[poppurple, dashed] (6,-2.5) -- (16,-2.5);
    \draw[poppurple, thick, ->] (16,-2.5) -- (20,-2.5);
    % Intervalle optique Delta
    \draw[<->, popgold, thick] (2,-3) -- (6,-3) node[midway, below] {$\Delta = 4\,\text{cm}$};
    % Distance O1O2
    \draw[<->, popmuted] (0,-4) -- (16,-4) node[midway, below] {$d = 16\,\text{cm}$};
\end{tikzpicture}
\legende{Microscope composé : image intermédiaire $A_1B_1$ au foyer objet $F_2$ de l'oculaire pour observation au repos (image finale à l'infini).}
\end{popfigure}

\begin{exercice}[Sujet type Probatoire : L'œil et ses corrections -- 14 points]
On modélise l'œil par un système optique unique (cornée + cristallin) assimilé à une lentille mince convergente de vergence variable $C_{\text{œil}}$, située à une distance fixe $d = 1,7\,\text{cm}$ de la rétine (plan image fixe). Un œil \textbf{emétrope} (normal) voit net un objet à l'infini sans accommodation ($C_{\text{œil}} = C_{\text{loin}}$) et voit net un objet à $25\,\text{cm}$ (point proche) avec une accommodation maximale ($C_{\text{œil}} = C_{\text{proche}}$).
\begin{enumerate}
    \item \textbf{Œil emétrope :} Calcule la vergence $C_{\text{loin}}$ nécessaire pour former l'image d'un objet à l'infini sur la rétine. Calcule la vergence $C_{\text{proche}}$ nécessaire pour former l'image d'un objet à $25\,\text{cm}$ sur la rétine. En déduire l'amplitude d'accommodation $\Delta C = C_{\text{proche}} - C_{\text{loin}}$ (en $\delta$).
    \item \textbf{Œil myope :} Un sujet myope a un point loin situé à $50\,\text{cm}$ devant son œil (il voit net sans accommodation jusqu'à $50\,\text{cm}$ max). Sa vergence au loin (au repos) est $C_{\text{loin, myope}}$.
    \begin{enumerate}
        \item Calcule $C_{\text{loin, myope}}$.
        \item Quelle est la vergence du verre correcteur (lentille mince placée au contact de l'œil, distance négligeable) pour lui permettre de voir net un objet à l'infini ? Précise la nature du verre.
        \item Ce sujet myope a une amplitude d'accommodation identique à l'emétrope ($\Delta C = 4\,\delta$). Quel est son point proche (distance minimale de vision nette) \textbf{sans correction} ? \textbf{Avec son verre correcteur} ?
    \end{enumerate}
    \item \textbf{Œil hypermétrope :} Un sujet hypermétrope a un point proche à $50\,\text{cm}$ (au lieu de $25\,\text{cm}$) avec son accommodation maximale. Sa vergence au loin (au repos) est $C_{\text{loin, hyper}}$.
    \begin{enumerate}
        \item Calcule $C_{\text{loin, hyper}}$.
        \item Détermine la vergence et la nature du verre correcteur pour lui permettre de lire à $25\,\text{cm}$.
    \end{enumerate}
    \item \textbf{Presbytie :} Avec l'âge, l'amplitude d'accommodation $\Delta C$ diminue. Un sujet emétrope de 50 ans a $\Delta C = 2\,\delta$. Quel est son point proche ? Quel verre additionnel (vergence et nature) doit-il porter pour lire à $25\,\text{cm}$ ?
\end{enumerate}
\end{exercice}

\begin{popfigure}
\begin{tikzpicture}[scale=1.2, >=Stealth]
    % Œil emétrope
    \begin{scope}[xshift=0cm]
        \draw[popline, thick] (-3,0) -- (3,0);
        % Lentille œil
        \draw[popblue, very thick] (0,-1) -- (0,1) node[above] {Œil};
        \node[popblue, below] at (0,-1.3) {$C_{\text{œil}}$};
        % Rétine
        \draw[popred, very thick] (1.7,-1.2) -- (1.7,1.2);
        \node[popred, right] at (1.8,0) {Rétine};
        % Objet infini
        \draw[popgreen, ->, thick] (-3,0.5) -- (0,0.5);
        \node[popgreen, left] at (-3,0.5) {Objet $\infty$};
        % Rayons
        \draw[popgreen, thick] (0,0.5) -- (1.7,0);
        \draw[popgreen, thick] (0,-0.5) -- (1.7,0);
        \node at (0,-2) {\textbf{Emmétropie}};
    \end{scope}
    % Œil myope (non corrigé)
    \begin{scope}[xshift=7cm]
        \draw[popline, thick] (-3,0) -- (3,0);
        \draw[popblue, very thick] (0,-1) -- (0,1) node[above] {Œil};
        \node[popblue, below] at (0,-1.3) {$C_{\text{myope}}$};
        \draw[popred, very thick] (1.7,-1.2) -- (1.7,1.2);
        \node[popred, right] at (1.8,0) {Rétine};
        % Point loin 50cm
        \draw[poporange, ->, thick] (-3,0.5) -- (0,0.5);
        \node[poporange, left] at (-3,0.5) {Point loin $50\,\text{cm}$};
        \draw[poporange, thick] (0,0.5) -- (1.7,0);
        \draw[poporange, thick] (0,-0.5) -- (1.7,0);
        % Correction
        \draw[poppurple, very thick, dashed] (-1.5,-1) -- (-1.5,1);
        \node[poppurple, above] at (-1.5,1.2) {Verre $-2\,\delta$};
        \draw[poppurple, ->, thick] (-3,0.5) -- (-1.5,0.5);
        \draw[poppurple, thick, dashed] (-1.5,0.5) -- (0,0.5); % virtuel
        \node at (0,-2) {\textbf{Myopie + Correction}};
    \end{scope}
    % Œil hypermétrope
    \begin{scope}[xshift=14cm]
        \draw[popline, thick] (-3,0) -- (3,0);
        \draw[popblue, very thick] (0,-1) -- (0,1) node[above] {Œil};
        \node[popblue, below] at (0,-1.3) {$C_{\text{hyper}}$};
        \draw[popred, very thick] (1.7,-1.2) -- (1.7,1.2);
        \node[popred, right] at (1.8,0) {Rétine};
        % Objet 25cm avec correction
        \draw[popgreen, ->, thick] (-3,0.5) -- (0,0.5);
        \node[popgreen, left] at (-3,0.5) {Objet $25\,\text{cm}$};
        % Verre correcteur
        \draw[poppurple, very thick, dashed] (-1.5,-1) -- (-1.5,1);
        \node[poppurple, above] at (-1.5,1.2) {Verre $+2\,\delta$};
        \draw[poppurple, thick] (-3,0.5) -- (-1.5,0.5) -- (0,0.8) -- (1.7,0);
        \draw[poppurple, thick] (-3,-0.5) -- (-1.5,-0.5) -- (0,-0.8) -- (1.7,0);
        \node at (0,-2) {\textbf{Hypermétropie + Correction}};
    \end{scope}
\end{tikzpicture}
\legende{Modélisation de l'œil : emmétropie, myopie (verge divergente) et hypermétropie (verge convergente). La distance lentille-rétine est fixée à $1,7\,\text{cm}$.}
\end{popfigure}

\newpage

\coursec{Corrigés des exercices}

\coursec{Corrigés détaillés — Lentilles minces}

\begin{corrige}[Exercice 1 — QCM : Reconnaissance et classification]
\begin{enumerate}
    \item \textbf{Réponse A.} La vergence $C$ et la distance focale image $f'$ (en mètres) sont liées par $C = \dfrac{1}{f'}$. Donc $f' = \dfrac{1}{C} = \dfrac{1}{5} = \SI{0,20}{m}$.
    \item \textbf{Réponse B.} C'est la propriété fondamentale du centre optique $O$ d'une lentille mince : tout rayon incident passant par $O$ n'est pas dévié (il poursuit sa route en ligne droite).
    \item \textbf{Réponse C.} Une lentille divergente ($C<0$) a un foyer image $F'$ virtuel (du côté de l'objet). Un faisceau incident parallèle à l'axe optique émerge en divergent comme s'il provenait de $F'$ (virtuel).
    \item \textbf{Réponse C.} Pour une lentille convergente, une image réelle, inversée et de même taille ($|g|=1$) s'obtient lorsque l'objet est au centre $2F$ (distance $2f'$ du centre optique), c'est-à-dire $p = -2f'$, $p' = +2f'$, $g = -1$.
\end{enumerate}
\end{corrige}

\begin{corrige}[Exercice 2 — Vrai ou Faux ? Justifie]
\begin{enumerate}
    \item \textbf{Vrai.} Par définition, la vergence $C$ (en dioptries $\delta$) est l'inverse de la distance focale image $f'$ exprimée en mètres : $C = \dfrac{1}{f'}\ (\text{avec } f' \text{ en m})$. L'analyse dimensionnelle donne $[\delta] = [\text{m}^{-1}]$.
    \item \textbf{Vrai.} Pour une lentille mince dans l'air ($n_1=n_2=1$), les distances focales objet et image sont opposées : $f = -f'$. Les foyers $F$ et $F'$ sont symétriques par rapport au centre optique $O$, qui est donc le milieu du segment $FF'$.
    \item \textbf{Vrai.} C'est le fonctionnement d'une loupe. Objet réel entre $F$ et $O$ ($-f' < p < 0$) $\Rightarrow$ image virtuelle ($p'<0$), droite ($g>0$) et agrandie ($|g|>1$).
    \item \textbf{Faux.} Le grandissement transverse est $g = \dfrac{p'}{p}$ (avec $p=\overline{OA}$, $p'=\overline{OA'}$). Pour une lentille convergente : si l'objet est réel ($p<0$) et l'image réelle ($p'>0$), alors $g<0$ (image inversée) ; si l'image est virtuelle ($p'<0$, cas de la loupe), alors $g>0$ (image droite). Le signe de $g$ n'est pas constant.
    \item \textbf{Vrai.} Théorème des vergences pour lentilles minces accolées : $C = C_1 + C_2 = +3\,\delta - 5\,\delta = -2\,\delta$. Le système équivalent est une lentille divergente.
\end{enumerate}
\end{corrige}

\begin{corrige}[Exercice 3 — Définitions et formules essentielles]
\begin{enumerate}
    \item L'unité SI de la vergence est le \textbf{dioptrie} (symbole $\delta$). La relation est $C = \dfrac{1}{f'}$ (avec $f'$ en mètres).
    \item Les trois conditions de Gauss pour une lentille mince convergente :
    \begin{enumerate}
        \item Tout rayon incident \textbf{parallèle à l'axe optique} émerge en passant par le \textbf{foyer image $F'$}.
        \item Tout rayon incident passant par le \textbf{centre optique $O$} n'est \textbf{pas dévié}.
        \item Tout rayon incident passant par le \textbf{foyer objet $F$} émerge \textbf{parallèle à l'axe optique}.
    \end{enumerate}
    \item La relation de Descartes (formule de position) : $\dfrac{1}{p'} - \dfrac{1}{p} = \dfrac{1}{f'}$, avec $p = \overline{OA}$ (distance algébrique objet), $p' = \overline{OA'}$ (distance algébrique image), $f' = \overline{OF'}$ (distance focale image).
    \item Grandissement transverse : $g = \dfrac{A'B'}{AB} = \dfrac{p'}{p}$.
    \begin{itemize}
        \item $g > 0$ : image \textbf{droite} (même orientation que l'objet).
        \item $g < 0$ : image \textbf{inversée} (orientation opposée).
        \item $|g| > 1$ : image agrandie ; $|g| < 1$ : image réduite ; $|g| = 1$ : même taille.
    \end{itemize}
    \item Théorème des vergences : Pour deux lentilles minces \textbf{accolées} (centres optiques confondus), la vergence du système équivalent est la somme des vergences : $C = C_1 + C_2$.
\end{enumerate}
\end{corrige}

\begin{corrige}[Exercice 4 — Calculs directs d'application]
\begin{enumerate}
    \item $f' = \SI{25}{cm} = \SI{0,25}{m}$.
    \[ C = \frac{1}{f'} = \frac{1}{0,25} = \SI{+4,0}{\delta}. \]
    \item Lentille divergente : $C = \SI{-2,5}{\delta} \Rightarrow f' = \dfrac{1}{C} = \SI{-0,40}{m} = \SI{-40}{cm}$.
    Objet réel : $p = \SI{-40}{cm}$.
    Relation de Descartes : $\dfrac{1}{p'} = \dfrac{1}{f'} + \dfrac{1}{p} = \dfrac{1}{-40} + \dfrac{1}{-40} = -\dfrac{2}{40} = -\dfrac{1}{20}$.
    \[ p' = \SI{-20}{cm}. \]
    Image \textbf{virtuelle} (même côté que l'objet, $p'<0$).
    Grandissement : $g = \dfrac{p'}{p} = \dfrac{-20}{-40} = +0,5$.
    Image \textbf{droite} ($g>0$) et \textbf{réduite} ($|g|<1$).
    \item $h = \SI{2}{cm}$, $p = \SI{-30}{cm}$, $f' = \SI{+10}{cm}$.
    \[ \frac{1}{p'} = \frac{1}{10} + \frac{1}{-30} = \frac{3-1}{30} = \frac{2}{30} = \frac{1}{15} \quad\Rightarrow\quad p' = \SI{+15}{cm}. \]
    Image \textbf{réelle} ($p'>0$), de l'autre côté de la lentille.
    Grandissement : $g = \dfrac{p'}{p} = \dfrac{+15}{-30} = -0,5$.
    Image \textbf{inversée} ($g<0$), \textbf{réduite} ($|g|=0,5$).
    Hauteur image : $h' = g \times h = -0,5 \times 2 = \SI{-1,0}{cm}$ (taille $\SI{1,0}{cm}$).
    \item $C_1 = \SI{+6}{\delta}$, $C_2 = \SI{-1,5}{\delta}$.
    \[ C = C_1 + C_2 = \SI{+4,5}{\delta}. \]
    \[ f' = \frac{1}{C} = \frac{1}{4,5} \approx \SI{0,222}{m} = \SI{22,2}{cm}. \]
    Système \textbf{convergent} ($C>0$).
\end{enumerate}
\end{corrige}

\begin{corrige}[Exercice 5 — Construction géométrique et analyse : La loupe]
\begin{enumerate}
    \item La construction (voir figure ci-dessous) utilise deux rayons particuliers issus de $B$ :
    \begin{itemize}
        \item Rayon parallèle à l'axe $\rightarrow$ réfracté par $F'$.
        \item Rayon passant par $O$ $\rightarrow$ non dévié.
    \end{itemize}
    Les rayons émergents sont divergents ; leurs prolongements en pointillés se coupent en $B'$ du même côté que l'objet.
    \begin{popfigure}
    \begin{tikzpicture}[scale=0.8, >=Stealth]
        % Axes
        \draw[->, thick] (-6,0) -- (3,0) node[right] {Axe optique};
        % Lentille
        \draw[thick, popblue] (0,-1.5) -- (0,1.5) node[above] {$L$};
        \node at (0, -1.8) {$O$};
        % Foyers
        \node[popblue] at (-5, -0.3) {$F$};
        \node[popblue] at (5, -0.3) {$F'$};
        \draw[popblue, dashed] (-5, -0.5) -- (-5, 0.5);
        \draw[popblue, dashed] (5, -0.5) -- (5, 0.5);
        % Objet
        \draw[thick, poporange, ->] (-4,0) -- (-4,1) node[above] {$B$};
        \node[poporange] at (-4, -0.3) {$A$};
        % Rayon 1 : parallèle -> F'
        \draw[poporange, ->] (-4,1) -- (0,1);
        \draw[poporange, ->] (0,1) -- (5,0);
        % Rayon 2 : passant par O
        \draw[poporange, ->] (-4,1) -- (0,0);
        \draw[poporange, ->] (0,0) -- (-20, 5); % prolongement virtuel
        % Image virtuelle (pointillés)
        \draw[popgreen, dashed, ->] (-20,0) -- (-20, 2.5) node[above] {$B'$};
        \node[popgreen] at (-20, -0.3) {$A'$};
        % Flèches pour distances
        \draw[<->, popmuted] (-4, -1) -- (0, -1) node[midway, below] {$|p|=4\,\text{cm}$};
        \draw[<->, popmuted] (0, -1.5) -- (5, -1.5) node[midway, below] {$f'=5\,\text{cm}$};
        \draw[<->, popmuted, dashed] (-20, -2) -- (0, -2) node[midway, below] {$|p'|=20\,\text{cm}$};
    \end{tikzpicture}
    \legende{Construction de l'image par une loupe (lentille convergente, objet entre F et O).}
    \end{popfigure}
    \item \textbf{Graphique :} $p' \approx \SI{-20}{cm}$, $h' \approx \SI{2,5}{cm}$.
    Nature : \textbf{virtuelle}. Sens : \textbf{droite}.
    \item \textbf{Calcul :} $f' = \SI{5}{cm}$, $p = \SI{-4}{cm}$.
    \[ \frac{1}{p'} = \frac{1}{5} + \frac{1}{-4} = 0,2 - 0,25 = -0,05\,\text{cm}^{-1} \quad\Rightarrow\quad p' = \SI{-20}{cm}. \]
    \[ g = \frac{p'}{p} = \frac{-20}{-4} = +5. \quad h' = g \times h = 5 \times 0,5 = \SI{2,5}{cm}. \]
    Résultats identiques à la construction.
    \item $g = +5$. L'image est \textbf{agrandie 5 fois}. L'observateur place son œil du côté opposé à l'objet (à droite de la lentille) et regarde \textbf{à travers} la lentille ; il voit l'image virtuelle du \textbf{même côté que l'objet} (à gauche de la lentille), comme si l'objet était plus grand et plus loin.
\end{enumerate}
\end{corrige}

\begin{corrige}[Exercice 6 — Détermination expérimentale de la vergence : Banc d'optique]
\begin{enumerate}
    \item Convention : Lumière de gauche vers droite. Origine au centre optique $O$ ($x_O=\SI{50,0}{cm}$). Objet à gauche $\Rightarrow p<0$. Image sur écran à droite $\Rightarrow p'>0$.
    \begin{center}
    \begin{tabular}{cccccc}
    \toprule
    $x_A$ (cm) & 30,0 & 25,0 & 20,0 & 15,0 & 10,0 \\
    \midrule
    $p = x_A - x_O$ (cm) & -20,0 & -25,0 & -30,0 & -35,0 & -40,0 \\
    \midrule
    $x_{A'}$ (cm) & 70,0 & 66,7 & 65,0 & 64,0 & 63,3 \\
    \midrule
    $p' = x_{A'} - x_O$ (cm) & +20,0 & +16,7 & +15,0 & +14,0 & +13,3 \\
    \bottomrule
    \end{tabular}
    \end{center}
    \item Inverses (en $\text{cm}^{-1}$) :
    \begin{center}
    \begin{tabular}{cccccc}
    \toprule
    $1/p$ & -0,0500 & -0,0400 & -0,0333 & -0,0286 & -0,0250 \\
    \midrule
    $1/p'$ & +0,0500 & +0,0599 & +0,0667 & +0,0714 & +0,0752 \\
    \bottomrule
    \end{tabular}
    \end{center}
    \item Relation de Descartes : $\dfrac{1}{p'} - \dfrac{1}{p} = \dfrac{1}{f'} \iff \dfrac{1}{p'} = \dfrac{1}{f'} + \dfrac{1}{p}$.
    Le graphique de $\dfrac{1}{p'}$ en fonction de $\dfrac{1}{p}$ est une \textbf{droite de pente $1$} et d'ordonnée à l'origine $\dfrac{1}{f'}$.
    \item \textbf{Graphique :} On reporte les points $(1/p,\, 1/p')$. La droite de régression (ou tracée à l'œil) coupe l'axe des ordonnées ($1/p=0$) pour $1/p' \approx \SI{0,100}{cm^{-1}}$.
    \[ \frac{1}{f'} = \SI{0,100}{cm^{-1}} \quad\Rightarrow\quad f' = \SI{10,0}{cm} = \SI{0,100}{m}. \]
    \[ C = \frac{1}{f'(\text{m})} = \frac{1}{0,100} = \SI{+10,0}{\delta}. \]
    \begin{popfigure}
    \begin{tikzpicture}
        \begin{axis}[
            width=\linewidth, height=0.5\linewidth,
            xlabel={$1/p$ ($\text{cm}^{-1}$)}, ylabel={$1/p'$ ($\text{cm}^{-1}$)},
            xmin=-0.06, xmax=0.01, ymin=0.04, ymax=0.09,
            grid=major, tick label style={/pgf/number format/fixed, /pgf/number format/precision=3},
            legend pos=north west
        ]
        \addplot[only marks, mark=*, popblue] coordinates {
            (-0.0500, 0.0500)
            (-0.0400, 0.0599)
            (-0.0333, 0.0667)
            (-0.0286, 0.0714)
            (-0.0250, 0.0752)
        };
        \addplot[poporange, domain=-0.06:0.01, samples=2] {x + 0.1};
        \addlegendentry{Données}
        \addlegendentry{Ajustement $y=x+1/f'$}
        \end{axis}
    \end{tikzpicture}
    \legende{Determination graphique de $1/f'$ (ordonnée à l'origine).}
    \end{popfigure}
    \item Incertitude de lecture $\Delta x = \pm \SI{0,1}{cm}$ (1 mm).
    Incertitude sur $p$ et $p'$ : $\Delta p \approx \Delta p' \approx \SI{0,1}{cm}$.
    Incertitude sur $1/p'$ : $\Delta(1/p') \approx \frac{\Delta p'}{p'^2}$. Pour $p'\approx \SI{15}{cm}$, $\Delta(1/p') \approx 0,1/225 \approx 4,4\times 10^{-4}\,\text{cm}^{-1}$.
    L'incertitude sur l'ordonnée à l'origine (par extrapolation) est de l'ordre de $\pm \SI{0,002}{cm^{-1}}$.
    Donc $1/f' = \SI{0,100 \pm 0,002}{cm^{-1}} \Rightarrow f' = \SI{10,0 \pm 0,2}{cm}$ et $C = \SI{10,0 \pm 0,2}{\delta}$.
\end{enumerate}
\end{corrige}

\begin{corrige}[Exercice 7 — Association de deux lentilles accolées : Correction d'amétropie]
\begin{enumerate}
    \item $C = C_1 + C_2 = \SI{+8}{\delta} - \SI{3}{\delta} = \SI{+5}{\delta}$.
    $f' = \dfrac{1}{C} = \dfrac{1}{5} = \SI{0,20}{m} = \SI{20}{cm}$.
    Système \textbf{convergent}.
    \item Objet réel : $p = \SI{-50}{cm} = \SI{-0,50}{m}$. $f' = \SI{+0,20}{m}$.
    \[ \frac{1}{p'} = \frac{1}{f'} + \frac{1}{p} = 5 + (-2) = 3\,\text{m}^{-1} \quad\Rightarrow\quad p' = \frac{1}{3}\,\text{m} \approx \SI{+33,3}{cm}. \]
    Image \textbf{réelle} ($p'>0$).
    $g = \dfrac{p'}{p} = \dfrac{+1/3}{-0,50} = -\dfrac{2}{3} \approx -0,667$.
    Image \textbf{inversée} ($g<0$), \textbf{réduite} ($|g|<1$).
    Hauteur $h' = g \times h = -0,667 \times 1 = \SI{-0,667}{cm}$ (taille $\SI{0,67}{cm}$).
    \item On veut $C_{\text{eq}} = \SI{+2}{\delta}$. $C_1 + C_2 + C_3 = +2 \Rightarrow +8 - 3 + C_3 = +2 \Rightarrow C_3 = \SI{-3}{\delta}$.
    $L_3$ est une lentille \textbf{divergente}.
    \item \textbf{Contexte ophtalmologique :} Myopie = convergence excessive de $\SI{3}{\delta}$ (donnée de l'exercice).
    Verre correcteur (au contact de l'œil) doit compenser cet excès : $C_{\text{verre}} = \SI{-3}{\delta}$.
    C'est un verre \textbf{divergent} (concave). Il forme l'image virtuelle de l'objet à l'infini au point loin de l'œil myope ($\SI{50}{cm}$).
\end{enumerate}
\end{corrige}

\begin{corrige}[Exercice 8 — Système de deux lentilles séparées : Projecteur]
Données : $f'_1 = \SI{20}{cm}$, $f'_2 = \SI{10}{cm}$, $d = O_1O_2 = \SI{60}{cm}$, $p_1 = \SI{-25}{cm}$, $h = \SI{2}{cm}$.
\begin{enumerate}
    \item \textbf{Image par $L_1$ :}
    \[ \frac{1}{p'_1} = \frac{1}{f'_1} + \frac{1}{p_1} = \frac{1}{20} + \frac{1}{-25} = \frac{5-4}{100} = \frac{1}{100} \quad\Rightarrow\quad p'_1 = \SI{+100}{cm}. \]
    Image $A_1B_1$ \textbf{réelle} (à droite de $L_1$).
    $g_1 = \dfrac{p'_1}{p_1} = \dfrac{100}{-25} = -4$.
    $h_1 = g_1 \times h = -4 \times 2 = \SI{-8}{cm}$ (inversée, agrandie 4 fois).
    \item \textbf{Objet pour $L_2$ :} $A_1$ est à $\SI{100}{cm}$ à droite de $O_1$. $O_2$ est à $\SI{60}{cm}$ à droite de $O_1$.
    $A_1$ est donc à $\SI{40}{cm}$ à \textbf{droite} de $O_2$.
    La lumière arrive de gauche sur $L_2$ ; l'objet $A_1B_1$ se trouve du côté émergent de $L_2$ : c'est un \textbf{objet virtuel} pour $L_2$.
    Distance algébrique : $p_2 = \overline{O_2A_1} = \SI{+40}{cm}$.
    \item \textbf{Image finale par $L_2$ :}
    \[ \frac{1}{p'_2} = \frac{1}{f'_2} + \frac{1}{p_2} = \frac{1}{10} + \frac{1}{40} = \frac{4+1}{40} = \frac{5}{40} = \frac{1}{8} \quad\Rightarrow\quad p'_2 = \SI{+8}{cm}. \]
    Image finale $A'B'$ \textbf{réelle} (à $\SI{8}{cm}$ à droite de $L_2$).
    $g_2 = \dfrac{p'_2}{p_2} = \dfrac{8}{40} = +0,2$.
    $h' = g_2 \times h_1 = 0,2 \times (-8) = \SI{-1,6}{cm}$.
    Image finale \textbf{inversée} par rapport à l'objet initial ($h'<0$), \textbf{réduite} ($|h'|=\SI{1,6}{cm} < \SI{2}{cm}$).
    \item Grandissement total : $G = g_1 \times g_2 = (-4) \times (+0,2) = -0,8$.
    Vérification : $h' = G \times h = -0,8 \times 2 = \SI{-1,6}{cm}$. \textbf{OK}.
    \item Schéma récapitulatif :
    \begin{popfigure}
    \begin{tikzpicture}[scale=0.5, >=Stealth]
        % Axes
        \draw[->, thick] (-5,0) -- (75,0) node[right] {Axe optique};
        % Lentille 1
        \draw[thick, popblue] (0,-3) -- (0,3) node[above] {$L_1$};
        \node at (0, -3.5) {$O_1$};
        \draw[popblue, dashed] (20,-0.5) -- (20,0.5); \node[popblue] at (20, -0.8) {$F'_1$};
        \draw[popblue, dashed] (-20,-0.5) -- (-20,0.5); \node[popblue] at (-20, -0.8) {$F_1$};
        % Lentille 2
        \draw[thick, poppurple] (60,-3) -- (60,3) node[above] {$L_2$};
        \node at (60, -3.5) {$O_2$};
        \draw[poppurple, dashed] (50,-0.5) -- (50,0.5); \node[poppurple] at (50, -0.8) {$F_2$};
        \draw[poppurple, dashed] (70,-0.5) -- (70,0.5); \node[poppurple] at (70, -0.8) {$F'_2$};
        % Objet initial
        \draw[thick, poporange, ->] (-25,0) -- (-25,2) node[above] {$B$};
        \node[poporange] at (-25, -0.5) {$A$};
        % Image 1 (réelle)
        \draw[thick, popgreen, ->] (100,0) -- (100,-8) node[below] {$B_1$};
        \node[popgreen] at (100, 0.5) {$A_1$};
        % Image finale
        \draw[thick, popred, ->] (68,0) -- (68,-1.6) node[below] {$B'$};
        \node[popred] at (68, 0.5) {$A'$};
        % Rayons principaux L1 (simplifiés)
        \draw[poporange, dashed, ->] (-25,2) -- (0,2) -- (20,0); % Parallèle -> F'1
        \draw[poporange, dashed, ->] (-25,2) -- (0,0) -- (100,-8); % Centre O1
        % Rayons principaux L2 (simplifiés)
        \draw[popgreen, dashed, ->] (100,-8) -- (60,-3.2); % Vers centre O2 (objet virtuel)
        \draw[popgreen, dashed, ->] (60,-3.2) -- (68,-1.6); % Émergent
        \draw[popgreen, dashed, ->] (100,-8) -- (50, -2); % Vers F2 -> parallèle
        \draw[popgreen, dashed, ->] (50, -2) -- (68, -1.6); % Émergent parallèle
    \end{tikzpicture}
    \legende{Schéma du projecteur à deux lentilles séparées.}
    \end{popfigure}
\end{enumerate}
\end{corrige}

\begin{corrige}[Exercice 9 — Sujet type Probatoire : Projecteur diapositive — 12 points]
\textbf{Barème indicatif :} Q1 (4 pts), Q2 (3 pts), Q3 (1 pt), Q4 (3 pts), Q5 (1 pt).
\textbf{Points de vigilance :} Conversion mm $\leftrightarrow$ cm/m ; signe des distances ($p<0, p'>0$) ; discussion du discriminant ; homogénéité dimensionnelle.
\begin{enumerate}
    \item Données : $h = \SI{24}{mm} = \SI{2,4}{cm}$, $f' = \SI{50}{mm} = \SI{5,0}{cm}$, $D = AA' = \SI{3,00}{m} = \SI{300}{cm}$.
    Lentille entre objet et écran : $p<0$, $p'>0$. Distance objet-écran : $p' - p = 300$ (car $p<0$).
    Système :
    \[ \begin{cases} \dfrac{1}{p'} - \dfrac{1}{p} = \dfrac{1}{5} \\ p' - p = 300 \end{cases} \iff \begin{cases} \dfrac{1}{p'} + \dfrac{1}{-p} = \dfrac{1}{5} \\ p' + (-p) = 300 \end{cases} \]
    Posons $u = -p > 0$ (distance objet-lentille), $v = p' > 0$ (distance lentille-écran).
    \[ \begin{cases} \dfrac{1}{u} + \dfrac{1}{v} = \dfrac{1}{5} \\ u + v = 300 \end{cases} \Rightarrow \dfrac{u+v}{uv} = \dfrac{1}{5} \Rightarrow \dfrac{300}{u(300-u)} = \dfrac{1}{5} \]
    \[ 1500 = 300u - u^2 \iff u^2 - 300u + 1500 = 0. \]
    Discriminant : $\Delta = 300^2 - 4\times 1500 = 90000 - 6000 = 84000$.
    $\sqrt{\Delta} = \sqrt{84000} = 100\sqrt{8,4} \approx 289,8\,\text{cm}$.
    \[ u_1 = \frac{300 - 289,8}{2} \approx 5,1\,\text{cm} \quad;\quad u_2 = \frac{300 + 289,8}{2} \approx 294,9\,\text{cm}. \]
    Positions du centre optique $O$ par rapport à la diapositive (objet) :
    \textbf{Position 1 :} $p_1 = -u_1 \approx \mathbf{-5,1\,cm}$ (lentille près de l'objet).
    \textbf{Position 2 :} $p_2 = -u_2 \approx \mathbf{-294,9\,cm}$ (lentille près de l'écran).
    \begin{popfigure}
    \begin{tikzpicture}[scale=0.4, >=Stealth]
        \draw[->, thick] (-5,0) -- (310,0) node[right] {Axe optique};
        % Objet
        \draw[thick, poporange, ->] (0,0) -- (0,5) node[above] {Diapo};
        % Écran
        \draw[thick, poppurple] (300,-1) -- (300,1) node[above] {Écran};
        % Position 1
        \draw[thick, popblue, dashed] (5.1,-3) -- (5.1,3) node[above] {$L_1$};
        \node[popblue] at (5.1, -3.5) {$p_1\approx -5,1$};
        % Position 2
        \draw[thick, popgreen, dashed] (294.9,-3) -- (294.9,3) node[above] {$L_2$};
        \node[popgreen] at (294.9, -3.5) {$p_2\approx -294,9$};
    \end{tikzpicture}
    \legende{Deux positions de la lentille pour une image nette sur l'écran (méthode de Bessel).}
    \end{popfigure}
    \item \textbf{Position 1 ($p_1=-5,1\,\text{cm}$) :} $p'_1 = 300 - 5,1 = 294,9\,\text{cm}$.
    $g_1 = \dfrac{p'_1}{p_1} = \dfrac{294,9}{-5,1} \approx -57,8$.
    $h'_1 = g_1 \times h = -57,8 \times 2,4 \approx \mathbf{-139\,mm} = \mathbf{-13,9\,cm}$.
    Image \textbf{inversée}, très agrandie.
    \textbf{Position 2 ($p_2=-294,9\,\text{cm}$) :} $p'_2 = 300 - 294,9 = 5,1\,\text{cm}$.
    $g_2 = \dfrac{5,1}{-294,9} \approx -0,0173$.
    $h'_2 = -0,0173 \times 2,4 \approx \mathbf{-0,041\,cm} = \mathbf{-0,41\,mm}$.
    Image \textbf{inversée}, très réduite.
    \item Choisir la \textbf{Position 1} (lentille près de la diapositive, $p \approx -5,1\,\text{cm}$).
    Justification : $|g_1| \approx 58 \gg |g_2| \approx 0,017$. L'image est beaucoup plus grande (taille $\approx \SI{14}{cm}$ contre $\SI{0,4}{mm}$).
    \item Nouveau $f' = \SI{100}{mm} = \SI{10}{cm}$.
    Équation : $u^2 - 300u + 3000 = 0$ (car $u(300-u) = 300 \times 10 = 3000$).
    $\Delta = 90000 - 12000 = 78000 > 0$.
    Deux solutions réelles positives existent ($u \approx 10,5\,\text{cm}$ et $u \approx 289,5\,\text{cm}$).
    \textbf{Oui}, il est encore possible d'obtenir une image nette (condition $D \ge 4f'$ vérifiée : $300 \ge 40$).
    \item \textbf{Analyse dimensionnelle :} $\left[\dfrac{1}{p'}\right] = \text{m}^{-1}$, $\left[\dfrac{1}{p}\right] = \text{m}^{-1}$, $\left[\dfrac{1}{f'}\right] = \text{m}^{-1}$.
    Tous les termes ont la dimension de l'inverse d'une longueur ($\text{L}^{-1}$). La relation est \textbf{homogène}.
\end{enumerate}
\end{corrige}

\begin{corrige}[Exercice 10 — Sujet type Probatoire : Microscope composé simplifié — 14 points]
\textbf{Barème indicatif :} Q1 (2 pts), Q2 (3 pts), Q3 (1 pt), Q4a (2 pts), Q4b (3 pts), Q4c (3 pts).
\textbf{Points de vigilance :} Conversion $\delta \leftrightarrow$ m/cm ; signe $p$ ; intervalle optique $\Delta = \overline{F'_1F_2}$ ; formule $G_v = \frac{\Delta}{f'_1}\frac{25}{f'_2}$ (image à l'infini).
\begin{enumerate}
    \item $C_1 = \SI{+50}{\delta} \Rightarrow f'_1 = \frac{1}{50} = \SI{0,02}{m} = \SI{2,0}{cm}$.
    $C_2 = \SI{+10}{\delta} \Rightarrow f'_2 = \frac{1}{10} = \SI{0,10}{m} = \SI{10}{cm}$.
    Lentilles accolées : $C = C_1 + C_2 = \SI{+60}{\delta}$.
    $f' = \frac{1}{60}\,\text{m} \approx \SI{0,01667}{m} = \mathbf{1,67\,cm}$.
    \item Système équivalent $L_{eq}$ : $f' = \SI{1,67}{cm}$, $p = \SI{-2,2}{cm}$.
    \[ \frac{1}{p'} = \frac{1}{f'} + \frac{1}{p} = 60 + \frac{1}{-0,022} = 60 - 45,45 = 14,55\,\text{m}^{-1}. \]
    $p' = \frac{1}{14,55} \approx \SI{0,0687}{m} = \mathbf{6,87\,cm}$.
    Image \textbf{réelle} ($p'>0$).
    $g = \frac{p'}{p} = \frac{+6,87}{-2,2} \approx \mathbf{-3,12}$.
    Image \textbf{inversée}, agrandie 3,12 fois.
    $h' = g \times h = -3,12 \times 0,10 = \mathbf{-0,312\,mm}$ (taille $\SI{0,31}{mm}$).
    \item Grandissement total $G = g = \mathbf{-3,12}$.
    \item \textbf{Réalisme (séparation $d=\SI{16}{cm}$) :}
    \begin{enumerate}
        \item Objectif seul : $p_1 = \SI{-2,2}{cm}$, $f'_1 = \SI{2,0}{cm}$.
        \[ \frac{1}{p'_1} = \frac{1}{2,0} + \frac{1}{-2,2} = \frac{2,2 - 2,0}{4,4} = \frac{0,2}{4,4} = \frac{1}{22} \quad\Rightarrow\quad p'_1 = \mathbf{+22\,cm}. \]
        Image intermédiaire $A_1B_1$ réelle, à $\SI{22}{cm}$ à droite de $O_1$.
        \item Pour observation au repos (image finale à l'infini), $A_1B_1$ doit être au foyer objet $F_2$ de l'oculaire.
        $F_2$ est à $f'_2 = \SI{10}{cm}$ à \textbf{gauche} de $O_2$.
        Distance $O_1O_2 = p'_1 + f'_2 = 22 + 10 = \mathbf{32\,cm}$.
        Donnée $d = \SI{16}{cm}$. \textbf{Incompatibilité} : la longueur de tube standard ($\SI{16}{cm}$) ne permet pas cette configuration avec ces focales pour un objet à $\SI{2,2}{cm}$. (L'objet devrait être plus loin de l'objectif).
        \item Intervalle optique $\Delta = \overline{F'_1F_2}$.
        $F'_1$ à $f'_1=\SI{2}{cm}$ à droite de $O_1$. $F_2$ à $f'_2=\SI{10}{cm}$ à gauche de $O_2$.
        $\Delta = d - f'_1 - f'_2 = 16 - 2 - 10 = \mathbf{4\,cm}$.
        Grandissement angulaire (pouvoir grossissant) :
        \[ G_v = \frac{\Delta}{f'_1} \times \frac{25\,\text{cm}}{f'_2} = \frac{4}{2} \times \frac{25}{10} = 2 \times 2,5 = \mathbf{5}. \]
        (Faible grossissement car focales longues pour un microscope).
    \end{enumerate}
\end{enumerate}
\end{corrige}

\begin{corrige}[Exercice 11 — Sujet type Probatoire : L'œil et ses corrections — 14 points]
\textbf{Barème indicatif :} Q1 (4 pts), Q2 (5 pts), Q3 (3 pts), Q4 (2 pts).
\textbf{Points de vigilance :} Modèle lentille mince unique ; distance lentille-rétine $p' = \SI{+1,7}{cm}$ fixe ; signe des vergences ; addition des vergences pour lentilles accolées (verre + œil).
\begin{enumerate}
    \item \textbf{Œil emmétrope :} $p' = \SI{+1,7}{cm} = \SI{0,017}{m}$.
    \begin{itemize}
        \item Objet à l'infini ($p=-\infty$) : $C_{\text{loin}} = \dfrac{1}{p'} = \dfrac{1}{0,017} \approx \mathbf{58,8\,\delta}$.
        \item Objet à $\SI{25}{cm}$ ($p=\SI{-0,25}{m}$) :
        $C_{\text{proche}} = \dfrac{1}{p'} - \dfrac{1}{p} = 58,8 - \dfrac{1}{-0,25} = 58,8 + 4,0 = \mathbf{62,8\,\delta}$.
        \item Amplitude d'accommodation : $\Delta C = C_{\text{proche}} - C_{\text{loin}} = \mathbf{4,0\,\delta}$.
    \end{itemize}
    \item \textbf{Œil myope :} Point loin à $\SI{50}{cm}$ ($p=\SI{-0,50}{m}$) vu net sans accommodation.
    \begin{enumerate}
        \item $C_{\text{loin, myope}} = \dfrac{1}{0,017} - \dfrac{1}{-0,50} = 58,8 + 2,0 = \mathbf{60,8\,\delta}$.
        \item Correction pour voir l'infini : Verre + Œil (accolés) doivent avoir vergence $C_{\text{loin}} = \SI{58,8}{\delta}$.
        $C_{\text{verre}} + C_{\text{loin, myope}} = 58,8 \Rightarrow C_{\text{verre}} = 58,8 - 60,8 = \mathbf{-2,0\,\delta}$.
        Verre \textbf{divergent} (concave).
        \item Amplitude $\Delta C = \SI{4}{\delta}$ (donnée). $C_{\text{proche, myope}} = C_{\text{loin, myope}} + 4 = \SI{64,8}{\delta}$.
        \textbf{Sans correction :} Point proche $p_{\text{proche}}$ tel que $\dfrac{1}{0,017} - \dfrac{1}{p_{\text{proche}}} = 64,8$.
        $58,8 - \dfrac{1}{p_{\text{proche}}} = 64,8 \Rightarrow -\dfrac{1}{p_{\text{proche}}} = 6,0 \Rightarrow p_{\text{proche}} = -\dfrac{1}{6} \approx \mathbf{-16,7\,cm}$.
        \textbf{Avec correction (verre $-2\,\delta$) :} Système corrigé a $C_{\text{loin}}=\SI{58,8}{\delta}$ et $\Delta C = \SI{4}{\delta}$ (inchangée).
        $C_{\text{proche, corrigé}} = 58,8 + 4 = \SI{62,8}{\delta}$.
        Point proche : $\dfrac{1}{0,017} - \dfrac{1}{p_{\text{proche, corr}}} = 62,8 \Rightarrow 58,8 - \dfrac{1}{p_{\text{proche, corr}}} = 62,8 \Rightarrow p_{\text{proche, corr}} = -\dfrac{1}{4} = \mathbf{-25\,cm}$.
        (Vision de près redevenue normale).
    \end{enumerate}
    \item \textbf{Œil hypermétrope :} Point proche à $\SI{50}{cm}$ avec accommodation max.
    \begin{enumerate}
        \item À accommodation max, il voit net à $\SI{50}{cm}$ : $C_{\text{proche, hyper}} = \dfrac{1}{0,017} - \dfrac{1}{-0,50} = 58,8 + 2,0 = \SI{60,8}{\delta}$.
        En supposant $\Delta C = \SI{4}{\delta}$ (physiologie identique) : $C_{\text{loin, hyper}} = C_{\text{proche, hyper}} - 4 = \mathbf{56,8\,\delta}$.
        (Vergence insuffisante au repos).
        \item Correction pour lire à $\SI{25}{cm}$ (avec accommodation max de l'œil, $C_{\text{œil}} = \SI{60,8}{\delta}$) :
        Verre + Œil (accolés) : $C_{\text{verre}} + 60,8 = \dfrac{1}{0,017} - \dfrac{1}{-0,25} = 62,8$.
        $C_{\text{verre}} = 62,8 - 60,8 = \mathbf{+2,0\,\delta}$.
        Verre \textbf{convergent} (convexe).
        \textit{Remarque :} Ce même verre $+2\,\delta$ corrige aussi la vision de loin (emmétropisation).
    \end{enumerate}
    \item \textbf{Presbytie :} Emmétrope, $\Delta C = \SI{2}{\delta}$.
    $C_{\text{loin}} = \SI{58,8}{\delta}$. $C_{\text{proche}} = 58,8 + 2 = \SI{60,8}{\delta}$.
    Point proche : $\dfrac{1}{0,017} - \dfrac{1}{p_{\text{proche}}} = 60,8 \Rightarrow 58,8 - \dfrac{1}{p_{\text{proche}}} = 60,8 \Rightarrow p_{\text{proche}} = -\dfrac{1}{2} = \mathbf{-50\,cm}$.
    Verre additionnel pour lire à $\SI{25}{cm}$ (œil au repos, $C_{\text{œil}} = \SI{58,8}{\delta}$) :
    $C_{\text{verre}} + 58,8 = 62,8 \Rightarrow C_{\text{verre}} = \mathbf{+4,0\,\delta}$.
    Verre \textbf{convergent} (loupe de lecture).
\end{enumerate}
\end{corrige}

\newpage

\coursec{Fiche bilan}

\begin{popfigure}
\begin{tikzpicture}[mindmap, grow cyclic, every node/.style={concept}, text width=2.6cm, align=center,
  root concept/.append style={concept color=popblue, font=\bfseries, text width=3cm},
  level 1/.append style={level distance=4.4cm, sibling angle=60, font=\small},
  level 1 concept/.append style={text width=2.9cm}]
\node[root concept]{Les lentilles minces}
  child[concept color=poporange]{ node {Définition et classification}
    [clockwise from=90, level 2/.append style={sibling angle=50, font=\footnotesize, text width=2.3cm}]
    child { node {Milieu transparent limité par 2 dioptres} }
    child { node {Convergentes : bords minces} }
    child { node {Divergentes : bords épais} }
  }
  child[concept color=popgreen]{ node {Éléments caractéristiques}
    [clockwise from=30, level 2/.append style={sibling angle=50, font=\footnotesize, text width=2.3cm}]
    child { node {Centre optique $O$, axe $\Delta$} }
    child { node {Foyers $F$, $F'$ ; $f'=\overline{OF'}$} }
    child { node {Vergence $C=\dfrac{1}{f'}$ (dioptrie)} }
  }
  child[concept color=popblue]{ node {Conditions de Gauss}
    [clockwise from=-30, level 2/.append style={sibling angle=50, font=\footnotesize, text width=2.3cm}]
    child { node {Rayons proches de l'axe} }
    child { node {Peu inclinés $\Rightarrow$ stigmatisme approché} }
  }
  child[concept color=popgold]{ node {Constructions géométriques}
    [clockwise from=-90, level 2/.append style={sibling angle=50, font=\footnotesize, text width=2.3cm}]
    child { node {Rayon par $O$ : non dévié} }
    child { node {Rayon $\parallel \Delta$ : passe par $F'$} }
    child { node {Rayon par $F$ : ressort $\parallel \Delta$} }
  }
  child[concept color=poppink]{ node {Formules}
    [clockwise from=-150, level 2/.append style={sibling angle=50, font=\footnotesize, text width=2.3cm}]
    child { node {$\dfrac{1}{\overline{OA'}}-\dfrac{1}{\overline{OA}}=\dfrac{1}{f'}$} }
    child { node {$\gamma=\dfrac{\overline{A'B'}}{\overline{AB}}=\dfrac{\overline{OA'}}{\overline{OA}}$} }
  }
  child[concept color=popgreen!70!black]{ node {Association et applications}
    [clockwise from=150, level 2/.append style={sibling angle=50, font=\footnotesize, text width=2.3cm}]
    child { node {$C=C_1+C_2$ (lentilles accolées)} }
    child { node {Loupe, œil, projecteur, appareil photo} }
  };
\end{tikzpicture}
\legende{Carte mentale de la leçon : les lentilles minces}
\end{popfigure}

\begin{bilan}
\textbf{Définitions clés}
\begin{itemize}
  \item \cle{Lentille mince} : milieu transparent limité par deux dioptres sphériques (ou un plan et un dioptre sphérique) dont l'épaisseur au centre est faible devant les rayons de courbure.
  \item \cle{Lentille convergente} (bords minces, symbole flèches vers l'extérieur) : transforme un faisceau parallèle en faisceau convergent ; $f'>0$, $C>0$.
  \item \cle{Lentille divergente} (bords épais, symbole flèches vers l'intérieur) : transforme un faisceau parallèle en faisceau divergent ; $f'<0$, $C<0$.
  \item \cle{Centre optique} $O$ : tout rayon passant par $O$ n'est pas dévié.
  \item \cle{Foyer image} $F'$ : point de convergence de l'image d'un objet à l'infini sur l'axe. \cle{Foyer objet} $F$ : point dont l'image est rejetée à l'infini. $F$ et $F'$ sont symétriques par rapport à $O$ : $\overline{OF}=-\overline{OF'}$.
  \item \cle{Distance focale} $f'=\overline{OF'}$ ; \cle{vergence} $C=\dfrac{1}{f'}$ en dioptries ($\delta$) si $f'$ est en mètres.
  \item \cle{Conditions de Gauss} : rayons peu inclinés sur l'axe et passant près du centre $\Rightarrow$ image nette (stigmatisme approché).
\end{itemize}

\textbf{Formules à connaître par cœur}
\begin{center}
\begin{tabularx}{\linewidth}{|l|X|l|}
\hline
\rowcolor{popblueL} \textbf{Grandeur / loi} & \textbf{Formule} & \textbf{Remarque} \\
\hline
Vergence & $C=\dfrac{1}{f'}$ & $f'$ en m, $C$ en $\delta$ \\
\hline
Formule de conjugaison (position) & $\dfrac{1}{\overline{OA'}}-\dfrac{1}{\overline{OA}}=\dfrac{1}{\overline{OF'}}=C$ & Grandeurs algébriques ! \\
\hline
Grandissement & $\gamma=\dfrac{\overline{A'B'}}{\overline{AB}}=\dfrac{\overline{OA'}}{\overline{OA}}$ & $\gamma>0$ : image droite ; $\gamma<0$ : renversée \\
\hline
Théorème des vergences & $C=C_1+C_2$ & Lentilles minces accolées \\
\hline
\end{tabularx}
\end{center}

\textbf{Les trois rayons utiles (construction d'image)}
\begin{enumerate}
  \item Rayon issu de $B$ passant par $O$ : \textbf{non dévié}.
  \item Rayon issu de $B$ parallèle à l'axe : émerge en passant par $F'$ (convergente) ou dont le prolongement passe par $F'$ (divergente).
  \item Rayon issu de $B$ passant par $F$ : émerge parallèle à l'axe.
\end{enumerate}

\textbf{Méthodes express}
\begin{itemize}
  \item \cle{Image réelle} : $\overline{OA'}>0$ (après la lentille), projetable sur un écran. \cle{Image virtuelle} : $\overline{OA'}<0$, visible à travers la lentille.
  \item Objet réel avant le foyer d'une convergente $\Rightarrow$ image réelle renversée ; objet entre $O$ et $F$ $\Rightarrow$ image virtuelle droite agrandie (loupe).
  \item Divergente + objet réel $\Rightarrow$ toujours image virtuelle, droite, réduite.
  \item Pour calculer $\overline{OA'}$ : isoler $\overline{OA'}=\dfrac{\overline{OA}\times f'}{\overline{OA}+f'}$ (avec les signes algébriques).
  \item Détermination expérimentale de $f'$ : méthode de l'objet à l'infini ($f'\approx$ distance lentille--image) ou méthode de Bessel / vérification de la formule de conjugaison au banc d'optique.
\end{itemize}
\end{bilan}

\begin{aretenir}[Checklist avant le Probatoire]
\begin{itemize}
  \item[$\square$] Je sais reconnaître une lentille convergente (bords minces) et une divergente (bords épais) à leur aspect et à leur symbole.
  \item[$\square$] Je sais placer $O$, $F$, $F'$, les plans focaux et donner le signe de $f'$ selon le type de lentille.
  \item[$\square$] Je connais la définition de la vergence $C=\dfrac{1}{f'}$ et son unité, la dioptrie ($\delta$), avec $f'$ en mètres.
  \item[$\square$] Je sais énoncer les conditions de Gauss et leur conséquence (stigmatisme approché, images nettes).
  \item[$\square$] Je sais tracer les trois rayons particuliers et construire l'image $A'B'$ d'un objet $AB$ pour toute position de l'objet.
  \item[$\square$] Je connais par cœur la formule de conjugaison $\dfrac{1}{\overline{OA'}}-\dfrac{1}{\overline{OA}}=\dfrac{1}{f'}$ et je respecte les signes des grandeurs algébriques.
  \item[$\square$] Je sais calculer et interpréter le grandissement $\gamma$ (image droite ou renversée, agrandie ou réduite).
  \item[$\square$] Je sais distinguer image réelle ($\overline{OA'}>0$, écran possible) et image virtuelle ($\overline{OA'}<0$).
  \item[$\square$] Je sais appliquer le théorème des vergences $C=C_1+C_2$ à deux lentilles accolées.
  \item[$\square$] Je sais construire l'image d'un objet à travers un système de deux lentilles (image intermédiaire servant d'objet pour la seconde).
  \item[$\square$] Je sais décrire une méthode expérimentale de détermination de $f'$ (objet à l'infini, banc d'optique).
  \item[$\square$] Je sais citer et expliquer des applications : loupe, correction de la vue (myopie, hypermétropie), appareil photographique, projecteur.
\end{itemize}
\end{aretenir}

\findecours

\end{document}
